\documentclass[twocolumn,prx,aps,superscriptaddress,longbibliography,floatfix]{revtex4-2}
\usepackage{amsmath,amssymb,amsthm}
\usepackage{graphicx}
\graphicspath{{figures/}}
\usepackage{booktabs}
\usepackage{hyperref}

\newtheorem{theorem}{Theorem}
\newtheorem{lemma}[theorem]{Lemma}
\newtheorem{proposition}[theorem]{Proposition}
\newtheorem{corollary}[theorem]{Corollary}
\newtheorem{definition}{Definition}
\newtheorem*{remark}{Remark}

% Claim markers: zero-width in production.
% \pclaim{id}{status}{scope}{backing} tags a claim for machine verification
% by the external tooling described in the verification repo; renders
% nothing in the PDF.
\newcommand{\pclaim}[4]{}

\emergencystretch=1em
\begin{document}

\title{Provable measurement design for quantum state tomography\\
from Hamming-scheme spectral theory}

\author{Andrew Craton}
\thanks{ORCID: \href{https://orcid.org/0009-0001-2269-8599}{0009-0001-2269-8599}}
\affiliation{Kootru Labs, Burlington, Wisconsin, USA}
\email{acraton@kootru.com}

\begin{abstract}
Quantum state tomography on near-term processors operates in an
underdetermined regime: a four-qubit system has $256$ Pauli
parameters but a practical budget of only ${\sim}\,100$
operators, a coverage fraction that decays as $4^{-n}$ with
qubit count, so any practical allocation must be
size-independent.
Which operators to measure matters far more than how many, yet
random Pauli sampling ignores operator weight, distributing
circuits without regard for which weight classes carry
reconstructable information.
%
We establish a correspondence, new to both combinatorics and
quantum information, between Hamming association schemes and
quantum measurement design.
\pclaim{prose:abstract:basin}{proved}{*}{thm:basin}%
Krawtchouk polynomials decompose the Pauli operator space into
weight-class eigenspaces and fix a parameter-free allocation:
a single integer $K^* = 5$ determines the full measurement
set for any $n \ge 4$ qubits, with no optimization and no
target-state input.
We prove that \emph{weight saturation}, a combinatorial
property with no prior quantum-operational content, is
sufficient for maximum-likelihood (MLE) basin separation:
every maximizer of the log-likelihood agrees on all
Pauli coordinates that carry state-dependent signal
(Theorem~\ref{thm:basin}),
with residual error $O(p/d)$ for states depolarized at
rate~$p$ in dimension $d = 2^n$ ($n$~qubits).
The Hilbert--Schmidt bound converts to an explicit fidelity
lower bound under an external near-purity hypothesis on the
MLE output (Corollary~\ref{cor:fidelity_lower_bound}), which
holds empirically with $\eta^{\star} \le 0.23$ across all
seven non-GHZ hardware runs tested.
The basin-separation guarantee extends to qudits of any local dimension~$q$.
\pclaim{prose:abstract:iff}{proved}{*}{lem:spectral_q_main}%
%
On superconducting hardware, structured selection improves
reconstruction fidelity by $\Delta F = +0.33 \pm 0.07$ over
random sampling (IBM \texttt{ibm\_fez}, four runs; corroborated
at $+0.25$ on Rigetti Ankaa-3), growing to
$+0.53 \pm 0.06$ at eight~qubits.
The advantage scales with system size: random selection
dilutes as ${\sim}\,4^{-n}$ per weight class while the
structured allocation is size-independent.
Among seven non-adaptive strategies at matched budget, $K^*$
ranks first (Table~\ref{tab:sota}).
To our knowledge, this provides the first non-adaptive
measurement design for quantum state tomography with provable
MLE identifiability guarantees, replacing heuristic operator
selection with an allocation derived from algebraic first
principles.
\end{abstract}

\keywords{quantum state tomography, Krawtchouk polynomials,
Hamming association scheme, maximum-likelihood estimation,
quantum measurement design}
\pacs{03.67.Ac, 03.65.Wj, 02.30.Zz}

\maketitle

%% ========================================================
\section{Introduction}
\label{sec:intro}

Quantum state tomography on near-term
processors~\cite{Vogel1989,James2001,Haah2017,Huang2020}
faces a measurement-allocation problem: the $4^n$ Pauli
operators far exceed any practical circuit budget, so the
choice of \emph{which} operators to measure determines
reconstruction quality.
Random Pauli sampling~\cite{Aaronson2018,Huang2020,Elben2023}
treats all weight classes equally, yet their information content
for reconstruction varies dramatically with operator weight.
No existing method provides a constructive algebraic formula
for operator selection at arbitrary qubit count; all scalable
approaches minimize a cost function or rely on randomization.
This paper shows that the allocation problem has a precise
solution rooted in the Hamming association scheme.

Krawtchouk polynomials are the spherical harmonics of the
Hamming cube~\cite{Delsarte1973}: they decompose the space of
$n$-qubit Pauli operators into $n + 1$ weight classes, each
carrying a definite eigenvalue under the Hamming association
scheme $H(n,q)$ (formal definitions in
Sec.~\ref{sec:selection}).
This decomposition is the discrete analogue of a spherical
harmonic expansion, with Hamming weight playing the role of
angular momentum.
Krawtchouk polynomials already appear across quantum
information---in error-correction
bounds~\cite{Ashikhmin2001,Rains1999}, perfect state
transfer~\cite{Vinet2012}, and Delsarte-type linear programming
bounds~\cite{MacWilliams1977,Okada2025}---but have not
previously been applied to measurement allocation for quantum
state tomography.

This paper establishes a correspondence between Hamming
association schemes and quantum measurement design, with
implications running in both directions.
The central result is a reconstruction guarantee.
\pclaim{prose:intro:basin}{proved}{*}{thm:basin}%
We prove (Theorem~\ref{thm:basin}) that \emph{weight saturation},
a combinatorial property of Pauli subsets with no prior
quantum-operational content, is sufficient for MLE basin
separation: every maximizer of the log-likelihood agrees on all
informative Pauli coordinates, so that the reconstruction is
unique on the physically relevant degrees of freedom regardless
of starting point.
As far as we are aware, this quantum-operational reading of weight
saturation is new.
Quantitative fidelity guarantees follow from the Hilbert--Schmidt
bound of Theorem~\ref{thm:basin}(iii) under an external near-purity
hypothesis on the MLE output (Corollary~\ref{cor:fidelity_lower_bound}).

From this theorem, the measurement-allocation rule follows
directly.
Krawtchouk eigenvalue masses determine how many operators to
measure at each weight class; the integer $K^*$ is the smallest
cutoff at which every weight class is saturated
(Sec.~\ref{sec:selection}).
At $n = 4$, $K^* = 5$ yields $137$~operators compiled into
$29$~circuits; the same $K^* = 5$ applies at any $n \ge 4$.
There is no optimization loop, no target-state input, no
parameters to tune.
\pclaim{prose:intro:iff}{proved}{*}{lem:spectral_q_main}%
The Krawtchouk spectral decomposition holds if and only if $q \le 3$
(Lemma~\ref{lem:spectral_q_main}); for $q \ge 4$, the full
Bose--Mesner diagonalization fails, but the weight-block-diagonal
structure that basin separation actually requires extends to
all~$q$ (Sec.~\ref{sec:universality}).

In the underdetermined regime (total measured operators
$M \ll 4^n$), the
correspondence has a direct practical payoff: the
Krawtchouk allocation concentrates measurements on the
weight classes that basin separation identifies as
load-bearing for nonlinear reconstruction.

\paragraph{Measurement allocation.}
The $4^n$ Pauli operators fall into $n+1$ weight classes.
The Krawtchouk decomposition assigns each class an eigenvalue
mass that determines how many operators to measure at that
weight (Fig.~\ref{fig:spectral};
Sec.~\ref{sec:selection} gives the explicit formula and recipe).
At $n = 4$, the allocation totals $137$ operators
across weights $0$--$4$ (Sec.~\ref{sec:selection};
Appendix~\ref{app:pauli-set}).
In simulation, the weight allocation by itself accounts for
roughly $81$\% of the advantage at $n = 4$; hardware data
corroborates the residual intra-weight effect
($76 \pm 12$\%, three seeds; Sec.~\ref{sec:weight-alloc}).

The sole experimental variable is \emph{which} Pauli operators
are measured.  Operator count, shot count, and reconstruction
algorithm are all held fixed between arms (circuit count may
differ due to basis-grouping efficiency; see
Appendix~\ref{app:pauli-set}).

\paragraph{Key results.}
The paper contributes at three levels:
\begin{itemize}
\item[\emph{(i)}] \emph{Theoretical.}
  \pclaim{prose:keyresults:basin}{proved}{*}{thm:basin}%
  Basin separation (Theorem~\ref{thm:basin}) is complemented by
  a tight purity bound $\varepsilon_{\mathrm{pos}} = O(p/d)$
  (Lemma~\ref{lem:purity_main})
  \pclaim{prose:keyresults:iff}{proved}{*}{lem:spectral_q_main}%
  and an iff spectral characterization: the Krawtchouk
  decomposition holds exactly when $q \le 3$
  (Lemma~\ref{lem:spectral_q_main}), while the
  weight-block-diagonal structure extends to all~$q$.
  The framework makes testable predictions:
  $K^* = 5$ from lattice analysis (proved for all $n \ge 4$),
  Greenberger--Horne--Zeilinger (GHZ) difficulty from weight-class mismatch ($F \approx 0.50$,
  recoverable to $0.887$ via automatic switching;
  Sec.~\ref{sec:ghz}), and
  near-sufficiency of weight allocation at $n = 4$.
\item[\emph{(ii)}] \emph{Experimental.}  $\Delta F = +0.33 \pm 0.07$
  for W~states on IBM \texttt{ibm\_fez} (four runs,
  run-to-run std.\ dev.\ $0.021$); corroborated at $+0.25$ on Rigetti
  Ankaa-3 (Sec.~\ref{sec:rigetti}).  Across $23$~test states, $K^*$ ranks first among
  seven non-adaptive strategies at matched budget
  (Table~\ref{tab:sota}).
\item[\emph{(iii)}] \emph{Operational.}  The $137$~operators compile
  into $29$~tensor-product circuits ($4.7\times$ compression).
  A compositional architecture extends the protocol to
  any linear-chain length at $\le 116$~circuits ($\Delta F = +0.53 \pm 0.06$
  at $8$~qubits, Sec.~\ref{sec:8qubit}); per-patch identifiability
  is guaranteed by Theorem~\ref{thm:basin}.
  The advantage grows with system size because random selection
  dilutes as $4^{-n}$ while the per-patch budget is
  size-independent.
\end{itemize}

\paragraph{Related work.}
Classical shadows~\cite{Aaronson2018,Huang2020,Acharya2021,King2025} and
their derandomized
variants~\cite{Huang2021,Hadfield2022} optimize measurement bases
for predicting target observables; the present work determines
operator \emph{selection} for full-state reconstruction, a
complementary task.  Overlapping tomography~\cite{Cotler2020}
reconstructs reduced density matrices from local measurements
but does not address full-state reconstruction.
Recent Pauli channel learning~\cite{ChenFlammia2024}
establishes tight bounds on noise characterization using structured
Pauli measurements; the present work applies analogous structure
to state tomography rather than process tomography.
Greedy optimal experimental design
(A/D/E-optimal)~\cite{Pukelsheim2006,Flammia2011,Granade2017} selects
operators sequentially by information gain but requires $O(M \cdot 4^n)$
computation; $K^*$ is a precomputed lookup.  Adaptive
tomography~\cite{Mahler2013,Huszar2012} relies on real-time
feedback between measurement rounds, which cloud-queue hardware
does not support.  Compressed sensing~\cite{Gross2010,Cramer2010} exploits
low-rank priors unavailable in blind certification.
Ma et al.~\cite{Ma2016} identified minimum Pauli subsets
that uniquely determine any pure state among all states (UDA)
at $n = 2$ ($11$~operators) and $n = 3$ ($31$~operators),
outperforming random subsets; their algebraic analysis
(anti-commutation structure of the complement, failing-set
classification) does not yield a constructive
formula for $n \ge 4$, where the combinatorial search becomes
intractable.
The present work addresses a different problem: not minimum
cardinality for identifiability, but spectrally prescribed
allocation for reconstruction fidelity under MLE. It provides
quantitative error bounds ($\varepsilon_{\mathrm{pos}} = O(p/d)$)
that the identifiability framework does not address.
Unlike all of these, $K^*$'s measurement set derives from
algebraic structure rather than variational optimization: every
other method in Table~\ref{tab:sota} minimizes a cost function or
maximizes an information measure, while $K^*$ has no objective
function.  Basin separation (Theorem~\ref{thm:basin}) then
provides a conditional fidelity guarantee under nonlinear
reconstruction, a regime where
frame-theoretic and Fisher-information analyses are provably
uninformative (Appendix~\ref{sec:fisher}).  Meanwhile, the
compositional architecture yields full local density matrices at
constant circuit cost on arbitrarily large processors, with an
advantage that grows as random selection becomes exponentially
dilute, a scaling property none of the above methods share.

The remainder of the paper is organized as follows.
Section~\ref{sec:method} defines the measurement selection,
reconstruction, and compositional architecture.
Section~\ref{sec:theory} presents the basin separation theorem,
purity bound, spectral dichotomy, and $q$-ary extension.
Section~\ref{sec:experiment} describes the experimental setup.
Section~\ref{sec:results} reports hardware and simulation results.
Section~\ref{sec:discussion} discusses mechanism, comparison to
existing methods, predictions, and limitations.

%% ========================================================
\section{Method}
\label{sec:method}

\subsection{Measurement selection}
\label{sec:selection}

Our structured measurement set
($M = 137$ operators for $n = 4$; Appendix~\ref{app:pauli-set}) has three
properties: (i)~\emph{non-adaptive}, requiring no mid-circuit feedback or
real-time computation; (ii)~\emph{support-complete}, meaning all $2^n$ Pauli
support classes are represented
(the lattice at $K^*$ populates every support class;
see Lemma~\ref{lem:completeness} below for a formal proof); and
(iii)~\emph{composable}, so that the same per-subsystem prescription applies
to overlapping local subsystems (Sec.~\ref{sec:compositional}).
Since $M = 137 < 255$, reconstruction requires regularized MLE
(Sec.~\ref{sec:reconstruction}).
A fourth property: (iv)~\emph{circuit-efficient}.  Because
the Krawtchouk allocation concentrates operators at low weights,
where tensor factors overlap heavily, the $137$ operators group
into only $29$ tensor-product bases via greedy set cover
($4.7\times$ compression; Table~\ref{tab:circuits}).  Each basis gives $15$ simultaneous
Pauli expectations, so $29$~circuits recover $196$ of $255$
nontrivial operators ($76.9$\%), with $60$ free additional operators
in the underdetermined weight-$3$/$4$ sectors.
By contrast, an equal-size random selection ($137$~operators,
seed-averaged) requires $50$~bases ($2.7\times$ compression),
reflecting the unstructured weight distribution; the Krawtchouk-selected
set is $1.7\times$ more circuit-efficient.
For the compositional
architecture: $\le 4\;\text{colors} \times 29 = 116$~circuits
(before basis deduplication) for any chain length
(Appendix~\ref{sec:circuits}).

\paragraph{Recipe.}
Only one input is needed: the qubit count~$n$.  Enumerate all
$m \in \mathbb{Z}^n$ with $|m|^2 \le K^*$ (where $K^* = 5$ for
$n \ge 4$; the optimality of this cutoff is proved in
Theorem~\ref{thm:spectral-char}(iv) below).
For $n \le 3$ the regime is overcomplete
($M/4^n \ge 42\%$) and weight saturation is not the
binding constraint; the spectral-completeness cutoff
$K = n$ (all $c_w > 0$) replaces $K^*$ as the
enumeration bound.
For each lattice point~$m$, reduce mod~$2$ to obtain a
\emph{parity-weight} vector $\bar m \in \mathbb{Z}_2^n$ of
Hamming weight~$w$.  The number of lattice points mapping to
weight~$w$ is the \emph{parity-weight count}
$c_w = |\{m : |m|^2 \le K^*,\; w(\bar m) = w\}|$.
Per-weight budget is then
$M_w = \min(c_w,\, A_w)$, where $A_w = 3^w \binom{n}{w}$ is
the total number of weight-$w$ Pauli operators.
Any excess $c_w - A_w > 0$ at saturated low-weight classes is
redistributed upward by a greedy procedure
(formalized in Theorem~\ref{thm:spectral-char}(iv)).
Within each weight class, $M_w$ operators are selected by any
deterministic convention (e.g., lexicographic;
Appendix~\ref{sec:layers}).
For $n = 4$, the result is a fixed set of $137$~operators
(Appendix~\ref{app:pauli-set}), independent of the target state,
the hardware platform, and the reconstruction algorithm.
Weight budget $\{1, 12, 54, 54, 16\}$ is uniquely determined
by lattice geometry; the intra-weight ordering is a reproducibility
convention (Supplemental Material~\cite{supplemental}).

\subsection{State reconstruction}
\label{sec:reconstruction}

Given expectation values $\{b_i\}_{i=1}^M$ from $M$ Pauli measurements
at $N_s$ shots each, we reconstruct the density matrix~$\rho$ via
an iterative maximum-likelihood estimator~\cite{Banaszek1999,Hradil1997,Blume-Kohout2010}
with three stabilization mechanisms (hedging, variance weighting,
and ratio clipping).

At each iteration, the update operator is
\begin{equation}\label{eq:mle}
  R = \zeta\, I + \frac{1}{M + \zeta\, d}
  \sum_{i=1}^{M} w_i\,
  \mathrm{clip}\!\left(\frac{b_i}{y_i},\, -5,\, 5\right) P_i,
\end{equation}
where $d = 2^n$, $P_i$ is the $i$th Pauli operator,
$b_i$ is the measured expectation value,
$y_i = \mathrm{tr}(P_i\,\rho)$ is the predicted expectation at the
current iterate, and
$\zeta = 0.05$ is a hedging parameter that regularizes~$\rho$
toward the maximally mixed state~$I/d$~\cite{Blume-Kohout2010}.
Ratios $b_i/y_i$ are clipped to $[-5, 5]$ to prevent divergence
when $y_i \approx 0$~\cite{Hradil1997,Paris2004}.
Per-measurement
weights $w_i = 1/\sigma_i^2$ (normalized to sum to~$M$)~\cite{Smolin2012}
incorporate shot-noise variance $\sigma_i^2 = (1 - \tilde{b}_i^2)/N_s$,
where $\tilde{b}_i = \mathrm{clip}(b_i, -0.999, 0.999)$ prevents
divergence when $|b_i| \to 1$ (the identity operator $I^{\otimes n}$
has $b = 1$ exactly; the clip keeps its weight finite).
Each iteration updates
$\rho \leftarrow \alpha\, R\rho R / \mathrm{tr}(R\rho R) +
(1 - \alpha)\,\rho$ with damping factor $\alpha = 0.5$,
increasing by $5$\% per iteration (up to $0.95$) once the
Frobenius norm of consecutive iterates falls below $0.01$ after a
$20$-iteration burn-in.  Iteration
terminates when $\|\rho^{(k+1)} - \rho^{(k)}\|_F < 10^{-7}$ or
after $1{,}000$~iterations.
Identical algorithm and hyperparameters are applied to both
structured and random measurement data.  The \emph{only} experimental
variable is the choice of which $M$~Pauli operators to measure.
Source code for the reconstruction is provided~\cite{supplemental}.
Empirically, the structured arm converges functionally (fidelity
plateaus by iteration~${\sim}300$) while the random arm oscillates
without stabilizing; convergence details and the
$n = 2$ identity-weight pathology are discussed in the Supplemental Material~\cite{supplemental}.
All hyperparameters $(\zeta, \alpha)$ are held constant across
all arms, states, and hardware platforms; a $78$-condition sweep
(7~hedge values $\times$ 6~states + 6~damping values $\times$ 6~states,
each axis varied independently) shows
that the $K^*$ advantage sign is positive across the entire tested
range (Supplemental Material~\cite{supplemental}, Sec.~VII\,H).

\subsection{Compositional architecture}
\label{sec:compositional}

For processors larger than the per-subsystem size, we decompose the
system into overlapping 4-qubit subsystems and reconstruct each
independently.  Subsystems are assigned to color groups via graph
coloring: subsystems that do not share qubits can be measured
simultaneously.  For a linear-chain topology, 4~color groups suffice.
Because the 137~$K^*$ operators compress into 29~tensor-product bases
(Sec.~\ref{sec:rigetti}; Appendix~\ref{sec:circuits}), each color
group contributes at most 29~circuits, giving an upper bound of
$4 \times 29 = 116$.  In practice, patches in different color
groups often map to the same global tensor-product basis after
extension to the full register, reducing the count further
(e.g., $116 \to 59$ unique bases for the 8-qubit linear chain;
Sec.~\ref{sec:8qubit}).
Circuit count is independent of chain length;
other topologies may require more color groups.  Figure~\ref{fig:compositional} illustrates the
architecture for an 8-qubit chain.
When patches overlap (qubits shared between adjacent subsystems),
each patch is reconstructed independently; consistency at patch
boundaries is enforced only through the shared marginals of the
per-patch density matrices, not through a global MLE.  For states
with high bipartite entanglement across a patch boundary, the
per-patch reconstructions capture the reduced state faithfully
(1-body RDM fidelities $\geq 0.98$ across all qubits in simulation;
hardware confirms $F_{\text{1-RDM}} = 0.985 \pm 0.004$,
$F_{\text{2-RDM}} = 0.771 \pm 0.031$;
Appendix~\ref{sec:rdm}), though the global state cannot be assembled
without additional cross-patch correlators.

\begin{figure}[tb]
\centering
\includegraphics[width=\columnwidth]{fig3_compositional.pdf}
\caption{Compositional architecture for an 8-qubit linear chain.
Five overlapping 4-qubit patches are assigned to 4~color groups;
patches in the same color group share no qubits and can be measured
in parallel.  With basis grouping, the total circuit count
($\le 4 \times 29 = 116$; basis deduplication reduces this
further, e.g.\ to 59 at 8~qubits) is independent of chain length.}
\label{fig:compositional}
\end{figure}

%% ========================================================
\section{Theoretical results}
\label{sec:theory}

This section presents the theoretical core: the basin separation
theorem (reverse direction, Sec.~\ref{sec:basin}), the purity bound
(Sec.~\ref{sec:purity-bound}), and the universality extension to
$q$-ary Hamming schemes (Sec.~\ref{sec:universality}).
Its forward direction (Krawtchouk eigenvalues $\to$ measurement
allocation) is in Sec.~\ref{sec:selection} and
Appendix~\ref{sec:layers}.

\subsection{Basin separation theorem}
\label{sec:basin}

We first establish the Fisher-information structure of the
reconstruction problem, then state and prove the main theorem.
The log-likelihood for quantum state tomography is concave on
the state space, so its maximizer set is convex; a \emph{basin}
is the region of initial conditions from which iterative MLE
converges to a given maximizer.
``Basin separation'' means that all basins lead to the same
reconstructed values on the informative Pauli coordinates,
so the reconstruction is unique where it matters.

Let $\mathcal{P}_n = \{I,X,Y,Z\}^{\otimes n}$ denote the $n$-qubit
Pauli group, with $\mathcal{P}_n^{(w)}$ the subset of weight-$w$
operators, and let $\mathcal{S} \subset \mathcal{P}_n \setminus \{I\}$ be
the measurement set.
Write $x_P = \mathrm{tr}(P\rho)$ for the true Pauli expectation values.
Call the unmeasured operators \emph{flat directions}
$\mathcal{F}(\mathcal{S}) = \mathcal{P}_n \setminus (\mathcal{S} \cup \{I\})$ and
the operators with nonzero expectation \emph{informative}:
$\mathcal{I}_0(\rho) = \{P : x_P \neq 0\}$.
By \emph{informative subspace} we mean the subspace of the free
Bloch parameter space spanned by the basis coordinates
$\{y_P : P \in \mathcal{I}_0(\rho) \setminus \{I\}\}$; we identify
it with its index set $\mathcal{I}_0(\rho) \setminus \{I\}$ via
this basis when context permits.

\pclaim{env:lem:hessian}{proved}{*}{lem:hessian}%
\begin{lemma}[Fisher--Hessian spectral decomposition]
\label{lem:hessian}
Let $\ell(\sigma) = \sum_{P \in \mathcal{S}} \sum_{\pm} n_{P,\pm}
\log p_{P,\pm}(\sigma)$ be the per-operator log-likelihood for
Pauli tomography with measurement set~$\mathcal{S}$, where
$p_{P,\pm} = (1 \pm \mathrm{tr}(P\sigma))/2$ and $n_{P,\pm}$ is
the $\pm 1$-outcome count from $N_s$ shots on~$P$.
In Pauli--Bloch coordinates $y_P = \mathrm{tr}(P\sigma)$, the
Fisher information matrix
$\mathcal{H}_{PP'} = -\mathbb{E}[\partial^2\ell/\partial y_P\,\partial y_{P'}]$,
evaluated at the true parameter $x_P = \mathrm{tr}(P\rho)$ with
$|x_P| < 1$ for every $P \in \mathcal{S}$, is diagonal in the
Pauli basis; its diagonal entries $h_P := \mathcal{H}_{PP}$ satisfy
\begin{equation}\label{eq:hessian}
  h_P = \frac{N_s}{1 - x_P^2}\,\mathbf{1}[P \in \mathcal{S}].
\end{equation}
Because the Pauli index set partitions by weight, equation~\eqref{eq:hessian}
yields the block decomposition
$\mathcal{H} = \bigoplus_{w=0}^{n} \mathcal{H}^{(w)}$ with
$\mathcal{H}^{(w)} = \mathrm{diag}(h_P)_{P \in \mathcal{P}_n^{(w)}}$,
which Theorem~\ref{thm:basin} uses to couple the weight grading
of~$\mathcal{S}$ to the Fisher spectrum.
The eigenvalues fall into three classes:
\begin{enumerate}
\item Informative measured operators ($P \in \mathcal{S} \cap \mathcal{I}_0$)
  have $h_P = N_s/(1 - x_P^2) > N_s$;
\item Uninformative measured operators ($P \in \mathcal{S},\, x_P = 0$)
  have $h_P = N_s$;
\item Flat directions ($P \in \mathcal{F}(\mathcal{S})$) have $h_P = 0$.
\end{enumerate}
Writing $x_{\min} = \min_{P \in \mathcal{S} \cap \mathcal{I}_0}|x_P|$ and
$x_{\max} = \max_{P \in \mathcal{S} \cap \mathcal{I}_0}|x_P|$, the restriction
$\mathcal{H}\big|_{\mathcal{S} \cap \mathcal{I}_0}$ has condition number
$\kappa_{\mathrm{info}} = (1 - x_{\min}^2)/(1 - x_{\max}^2)$,
which depends only on the extremal $|x_P|$ over the measured informative set.
\end{lemma}

\begin{proof}
\emph{Diagonality.}
The outcome probabilities $p_{P,\pm} = (1 \pm y_P)/2$
depend on $\sigma$ only through the single coordinate
$y_P = \mathrm{tr}(P\sigma)$, so
$\partial p_{P,\pm}/\partial y_{P'} =
\pm\delta_{PP'}/2$.
This factorization is the \emph{operational} consequence of
Schur orthogonality for the character group
$\mathbb{Z}_q^n$~\cite{BannaiIto1984,Holevo2011}: the Pauli
operators are the characters of $\mathbb{Z}_q^n$ (at $q = 2$,
the $n$-qubit Pauli group modulo phases), and the character
orthogonality relation
$\mathrm{tr}(PP') = d\,\delta_{PP'}$ ensures that each Pauli
measurement constrains exactly one Bloch coordinate, with zero
cross-information to all others.
Differentiating $\ell$ twice gives $-\partial^2\ell/\partial y_P^2 =
n_{P,+}/(1+y_P)^2 + n_{P,-}/(1-y_P)^2$; evaluating at $y = x$ and using
$\mathbb{E}[n_{P,\pm}] = N_s(1\pm x_P)/2$, the Fisher information reduces to
$\mathcal{H}_{PP'} = N_s\,\delta_{PP'}/(1 - x_P^2)$ for
$P, P' \in \mathcal{S}$ and zero otherwise, giving~\eqref{eq:hessian}.

\emph{Condition number.}
$\kappa_{\mathrm{info}} = \max(h_P)/\min(h_P)
= (1 - x_{\max}^2)^{-1}/(1 - x_{\min}^2)^{-1}$ over the
informative measured subset.
\end{proof}

\begin{remark}[Finite-sample independence and model scope]
Each measurement outcome $n_{P,+} \sim
\mathrm{Binomial}(N_s, (1{+}x_P)/2)$ is statistically
independent across distinct operators $P, P' \in \mathcal{S}$
(since each is measured in a separate basis).
By Hoeffding's inequality, the estimates
$\hat{y}_P = 2n_{P,+}/N_s - 1$ satisfy
\begin{equation}\label{eq:hoeffding}
  \Pr(|\hat{y}_P - x_P| > t) \le 2\exp(-N_s t^2/2)
\end{equation}
independently for each $P \in \mathcal{S}$.
When each operator~$P$ is measured in a dedicated basis,
this independence is \emph{exact} at finite~$N_s$,
not merely asymptotic.
Under tensor-product basis compression (where multiple
Paulis are extracted from common shot data), the
\emph{marginal} likelihood for each operator~$P$ still
factors as a Bernoulli in~$y_P$, because the $\pm 1$
outcome for~$P$ is a function of the bitstring parity
along $P$'s support.
The \emph{joint} likelihood within a single basis is
multinomial in the $2^n$ outcomes, which differs from
the product of per-operator marginals.
All theoretical results in Sec.~\ref{sec:basin}
(Theorem~\ref{thm:basin}, Lemma~\ref{lem:hessian})
are stated and proved under the per-operator marginal
model, and our MLE implementation (Sec.~\ref{sec:method})
takes as input the same marginal estimates and shares the
same fixed-point set (Supplemental
Material~\cite{supplemental}, Sec.~VII\,E): it receives
the vector of marginal estimates
$\hat{y}_P = 2n_{P,+}/N_s - 1$, which is identical
whether the estimates come from independent
single-operator circuits or from grouped-basis bitstring
parities.
Basis grouping therefore affects only the covariance
of the input estimates (through multinomial correlations
between co-measured operators), not the objective function
that the MLE maximizes.
Informative-subspace identifiability
(Theorem~\ref{thm:basin}) depends on the objective
being the per-operator marginal log-likelihood, which
is the model used throughout.

Diagonality of $\mathcal{H}$ is a consequence of
character orthogonality for the abelian group
$\mathbb{Z}_q^n$, the same Schur relation that underlies the
Hamming-scheme spectral theory
(Lemma~\ref{lem:spectral_q_main}).
What Lemma~\ref{lem:hessian} shows is that the weight-class
makeup of~$\mathcal{S}$ controls the spectral profile of the Fisher
information: because the Hessian is diagonal in the weight-class
basis, each eigenspace is active (nonzero eigenvalue) if and only
if $\mathcal{S}$ covers that weight class, while the per-operator
signal strength $x_P$ determines the eigenvalue magnitude.
Weight saturation therefore determines \emph{which} eigenspaces
see nonzero Fisher information; the Krawtchouk eigenvalue structure
determines \emph{how much} each weight class gets.
\end{remark}

\begin{definition}[Weight-$k$ supported state]
\label{def:weight-k-supported}
A pure $n$-qubit state $\rho$ is \emph{weight-$k$ supported}
if its weight support satisfies $W(\rho) = k$, equivalently
$x_P = 0$ for every Pauli operator $P$ with
$\mathrm{wt}(P) > k$.
\end{definition}

\pclaim{env:thm:basin}{proved}{*}{thm:basin}%
\begin{theorem}[Basin separation]
\label{thm:basin}
Let $\rho$ be an $n$-qubit state with weight support
$W(\rho) = \max\{w : \exists\, P,\; \mathrm{wt}(P) = w,\;
x_P \neq 0\}$, and let $\mathcal{S}$ be a measurement set with weight
saturation level
$w_{\mathrm{sat}}(\mathcal{S}) = \max\{k : \mathcal{P}_n^{(w)} \subseteq \mathcal{S}
\;\forall\, 1 \le w \le k\}$.

\medskip\noindent
\textbf{(i) Structured case.}
If\/ $w_{\mathrm{sat}}(\mathcal{S}) \ge W(\rho)$, every flat direction is
uninformative:
$\mathcal{F}(\mathcal{S}) \cap \mathcal{I}_0(\rho) = \varnothing$.
Fisher information $\mathcal{H}$ is positive definite on
the informative subspace $\mathcal{I}_0(\rho) \setminus \{I\}$, with spectrum
$\{N_s/(1 - x_P^2)\}_{P \in \mathcal{I}_0 \setminus \{I\}}$ and condition number
$\kappa_{\mathrm{info}} = (1{-}x_{\min}^2)/(1{-}x_{\max}^2)$,
where $x_{\min} = \min_{P \in \mathcal{I}_0 \setminus \{I\}} |x_P|$
and $x_{\max} = \max_{P \in \mathcal{I}_0 \setminus \{I\}} |x_P|$.

\medskip\noindent
\textbf{(ii) Random case.}
For a uniform random subset of size $M$, the expected number of
missing weight-$w$ operators is
$|\mathcal{P}_n^{(w)}| \cdot (4^n{-}1{-}M)/(4^n{-}1)$.
For any state with nonzero weight-$w$ signal and any
weight class with $\mu_w < 1$, this produces
informative-but-flat directions.

\medskip\noindent
\textbf{(iii) Reconstruction error (finite-sample bound).}
For pure weight-$k$ supported target states
(Definition~\ref{def:weight-k-supported}),
let $\hat\rho$ be the PSD-constrained MLE
($\hat\rho \in \arg\max_{\sigma \ge 0,\,\mathrm{tr}(\sigma)=1}
L(\sigma)$, where $L$ depends only on measured
coordinates $\{y_P\}_{P \in \mathcal{S}}$).
Denote by $\hat\rho_{\mathcal{S}}$ the projected
reconstruction with Bloch coordinates
$y_P = \mathrm{tr}(P\hat\rho)$ for
$P \in \mathcal{S}$ and $y_P = 0$ for
$P \notin \mathcal{S}$ (minimum-norm convention;
$\hat\rho_{\mathcal{S}}$ need not be PSD).
With probability at least $1 - \delta$ over the measurement
outcomes ($\delta \in (0,1)$), provided the MLE output
satisfies the bounded-Bloch regime
$|y_P(\hat\rho)| \le t_\star$ for all $P \in \mathcal{S}$,
where
$t_\star = \max_P |\hat{y}_P| + \sqrt{2\ln(2M/\delta)/N_s} < 1$
is the bounded-Bloch parameter, computable from the sample means
$\hat{y}_P$ (under $\mathcal{E}_\delta$,
$\max_P |x_P| \le t_\star$ is guaranteed;
Proposition~4 in the Supplemental Material~\cite{supplemental}),
the Hilbert--Schmidt
reconstruction error satisfies:
\begin{align}
  \|\rho - \hat\rho_{\mathcal{S}}\|_F^2 &\le
  d\,\varepsilon_{\mathrm{pos}}
  + \frac{8(1{+}t_\star^2)}{(1{-}t_\star)^2}
    \,\frac{M\ln(2M/\delta)}{d\,N_s}
  & &\text{(struct.)}, \label{eq:fid-structured} \\
  \|\rho - \hat\rho_{\mathcal{S}}\|_F^2 &\le
  d\,\varepsilon_{\mathrm{flat}}
  + d\,\varepsilon_{\mathrm{pos}}
  + \frac{8(1{+}t_\star^2)}{(1{-}t_\star)^2}
    \,\frac{M\ln(2M/\delta)}{d\,N_s}
  & &\text{(rand.)}, \label{eq:fid-random}
\end{align}
where $M = |\mathcal{S}|$ is the number of measured operators,
$\varepsilon_{\mathrm{flat}} = d^{-2}\sum_{P \in
\mathcal{F} \cap \mathcal{I}_0} x_P^2$ is irreducible
(independent of~$N_s$; depends only on $\mathcal{S}$ and $\rho$), and
$\varepsilon_{\mathrm{pos}} = d^{-2}\sum_{\mathrm{wt}(P)>k}
[\mathrm{tr}(P\hat\rho)]^2$
is the positivity excursion of the MLE output,
bounded by $\mathrm{tr}(\hat\rho^2) \le 1$
(Lemma~\ref{lem:purity_main}; computable from $\hat\rho$).
The condition $t_\star < 1$ requires
$\max_P |x_P| < 1$ and sufficient shots
$N_s > 2\ln(2M/\delta)/(1{-}\max_P|x_P|)^2$.
In the shot-noise-limited regime ($\max_P |x_P| \ll 1$ and $N_s \gg \ln(M/\delta)$),
$t_\star \to 0$ and the coefficient
$8(1{+}t_\star^2)/(1{-}t_\star)^2 = 8 + O(t_\star)$.
Both bounds require only the
Hoeffding event; conversion to fidelity $F(\rho,\hat\rho)$
requires an additional near-purity hypothesis on the MLE
output (Definition~\ref{def:near-pure}) and is carried out by
Corollary~\ref{cor:fidelity_lower_bound}.
\end{theorem}

\begin{proof}
\textbf{Part~(i).}
Weight saturation $w_{\mathrm{sat}}(\mathcal{S}) \ge W(\rho)$ means
$\mathcal{P}_n^{(w)} \subseteq \mathcal{S}$ for all $1 \le w \le W(\rho)$.
Every non-identity informative operator $P \in \mathcal{I}_0(\rho)
\setminus \{I\}$ has $1 \le \mathrm{wt}(P) \le W(\rho)$, hence
$P \in \mathcal{S}$; the identity $I$ is excluded from
$\mathcal{F}(\mathcal{S})$ by definition.
Therefore $\mathcal{F}(\mathcal{S}) \cap \mathcal{I}_0(\rho) = \varnothing$.
By Lemma~\ref{lem:hessian}, the restriction
$\mathcal{H}\big|_{\mathcal{I}_0(\rho) \setminus \{I\}}$ is diagonal
with all entries $h_P = N_s/(1 - x_P^2) > N_s > 0$,
hence positive definite on the informative subspace.
Since the log-likelihood is concave on the convex state
space~\cite{Blume-Kohout2010}, its maximizer set is convex
and every maximizer yields identical values of
$\{y_P\}_{P \in \mathcal{I}_0}$
(strict concavity of each $\log p_{P,\pm}$ in $y_P$;
see Supplemental Material~\cite{supplemental} for the
full argument).
Non-uniqueness can arise only in flat directions
$P \in \mathcal{F}(\mathcal{S})$; weight saturation ensures
$\mathcal{F} \cap \mathcal{I}_0 = \varnothing$,
so all maximizers agree on every informative coordinate.
For the random bound~\eqref{eq:fid-random}, the minimum-norm
convention ($y_P = 0$ for $P \notin \mathcal{S}$) resolves
flat-direction ambiguity and makes
$\varepsilon_{\mathrm{flat}} = d^{-2}\sum_{P \in \mathcal{F}
\cap \mathcal{I}_0} x_P^2$ exact.

\medskip\noindent
\textbf{Part~(ii).}
Indicators $\{\mathbf{1}[P \notin \mathcal{S}]\}_{P \neq I}$ are
exchangeable under uniform sampling of $M$ elements from $4^n - 1$
nontrivial Paulis.
$\mathbb{E}[\mathbf{1}[P \notin \mathcal{S}]] = (4^n{-}1{-}M)/(4^n{-}1)$;
linearity gives the stated weight-class expectation.
At $n = 4$, $M = 137$: the probability that all $12$ weight-$1$
operators survive is $\binom{243}{125}/\binom{255}{137}
< 5 \times 10^{-4}$ (Table~\ref{tab:weight-coverage}),
so at least one informative-but-flat direction exists with
probability $> 0.999$.

\medskip\noindent
\textbf{Part~(iii).}
The proof proceeds in three steps: an algebraic decomposition
of the error, a finite-sample concentration bound on the
measured coordinates, and a PSD contraction argument relating
the MLE output to the sample means.

\emph{Step 1: Algebraic decomposition.}
Decompose the squared Hilbert--Schmidt distance via the Pauli
expansion~\cite{Paris2004,Holevo2011}:
\begin{align}
  \|\rho - \hat\rho_{\mathcal{S}}\|_F^2
  &= \frac{1}{d}\sum_{P \neq I} (x_P - y_P)^2 \notag \\
  &= \frac{1}{d}\Biggl[\,
  \underbrace{\sum_{P \in \mathcal{S} \cap \mathcal{I}_0}
  \!\!\!(x_P - y_P)^2}_{\text{est.\ variance}} \notag \\
  &\qquad
  +\underbrace{\sum_{P \in \mathcal{F} \cap \mathcal{I}_0}
  \!\!\!x_P^2}_{\varepsilon_{\mathrm{flat}}\,d^2}
  +\underbrace{\sum_{\substack{P \in \mathcal{S},\; x_P = 0 \\
  \mathrm{wt}(P) \le k}}
  \!\!y_P^2}_{\text{meas.-uninf.}}
  +\underbrace{\sum_{\mathrm{wt}(P)>k}
  \!\!y_P^2}_{\le\,\varepsilon_{\mathrm{pos}}\,d^2}
  \Biggr].
  \label{eq:fidelity-decomp}
\end{align}
For weight-$k$ supported targets ($x_P = 0$ when $\mathrm{wt}(P) > k$),
the fourth sum reduces to $\sum y_P^2$ exactly.
In the structured case $\mathcal{F}\cap\mathcal{I}_0 = \varnothing$,
the second sum vanishes and the decomposition is exact.
In the random case, the minimum-norm convention
$y_P = 0$ for $P \notin \mathcal{S}$ makes the second sum
exact: $(x_P - 0)^2 = x_P^2$ for each flat informative
operator, and $y_P = 0$ for all unmeasured operators
eliminates cross-terms.
The second sum is
deterministic (depends on the measurement set, not the data).
The fourth sum is bounded by purity of the MLE output~$\hat\rho$.
Since $y_P = y_P(\hat\rho)$ for $P \in \mathcal{S}$
and $y_P = 0$ otherwise, the fourth sum satisfies
$\sum_{\mathrm{wt}(P)>k} y_P^2
\le \sum_{\mathrm{wt}(P)>k} [y_P(\hat\rho)]^2$.
Now $\hat\rho \ge 0$, $\mathrm{tr}(\hat\rho) = 1$ imply
$\mathrm{tr}(\hat\rho^2) \le 1$, hence
$\sum_{\mathrm{wt}(P)>k} [y_P(\hat\rho)]^2 \le d - 1 -
\sum_{\mathrm{wt}(P) \le k,\, P \neq I} [y_P(\hat\rho)]^2$
(the Pauli-expansion identity
$\mathrm{tr}(\hat\rho^2) = d^{-1}(1 + \sum_{P\neq I} [y_P(\hat\rho)]^2)$;
no moment-matching hypothesis
is needed, only $\hat\rho \ge 0$ and $\mathrm{tr}(\hat\rho)=1$).
This gives $d\,\varepsilon_{\mathrm{pos}} =
d^{-1}\sum_{\mathrm{wt}(P)>k} [y_P(\hat\rho)]^2$, a deterministic
quantity computable from the MLE output.
It remains to bound the first and third sums.

\emph{Step 2: Hoeffding concentration on sample means.}
Under the per-operator marginal model, each $\hat{y}_P$ is
a sample mean of $N_s$ independent $\pm 1$ outcomes, so
$\mathrm{Var}(\hat{y}_P) = (1 - x_P^2)/N_s$ \emph{exactly}
(not asymptotically)~\cite{Helstrom1976,Holevo2011,GillMassar2000}.
By Hoeffding's inequality~\eqref{eq:hoeffding} and a union
bound over all $M = |\mathcal{S}|$ measured operators, for any
$\delta \in (0,1)$:
\begin{equation}\label{eq:variance-bound}
  \Pr\Bigl(\exists\, P \in \mathcal{S}:\;
  |\hat{y}_P - x_P| > t \Bigr) \le \delta,
  \quad t = \sqrt{\frac{2\ln(2M/\delta)}{N_s}}.
\end{equation}
Under the event $\mathcal{E}_\delta = \{\forall\, P \in \mathcal{S}:\;
|\hat{y}_P - x_P| \le t\}$, which holds with probability
$\ge 1 - \delta$, the first sum (estimation variance on
sample means) satisfies
$\sum_{P \in \mathcal{S} \cap \mathcal{I}_0}
(x_P - \hat{y}_P)^2 \le M \cdot 2\ln(2M/\delta)/N_s$.
The third sum (measured-uninformative, $x_P = 0$) satisfies
the same bound since $\hat{y}_P^2 \le t^2$ when $x_P = 0$.
Both are bounded by $2M\ln(2M/\delta)/N_s$.

\emph{Step 3: PSD contraction.}
The MLE output $\hat\rho$ maximizes the per-operator
log-likelihood $L(\sigma) = \sum_{P \in \mathcal{S}} \ell_P(y_P)$
subject to $\sigma \ge 0$, $\mathrm{tr}(\sigma) = 1$,
where $\ell_P(y) = n_{P,+}\!\log\frac{1{+}y}{2}
+ n_{P,-}\!\log\frac{1{-}y}{2}$.
Each $\ell_P$ is strictly concave on $(-1,1)$ with
$\ell_P''(y) = -n_{P,+}/(1{+}y)^2 - n_{P,-}/(1{-}y)^2$.
On $|y| \le t_\star < 1$, each factor $(1 \pm y)^2$ lies in
$[(1{-}t_\star)^2,\,(1{+}t_\star)^2]$, so
$N_s/(1{+}t_\star)^2 \le |\ell_P''(y)| \le N_s/(1{-}t_\star)^2$.
By Taylor's theorem with Lagrange remainder around the
unconstrained maximizer~$\hat{y}_P$
(using $\ell_P'(\hat{y}_P) = 0$, valid since
$|\hat{y}_P| \le t_\star < 1$ in the bounded-Bloch
regime below, and the lower Hessian bound):
$\ell_P(y) \le \ell_P(\hat{y}_P) -
\tfrac{N_s}{2(1+t_\star)^2}(y - \hat{y}_P)^2$.
Under the bounded-Bloch regime
$|x_P|, |\hat{y}_P|, |y_P(\hat\rho)| \le t_\star < 1$
(the first two bounds hold under $\mathcal{E}_\delta$ with
$t_\star = \max_P |\hat{y}_P| + \sqrt{2\ln(2M/\delta)/N_s}$
provided $\max_P |\hat{y}_P| < 1$; the third is a verifiable
output property of the PSD-constrained MLE;
see Proposition~4 in the
Supplemental Material~\cite{supplemental}), the upper Hessian bound
$|\ell_P''| \le N_s/(1 - t_\star)^2$ yields the lower
Taylor bound $\ell_P(y) \ge \ell_P(\hat{y}_P)
- \tfrac{N_s}{2(1-t_\star)^2}(y - \hat{y}_P)^2$ on the same interval.
Applying the upper bound at $y = y_P(\hat\rho)$ and the lower bound
at $y = x_P$, then using $L(\hat\rho) \ge L(\rho)$
(since $\rho$ is feasible):
\begin{equation}\label{eq:psd-contraction}
  \sum_{P \in \mathcal{S}}
  \bigl(y_P(\hat\rho) - \hat{y}_P\bigr)^2
  \;\le\;
  \left(\frac{1+t_\star}{1-t_\star}\right)^{\!2}
  \sum_{P \in \mathcal{S}}
  \bigl(x_P - \hat{y}_P\bigr)^2,
\end{equation}
where $t_\star = \max_P |\hat{y}_P| + O(\sqrt{\ln(M/\delta)/N_s})$.
In the shot-noise-limited regime ($\max_P |x_P| \ll 1$ and $N_s \gg \ln(M/\delta)$),
$t_\star \to 0$ and the prefactor approaches unity;
the exact $t_\star$-dependence is retained in
Theorem~\ref{thm:basin}(iii).
Equation~\eqref{eq:psd-contraction} is a finite-sample,
deterministic bound: the PSD-constrained MLE's squared
deviation from the unconstrained optimum is at most
$\bigl(\frac{1+t_\star}{1-t_\star}\bigr)^{\!2}$ times
the true state's, on measured coordinates
(Supplemental Material~\cite{supplemental},
Proposition~4).  The contraction holds for the exact
constrained MLE; a regularized (hedged) estimator
satisfies the overall HS bound empirically but need
not satisfy~\eqref{eq:psd-contraction} exactly.
By the triangle inequality
$|x_P - y_P(\hat\rho)| \le |x_P - \hat{y}_P|
+ |\hat{y}_P - y_P(\hat\rho)|$
and~\eqref{eq:psd-contraction}:
\begin{equation}\label{eq:mle-variance-bound}
  \sum_{P \in \mathcal{S}}
  \bigl(x_P - y_P(\hat\rho)\bigr)^2
  \;\le\;
  \frac{4(1{+}t_\star^2)}{(1{-}t_\star)^2}
  \sum_{P \in \mathcal{S}}
  \bigl(x_P - \hat{y}_P\bigr)^2.
\end{equation}
In the shot-noise-limited regime ($\max_P |x_P| \ll 1$ and $N_s \gg \ln(M/\delta)$),
$t_\star \to 0$ and the prefactor
$4(1{+}t_\star^2)/(1{-}t_\star)^2 = 4 + O(t_\star)$;
the prefactor $4(1{+}t_\star^2)/(1{-}t_\star)^2$
combines with the Hoeffding factor~$2$
to give the theorem's coefficient
$8(1{+}t_\star^2)/(1{-}t_\star)^2$.
Under $\mathcal{E}_\delta$, this gives
$\frac{1}{d}\sum_{P \in \mathcal{S}}(x_P - y_P(\hat\rho))^2
\le \frac{8(1{+}t_\star^2)}{(1{-}t_\star)^2}
\frac{M\ln(2M/\delta)}{d\,N_s}$,
which absorbs the first, third, and fourth sums
(the fourth sum consists entirely of measured high-weight
coordinates $\{P \in \mathcal{S} : \mathrm{wt}(P) > k\}$,
since unmeasured high-weight coordinates are zero by
the minimum-norm convention).

\emph{Assembly.}
In the structured case~(i), $\mathcal{F} \cap \mathcal{I}_0 =
\varnothing$, so $\varepsilon_{\mathrm{flat}} = 0$ and the second
sum vanishes.
The fourth sum (measured high-weight coordinates only)
is bounded by the positivity excursion
$d\,\varepsilon_{\mathrm{pos}} =
d^{-1}\sum_{\mathrm{wt}(P)>k} y_P(\hat\rho)^2$
(which includes all high-weight coordinates of~$\hat\rho$),
itself bounded by the purity identity (Step~1).
Because the PSD contraction over~$\mathcal{S}$ already
subsumes the fourth sum, the
$d\,\varepsilon_{\mathrm{pos}}$ addend in
bounds~\eqref{eq:fid-structured}--\eqref{eq:fid-random}
is conservative; we retain it as a separate term because
it isolates the noise-dependent contribution
($d\,\varepsilon_{\mathrm{pos}} = O(p)$ under depolarizing
noise at rate~$p$; vanishes for pure weight-$k$ supported targets)
and is directly computable from the MLE output.
In the random case, $\varepsilon_{\mathrm{flat}} > 0$ generically:
the missing informative operators contribute a bias
independent of $N_s$.

\end{proof}

\begin{remark}[Random-case illustration]
At $n = 4$, $M = 137$: $46.3$\% of weight-1 operators are
expected flat, with probability $> 0.999$ that at least one
weight-1 operator is missing
(Table~\ref{tab:weight-coverage}).
For any state with nonzero
weight-1 signal (i.e., generically), this produces at least
one informative-but-flat direction.
\end{remark}

\begin{remark}[Shot-noise illustration]
Shot noise is explicit and deterministic:
at $n = 4$ ($d = 16$), $M = 137$, $N_s = 10{,}000$,
$\delta = 0.01$, and $\max_P |x_P| \approx 0$
(shot-noise-limited regime),
$t_\star = 0.045$ and the
shot-noise term evaluates to~$0.077$
($0.070$ at leading order).
For comparison, full-tomography ($M = d^2{-}1$) expected-value
shot-noise $\mathrm{E}[\cdot] \le (d^2{-}1)/(d\,N_s)$
(from the Bernoulli variance; exact at $x_P = 0$)
is not a high-probability guarantee, yet is
$\approx 44\times$ tighter at leading order
(ratio $8M\ln(2M/\delta)/(d^2{-}1)$, independent
of~$N_s$ at leading order; at $M = 137$, $\delta = 0.01$).
Under depolarizing noise,
$d\,\varepsilon_{\mathrm{pos}} = O(p)$, vanishing
in the noiseless limit;
for pure weight-$k$ supported targets, $\varepsilon_{\mathrm{pos}} = 0$
when the MLE exactly recovers the low-weight moments.
\end{remark}

\begin{remark}[Tightness of Theorem~\ref{thm:basin}(iii)]
The bound is conservative for PSD estimators
because the worst-case reflection $y_P = 2\hat{y}_P - x_P$ is
generically not PSD.
The high-probability bound trades tightness for
a deterministic guarantee at each confidence level~$\delta$.
\end{remark}

\begin{definition}[Near-pure MLE output]
\label{def:near-pure}
A PSD-constrained MLE output $\hat\rho$ is \emph{near-pure with
deficit $\eta \in [0, 1)$} if $\mathrm{tr}(\hat\rho^2) \ge 1 - \eta$.
The tightest admissible deficit is
$\eta^{\star}(\hat\rho) := 1 - \mathrm{tr}(\hat\rho^2)$,
a quantity directly computable from the MLE output.
\end{definition}

\pclaim{env:cor:fidelity_lower_bound}{proved}{*}{cor:fidelity_lower_bound}%
\begin{corollary}[Fidelity lower bound under MLE near-purity]
\label{cor:fidelity_lower_bound}
Fix the setup of Theorem~\ref{thm:basin}(iii), structured case, for
a pure weight-$k$ supported target $\rho$
(Definition~\ref{def:weight-k-supported}).
Assume additionally that the PSD-constrained MLE output is
near-pure with deficit $\eta$ (Definition~\ref{def:near-pure}).
Then, with probability at least $1 - \delta$, the fidelity
$F(\rho,\hat\rho) = \mathrm{tr}(\rho\hat\rho)$ satisfies
\begin{equation}\label{eq:fidelity-lower-bound}
  F(\rho, \hat\rho)
  \ \ge\ 1 - \frac{\eta}{2}
  - d\,\varepsilon_{\mathrm{pos}}
  - \frac{4(1{+}t_\star^2)}{(1{-}t_\star)^2}
     \,\frac{M\ln(2M/\delta)}{d\,N_s},
\end{equation}
where $\varepsilon_{\mathrm{pos}}$ and $t_\star$ are as in
Theorem~\ref{thm:basin}(iii).
The bound is tight in~$\eta$: the smallest admissible value is
$\eta^{\star}(\hat\rho) = 1 - \mathrm{tr}(\hat\rho^2)$
(Definition~\ref{def:near-pure}), which yields the strongest
guarantee at the observed purity.
The near-purity hypothesis is external to Theorem~\ref{thm:basin};
it is not proved in this paper.
The random-case analogue of eq.~\eqref{eq:fidelity-lower-bound}
incurs an additional $-d\,\varepsilon_{\mathrm{flat}}$ term on the
right-hand side, because unmeasured informative operators
($P \notin \mathcal{S}$ with $x_P \neq 0$) no longer vanish;
we state only the structured case here.
\end{corollary}
\begin{proof}
For pure $\rho$, the Hilbert--Schmidt identity gives
$\|\rho - \hat\rho\|_F^2
= \mathrm{tr}(\rho^2) - 2\mathrm{tr}(\rho\hat\rho) + \mathrm{tr}(\hat\rho^2)
= 1 - 2F + \mathrm{tr}(\hat\rho^2)$, so
$F = 1 - \tfrac{1}{2}(1 - \mathrm{tr}(\hat\rho^2))
      - \tfrac{1}{2}\|\rho - \hat\rho\|_F^2$.
Near-purity gives
$\tfrac{1}{2}(1 - \mathrm{tr}(\hat\rho^2)) \le \eta/2$, so
$F \ge 1 - \eta/2 - \tfrac{1}{2}\|\rho - \hat\rho\|_F^2$.
It remains to bound $\|\rho - \hat\rho\|_F^2$ (not the projected
norm $\|\rho - \hat\rho_{\mathcal{S}}\|_F^2$ appearing in
Theorem~\ref{thm:basin}(iii)).  In the structured case, every
unmeasured operator $P \notin \mathcal{S}$ has $\mathrm{wt}(P) > k$
(hence $x_P = 0$), so
$\|\rho - \hat\rho\|_F^2
= \|\rho - \hat\rho_{\mathcal{S}}\|_F^2
+ d^{-1}\!\sum_{P \notin \mathcal{S}} [\mathrm{tr}(P\hat\rho)]^2
\le \|\rho - \hat\rho_{\mathcal{S}}\|_F^2
+ d\,\varepsilon_{\mathrm{pos}}$.
Applying bound~\eqref{eq:fid-structured} then yields
$\|\rho - \hat\rho\|_F^2 \le 2\,d\,\varepsilon_{\mathrm{pos}}
+ \tfrac{8(1+t_\star^2)}{(1-t_\star)^2}\tfrac{M\ln(2M/\delta)}{d\,N_s}$.
Halving and substituting into $F \ge 1 - \eta/2 - \tfrac{1}{2}\|\rho-\hat\rho\|_F^2$
gives equation~\eqref{eq:fidelity-lower-bound}.
\end{proof}

\begin{remark}[Empirical status of the near-purity hypothesis]
On hardware, the near-purity hypothesis of
Corollary~\ref{cor:fidelity_lower_bound} holds with a comfortable
margin across every non-GHZ state we tested:
$\mathrm{tr}(\hat\rho^2) \ge 0.77$ across seven runs
(four W-state, three allocation-fraction) plus product states
(with $\mathrm{tr}(\hat\rho^2) = 0.996$ for
$|{+}\rangle^{\otimes 4}$), corresponding to $\eta^{\star} \le 0.23$.
GHZ is the sole exception: the information-budget mismatch
(Sec.~\ref{sec:ghz}) yields $\mathrm{tr}(\hat\rho^2) \approx 0.53$
($\eta^{\star} \approx 0.47$), and
Corollary~\ref{cor:fidelity_lower_bound} correspondingly yields a
weaker fidelity guarantee consistent with the observed degradation.
\end{remark}

\pclaim{env:cor:approx_local}{proved}{*}{cor:approx_local}%
\begin{corollary}[Approximate locality]
\label{cor:approx_local}
For any pure state~$\rho$ (not necessarily weight-$k$ supported) and any
measurement set with weight-class coverage
$\mu_w = |\mathcal{S} \cap \mathcal{P}_n^{(w)}|/A_w = 1$
for all $w \le k$,
\begin{equation}\label{eq:approx-local}
  \|\rho - \hat\rho_{\mathcal{S}}\|_F^2 \;\le\;
  2d\,(\varepsilon_{\mathrm{pos}} + \varepsilon_{\mathrm{tail}})
  + \frac{8(1{+}t_\star^2)}{(1{-}t_\star)^2}
    \,\frac{M\ln(2M/\delta)}{d\,N_s},
\end{equation}
where $\varepsilon_{\mathrm{tail}} = d^{-2}
\sum_{\mathrm{wt}(P)>k} [\mathrm{tr}(P\rho)]^2$
is the squared Frobenius norm of the Pauli expansion above
weight~$k$.
\end{corollary}
\begin{proof}
Decompose the reconstruction error~\eqref{eq:fidelity-decomp}
into three sums: statistical (bounded by
$\tfrac{8(1+t_\star^2)}{(1-t_\star)^2}\,M\ln(2M/\delta)/(d\,N_s)$
via Hoeffding concentration
and PSD contraction~\eqref{eq:psd-contraction}),
positivity excursion ($\varepsilon_{\mathrm{pos}}$ by
Lemma~\ref{lem:purity_main}), and tail leakage
($\varepsilon_{\mathrm{tail}}$, the contribution of
high-weight operators).
Weight saturation gives $\varepsilon_{\mathrm{flat}} = 0$.
For states with $W(\rho) > k$, the high-weight error
$\sum_{\mathrm{wt}(P)>k}(x_P - y_P)^2
\le 2\bigl(\sum x_P^2 + \sum y_P^2\bigr)$
by the elementary inequality $(a{-}b)^2 \le 2(a^2 + b^2)$,
giving the factor~$2$.
\end{proof}

\noindent
As a numerical check of Corollary~\ref{cor:approx_local},
for the W~state at $k = 2$, the high-weight tail is
$\varepsilon_{\mathrm{tail}} = 4.3 \times 10^{-2}$
(41 operators above weight~2 carry nonzero expectation),
so the tail contribution is small relative to the total
Pauli energy even though $W(\rho) = 4$; the full bound is
loose at $d = 16$ (the $2d$ prefactor dominates) but
tightens as $d$ grows at fixed~$k$.
All $56$ informative Pauli operators with
$0 < |x_P| < 1$ share $|x_P| = 1/2$, giving
$\kappa_{\mathrm{info}} = 1$ (perfectly conditioned
within the non-degenerate subspace).

\paragraph{Interpretation for coding theory.}
Weight saturation (the number of Hamming weight classes fully
covered by a subset of the Hamming cube) is a combinatorial
quantity that appears in Delsarte's linear-programming
bounds~\cite{Delsarte1973,BannaiIto1984} and in covering-code
theory~\cite{MacWilliams1977}.
Theorem~\ref{thm:basin} gives it a new operational meaning:
weight saturation $w_{\mathrm{sat}} \ge k$ is \emph{sufficient}
for high-fidelity nonlinear reconstruction of any state whose Pauli
expansion is supported at weight $\le k$, with the spectral gap of
the Fisher information (Lemma~\ref{lem:hessian}) guaranteeing
identifiability on the informative subspace.  Violation ($w_{\mathrm{sat}} < k$) leaves
informative operators unmeasured (flat directions in
$\mathcal{I}_0$), producing an irreducible $O(1)$ bias
(the $\varepsilon_{\mathrm{flat}}$ term in
Theorem~\ref{thm:basin}(iii)) that no increase in sample size
can remove.
Three concrete results flow back to the association-scheme
setting.  First, the sharp dichotomy
(Theorem~\ref{thm:spectral-char}(iii)): weight saturation
admits no soft intermediate: worst-case reconstruction fidelity
is either~$\le 1/2$ or~$\ge 1 - O(d^{-1})$, a rigidity
property of Hamming-scheme subsets that has no analog in
classical covering-code
theory~\cite{Delsarte1973,MacWilliams1977}.
Second, the spectral convergence threshold $K^* = q^2 + 1$ ($q = 2$)
or $q^2$ ($q \ge 3$)
(Lemma~\ref{lem:spectral_q_main}): the Krawtchouk lattice
prescription achieves weight saturation at shell radius~$q^2$
(plus one additional shell for~$q = 2$)
universally across $H(n,q)$, determined by the mod-$q$ invisible
vectors, a structural fact about Hamming-scheme spectral geometry
independent of its quantum motivation.
Third, the primitive idempotents $E_w$ of the Bose--Mesner algebra
(defined in Eq.~\ref{eq:delsarte-cert})
serve as Delsarte dual certificates of reconstruction failure
(Theorem~\ref{thm:spectral-char}(i)), the same algebraic objects
used in LP upper bounds on code sizes~\cite{Delsarte1973} but here
applied to a reconstruction-theoretic obstruction rather than a
distance-theoretic one.
These results give weight saturation new operational content but
do not, by themselves, improve existing combinatorial bounds;
whether the reconstruction dichotomy yields tighter LP constraints
on covering-code parameters remains an open question for the
Delsarte theory.

\subsection{Purity bound on positivity excursion}
\label{sec:purity-bound}

MLE reconstructs a density matrix $\sigma \ge 0$ consistent
with the measured expectations.  Since the measurement set
underdetermines the state, the positivity constraint can force
spectral weight into unmeasured high-weight Pauli directions;
the \emph{positivity excursion} $\varepsilon_{\mathrm{pos}}$
quantifies this leakage.

\pclaim{env:lem:purity_main}{proved}{*}{lem:purity_main}%
\begin{lemma}[Purity bound on positivity excursion]
\label{lem:purity_main}
Let $\sigma \in \Sigma_k(\rho) = \{\sigma \ge 0 :
\mathrm{tr}(P\sigma) = \mathrm{tr}(P\rho)\;\forall\,
\mathrm{wt}(P) \le k\}$ be any density matrix matching $\rho$
on all Pauli operators of weight $\le k$.

\textup{(i)} \emph{Bound.}
The positivity excursion satisfies
\begin{equation}\label{eq:purity}
  \varepsilon_{\mathrm{pos}} := \frac{1}{d^2}
  \sum_{\mathrm{wt}(P) > k} [\mathrm{tr}(P\sigma)]^2
  \le \frac{d - 1 - S_k}{d^2},
\end{equation}
where $S_k = \sum_{w \le k,\, P \neq I} [\mathrm{tr}(P\rho)]^2$.
For depolarized weight-$k$ supported states (noise~$p$),
$\varepsilon_{\mathrm{pos}} \le 2p(d{-}1)/d^2 = O(p/d)$.

\textup{(ii)} \emph{Tightness.}
For pure weight-$k$ supported states, $\varepsilon_{\mathrm{pos}} = 0$
(the bound is saturated).  Conversely, if the measurement set covers all
weight-${\le}\,k$ classes but leaves at least one
weight-${>}\,k$ operator unmeasured, there exists a Hermitian
trace-$1$ operator $\sigma \in \Sigma_k(\rho)$ achieving
equality in~\eqref{eq:purity}.  Whether $\sigma \ge 0$
(physicality) is a quantum marginal feasibility
question; the upper bound holds for all physical states
regardless.
\end{lemma}

\begin{proof}
\emph{Part~(i).}
Writing $\sigma = (I + \sum_{P \neq I} y_P\, P)/d$ and using
Pauli orthogonality $\mathrm{tr}(PP') = d\,\delta_{PP'}$
gives~\cite{Paris2004,Holevo2011}
\begin{equation}\label{eq:bloch-purity}
  d\,\mathrm{tr}(\sigma^2) = 1 + \sum_{P \neq I} y_P^2 \le d,
\end{equation}
since $\sigma \ge 0$ and $\mathrm{tr}(\sigma) = 1$ imply
$\mathrm{tr}(\sigma^2) \le 1$.
Marginal matching $\sigma \in \Sigma_k(\rho)$ fixes
$y_P = x_P$ for $\mathrm{wt}(P) \le k$.  Partitioning
by weight:
\[
\underbrace{\sum_{\substack{\mathrm{wt}(P) \le k \\ P \neq I}} x_P^2}_{= S_k}
+ \sum_{\mathrm{wt}(P) > k} y_P^2 \le d - 1.
\]
Dividing the second sum by $d^2$ yields~\eqref{eq:purity}.
Under depolarization $x_P \to (1{-}p)x_P$,
$S_k^{(\mathrm{noisy})} = (1{-}p)^2 S_k^{(\mathrm{pure})}$:
$\varepsilon_{\mathrm{pos}} \le
[d{-}1 - (1{-}p)^2(d{-}1)]/d^2 \le 2p(d{-}1)/d^2$.

\emph{Part~(ii).}
For a pure weight-$k$ supported state, $\mathrm{tr}(\rho^2) = 1$ and
$x_P = 0$ for $\mathrm{wt}(P) > k$, so
$S_k = d - 1$ and $\varepsilon_{\mathrm{pos}} = 0$.

For the converse, we construct $\sigma$ explicitly.
Let $P_0 \notin \mathcal{S}$ be any unmeasured operator with
$\mathrm{wt}(P_0) > k$
(which exists whenever the measurement set leaves at least one
high-weight operator uncovered; in particular, whenever
$|\mathcal{S}| < d^2 - 1$ and $\mathcal{S}$ spans all
weight-${\le}\,k$ classes, the regime of interest for
Theorem~\ref{thm:basin}),
and define
\[
  \sigma_\alpha = \frac{I + \sum_{\substack{\mathrm{wt}(P) \le k \\ P \neq I}}
  x_P\, P + \alpha\, P_0}{d},
\]
with $\alpha$ chosen to saturate purity:
$\alpha^2 = d - 1 - S_k$.  Then $\sigma_\alpha$ matches
$\rho$ on all weight-${\le}\,k$ operators,
$\mathrm{tr}(\sigma_\alpha^2) = (1 + S_k + \alpha^2)/d = 1$,
and $\varepsilon_{\mathrm{pos}} = \alpha^2/d^2
= (d{-}1{-}S_k)/d^2$, achieving equality among
Hermitian trace-$1$ operators in~$\Sigma_k(\rho)$.
Whether $\sigma_\alpha \ge 0$ (i.e., whether the
bound is attained by a physical state) is a quantum
marginal feasibility problem, known to be
$\mathrm{QMA}$-complete in
general~\cite{Liu2006,Wei2010}.
This does not affect
the \emph{upper} bound: every $\sigma \ge 0$ in
$\Sigma_k(\rho)$ satisfies
$\varepsilon_{\mathrm{pos}} \le (d{-}1{-}S_k)/d^2$
regardless, and this is the only direction used in the
fidelity guarantee (Theorem~\ref{thm:basin}(iii)).
\end{proof}

\begin{remark}
The bound exploits the spectral stratification of the
Bloch body by Pauli weight~\cite{BannaiIto1984}: orthogonality
$\mathrm{tr}(PP') = d\,\delta_{PP'}$ decomposes the purity
constraint $\|\mathbf{y}\|^2 \le d - 1$ into independent
contributions from each weight class, reducing the quantum
marginal problem ($\mathrm{QMA}$-complete in
general~\cite{Liu2006,Wei2010}) to an arithmetic inequality
on the spectral budget.
At $n = 4$, $p = 0.03$:
$\varepsilon_{\mathrm{pos}} \le 0.0035$, consistent with
numerical verification
(Supplemental Material~\cite{supplemental}).
\end{remark}

\subsection{Spectral necessity and dichotomy}
\label{sec:spectral-necessity}

We first define the central algebraic object.
All $q^n$ generalized Pauli support classes are indexed by
$s \in \mathbb{Z}_q^n$.  For each integer vector
$m \in \mathbb{Z}^n$, reduction mod~$q$ gives a support class
$\bar{m} = m \bmod q \in \mathbb{Z}_q^n$; the Hamming weight
$w(\bar{m})$ equals the number of nonzero coordinates of~$\bar{m}$,
i.e., the Pauli weight of the corresponding operator.
A \emph{Gram matrix} at cutoff~$K$ is defined by the truncated character sum
\begin{equation}\label{eq:gram-def}
  G_{ss'}(K)
  = \sum_{\substack{m \in \mathbb{Z}^n \\ |m|^2 \le K}}
  \omega^{\bar{m} \cdot (s - s')},
\end{equation}
where $\omega = e^{2\pi i/q}$.
For $q \le 3$, the balanced lift
$\{-(q{-}1)/2,\dots,(q{-}1)/2\}$ assigns equal squared norm to
every nonzero residue, so $G(K)$ is distance-class-constant and
lies in the Bose--Mesner algebra of $H(n,q)$
(the commutative matrix algebra spanned by the distance-class
adjacency matrices of the association
scheme~\cite{Delsarte1973,BannaiIto1984}); the Krawtchouk
polynomials diagonalize it in the idempotent basis.
For $q \ge 4$, the Euclidean cutoff distinguishes nonzero residues
with different squared-lift norms, so $G(K)$ is not
distance-class-constant; it remains weight-block-diagonal
($S_n$-invariant), which suffices for basin separation
(Lemma~\ref{lem:spectral_q_main}).
$G(K)$ measures how well the lattice modes up to norm~$K$ distinguish
between support classes (Fig.~\ref{fig:spectral}): $G(K)$ has full
rank exactly when all support classes are spectrally resolved.
Euclidean cutoff $|m|^2 \le K$ respects the coordinate
permutation symmetry of~$H(n,q)$ and connects to classical
lattice-counting results
(Jacobi's four-square theorem~\cite{HardyWright2008} at $n = 4$),
giving closed-form mode counts.
Full rank of~$G(K)$ requires $c_w(K) > 0$ for every weight~$w$,
which is the spectral condition underlying weight saturation
(Theorem~\ref{thm:basin}).
Basin separation (Theorem~\ref{thm:basin}) shows weight saturation is sufficient
for reconstruction; Theorem~\ref{thm:spectral-char} below shows
it is necessary.

\begin{figure}[tb]
\centering
\includegraphics[width=\columnwidth]{fig2_spectral_correspondence.pdf}
\caption{Krawtchouk spectral correspondence at $n = 4$, $q = 2$.
(a)~Lattice shells $|m|^2 \le K^* = 5$ in $\mathbb{Z}^4$ map to Pauli
weight classes via mod-$2$ reduction $m \mapsto \bar{m}$.  Shell
multiplicities $r_4(k)$ (Jacobi's four-square theorem) and operator
counts $M_w/A_w$ are shown; weights $w \le 2$ are fully saturated
($M_w = A_w$), while $w = 3, 4$ are subsampled by the lattice cutoff.
Thick arrows indicate dominant flows; thin arrows show secondary
contributions from mixed-parity shells.
(b)~Gram eigenvalues $\lambda_w = 16\, c_w / \binom{4}{w}$
from the idempotent decomposition~\eqref{eq:spectral_q_main}.
Weight~$1$ carries the largest mass ($40.9$\% of $\mathrm{tr}\,G$)
despite having only $12$ operators; the nonuniform eigenvalue
distribution drives the $K^*$ measurement allocation.}
\label{fig:spectral}
\end{figure}

The eigenvalue structure of the Gram matrix improves
monotonically with the spectral cutoff:

\pclaim{env:lem:monotone}{proved}{*}{lem:monotone}%
\begin{lemma}[Eigenvalue monotonicity]
\label{lem:monotone}
For the Gram matrix $G(K) = \sum_w \lambda_w(K)\, E_w$ of
the Hamming scheme $H(n,q)$:
\begin{enumerate}
\item $\lambda_w(K)$ is non-decreasing in $K$ for every
  weight class~$w$.
\item $\lambda_w(K') > \lambda_w(K)$ if and only if
  the shell at $|m|^2 = K'$ contains at least one
  lattice vector of parity weight~$w$.
\item The redistributed budget $M_w(K)$
  (defined by assigning $\min(c_w{+}e,\, A_w)$ to each
  weight class $w = 0,1,\ldots,n$ in order, with running
  excess $e$ initialized to zero) is
  non-decreasing in $K$ for each $w$.
\end{enumerate}
\end{lemma}
\begin{proof}
(1)~Each new shell adds non-negative counts:
$c_w(K') = c_w(K) + |\{m : |m|^2 = K',\; w(\bar{m}) = w\}|
\ge c_w(K)$.  Since $\lambda_w = q^n c_w / v_w$
(where $v_w = \binom{n}{w}(q{-}1)^w$ is the Hamming-scheme multiplicity),
monotonicity follows.
(2)~Strict inequality holds iff the new shell contributes
at least one vector at parity weight~$w$, which increases
$c_w$ and hence $\lambda_w$.
(3)~The greedy redistribution processes weight classes
in order $w = 0, 1, \ldots, n$, maintaining a running
excess $e$.  Initialize $e = 0$.  At each~$w$:
$M_w(K) = \min(c_w(K) + e,\; A_w)$ and
$e \leftarrow c_w(K) + e - M_w(K)$.
Since $c_w(K)$ is non-decreasing in~$K$ for every~$w$
by~(1), the input to each stage is non-decreasing, so
each~$M_w(K)$ is non-decreasing and each running
excess is non-decreasing.  (Formally: induction on~$w$
using the fact that $\min(a,b)$ and $a - \min(a,b)$
are both non-decreasing when~$a$ is non-decreasing
and~$b$ is fixed.)
\end{proof}

Theorem~\ref{thm:basin} showed that weight saturation is
\emph{sufficient} for basin separation.  The next result shows it
is also \emph{necessary}: if any weight class is incomplete,
one can construct a pair of states that are indistinguishable
from the measured data.  For $q \le 3$, the argument constructs Delsarte dual certificates
of spectral insufficiency in the Bose--Mesner algebra
of~$H(n,q)$~\cite{Delsarte1973,BannaiIto1984}, the same
algebraic machinery that yields linear-programming upper bounds
on code sizes, here repurposed as reconstruction obstructions.
For $q \ge 4$, the weight-block-diagonal structure yields an
analogous obstruction via weight-class projectors.

\pclaim{env:thm:spectral-char}{proved}{n>=2,q=2}{thm:spectral-char}%
\begin{theorem}[Spectral characterization of basin separation]
\label{thm:spectral-char}
For any $n \ge 2$ and for any measurement set
$\mathcal{S} \subset \mathcal{P}_n \setminus \{I\}$,
define the weight-class coverage
$\mu_w = |\mathcal{S} \cap \mathcal{P}_n^{(w)}| / A_w$,
where $A_w = \binom{n}{w}(q^2{-}1)^w$
(with $q = 2$ throughout parts (i)--(iii); part (iv) also
addresses $q \ge 3$).

\textup{(i)} \emph{Spectral obstruction (necessity).}
If $\mu_{w_0} < 1$ for some $w_0 \le k$, there exist states
$\rho_1, \rho_2$ with $W(\rho_i) = w_0 \le k$ such that
$\rho_1$ and $\rho_2$ produce identical measurement
data on $\mathcal{S}$, yet $F(\rho_1, \rho_2) = 0$.
No estimator (linear, MLE, or otherwise) achieves worst-case
fidelity $> 1/2$ over the class~$\mathcal{W}_k = \{\rho : W(\rho) \le k\}$ of weight-$k$ supported states.

\textup{(ii)} \emph{Spectral resolution (sufficiency).}
If $\mu_w = 1$ for all $w \le k$ and the Krawtchouk eigenvalues
$\lambda_w > 0$ for all~$w$
(the latter follows from $\mu_w = 1$ for $w \le k$;
it is an additional condition for $w > k$, ensuring
support-completeness of the Gram matrix),
then for all pure $\rho$ with
$W(\rho) \le k$:
\begin{equation}\label{eq:spectral-resolution}
  \|\rho - \hat\rho_{\mathcal{S}}\|_F^2 \le \frac{d{-}1{-}S_k}{d}
  + \frac{8(1{+}t_\star^2)}{(1{-}t_\star)^2}
    \,\frac{M\ln(2M/\delta)}{d\,N_s},
\end{equation}
with probability $\ge 1 - \delta$.

\textup{(iii)} \emph{Sharp dichotomy.}
Parts~(i) and~(ii) combine into a dichotomy:
the impossibility branch~(i) holds over all
of~$\mathcal{W}_k$, while the high-fidelity
branch~(ii) holds for pure weight-$k$ supported states
conditional on near-purity of the MLE output
and sufficient shots $N_s = \Omega(d^2)$.
Worst-case fidelity is either~$\le 1/2$
or~$\ge 1 - O(d^{-1})$, with no intermediate regime
(over measurement sets).
The fidelity gap $1 - O(d^{-1}) \to 1$ as $d\to\infty$.

\textup{(iv)} \emph{Spectral progression and optimality of $K^*$.}
Define the Krawtchouk weight budget at cutoff~$K$ as
$c_w(K) = |\{m \in \mathbb{Z}^n : |m|^2 \le K,\;
w(\bar{m}) = w\}|$, and let the \emph{redistributed
budget} $M_w(K)$ be the allocation that saturates the
lowest weight classes first (assigning $\min(c_w, A_w)$
to each class in order, with excess flowing to
higher-weight classes).
The \emph{full-saturation threshold}
\begin{equation}\label{eq:ksat-def}
  K_{\mathrm{sat}} = \min\bigl\{K : M_w(K) = A_w\;\;
  \forall\, w \le \lfloor n/2 \rfloor\bigr\}
\end{equation}
is the minimal spectral cutoff at which the redistributed
budget achieves weight saturation for all
weight classes relevant to weight-$k$ supported tomography.

For $q = 2$, $K^* = 5$ saturates weight classes $1$ and $2$ for
every $n \ge 4$ (proved by direct lattice enumeration:
the parity-weight count $c_w(5)$ exceeds $A_w$ for $w \in \{1,2\}$
at all $n \ge 4$; Supplemental Material~\cite{supplemental}).
Full saturation to $\lfloor n/2 \rfloor$ holds only for
$n \le 5$, where $\lfloor n/2 \rfloor \le 2$ and the
weight-$1,2$ guarantee suffices.
For $n \ge 6$, the monolithic full-saturation cutoff $K_{\mathrm{full}}(n)$
grows with~$n$ and the operator budget scales rapidly
($137$ at $n{=}4$; $1{,}341$ at $n{=}6$, $K_{\mathrm{full}}{=}6$;
$12{,}369$ at $n{=}8$, $K_{\mathrm{full}}{=}7$;
Sec.~\ref{sec:limitations}).
The compositional architecture (Sec.~\ref{sec:8qubit})
circumvents this scaling by decomposing into $4$-qubit patches,
each at $K^* = 5$ with full saturation, providing an $18\times$
operator-count saving at $n = 8$.

For $q \ge 3$, the Krawtchouk-prescribed cutoff is the
\emph{spectral mass threshold} $K_{\mathrm{mass}} = q^2$
(Lemma~\ref{lem:spectral_q_main}), where the mod-$q$
invisible vectors enter the lattice and the eigenvalue
distribution undergoes a structural transition.
At $K_{\mathrm{mass}}$, every weight class has positive
budget ($c_w > 0$ for all~$w$), but weight saturation
is partial: e.g., $M_1 = 16$ vs.\ $A_1 = 32$ ($50$\%) at
$q = 3$, $n = 4$.  Full weight saturation~\eqref{eq:ksat-def}
requires $K_{\mathrm{sat}} > K_{\mathrm{mass}}$ (e.g.,
$K_{\mathrm{sat}} = 14$ at $q = 3$, $n = 4$;
Supplemental Material~\cite{supplemental}).
Throughout this paper, $K^*$ denotes the Krawtchouk-prescribed
cutoff: $q^2 + 1$ for $q = 2$ (where it equals the operational
threshold for $n \le 5$) and $K_{\mathrm{mass}} = q^2$ for
$q \ge 3$ (where it gives spectral completeness but not full
weight saturation).
The compositional architecture compensates for partial
saturation (whether from large~$n$ at $q = 2$ or from $q \ge 3$)
by reducing the effective patch size
(Supplemental Material~\cite{supplemental}).

Below $K^*$, insufficiency is established by two
mechanisms:
\begin{itemize}
\item For $K \le n{-}1$: $c_{K+1}(K) = 0$, so the
  primitive idempotent $E_{K+1}$ of $H(n,q)$
  provides a Delsarte dual certificate of spectral
  incompleteness (no estimator can resolve
  weight-$(K{+}1)$ states).
\item At $K = n$: all $c_w > 0$ (spectral completeness),
  but the Krawtchouk budget underallocates low-weight
  classes.  Any reallocation that achieves saturation
  must deviate from the spectral prescription by
  $\ge M_{\lfloor n/2\rfloor}^{\mathrm{deficit}}(n)$
  operators, breaking the correspondence between
  lattice geometry and measurement selection.
\end{itemize}
For $q \ge 3$, $K^*$ achieves spectral completeness with partial
weight coverage; the gap to full saturation ($K_{\mathrm{sat}}$)
is bridged by the compositional architecture
(verified computationally at $q = 3, 5, 7$).
\end{theorem}

The sharpness of part~(iii) for random measurement sets
is quantified by a coupon-collector bound:

\pclaim{env:lem:coupon}{proved}{*}{lem:coupon}%
\begin{lemma}[Coupon-collector threshold]
\label{lem:coupon}
Let $M$ operators be drawn uniformly without replacement from
the $4^n - 1$ non-identity Paulis.  For each weight class~$w$
with $A_w = \binom{n}{w}3^w$ elements:
(a)~If $M < A_w \ln A_w - A_w \ln\ln A_w$, then
  $\Pr[\text{weight $w$ saturated}] \to 0$ as $n \to \infty$.
(b)~The probability that a random $M$-element set
  saturates weight class~$w$ satisfies
  $\Pr[\mu_w = 1] = \prod_{j=0}^{A_w-1}(M{-}j)/(N{-}j)
  \le (M/N)^{A_w}$.
  When $M/N \le c < 1$, this bound vanishes as
  $A_w \to \infty$ (i.e., as $n \to \infty$ for
  fixed~$w$).
\end{lemma}
\begin{proof}
The $M$ operators are drawn without replacement from
$N = 4^n - 1$ non-identity Paulis, so the number of
weight-$w$ operators included follows a hypergeometric
distribution.  The probability that all~$A_w$ weight-$w$
operators are captured is
$\prod_{j=0}^{A_w-1}(M{-}j)/(N{-}j) \le (M/N)^{A_w}$,
since $(M{-}j)/(N{-}j) \le M/N$ for every $j \ge 0$.
When $M/N \le c < 1$ for some constant~$c$
(as holds for the $K^*$ budget at $q = 2$, $n = 2, \ldots, 9$),
this gives $\Pr[\mu_w = 1] \le c^{A_w} \to 0$, proving~(b).
Part~(a) is a corollary: $M < A_w \ln A_w$ is
far below the saturation threshold $M \approx N$,
so $\Pr[\mu_w = 1] \to 0$ a
fortiori~\cite{MitzenmacherUpfal2005}.
\end{proof}

\begin{proof}[Proof of Theorem~\ref{thm:spectral-char}]
\emph{Part~(i).}
Let $P_0 \in \mathcal{P}_n^{(w_0)}$ with $P_0 \notin
\mathcal{S}$, and define the Pauli-axis states
$\rho^\pm = (I \pm P_0)/d$.
For qubits ($q = 2$), every $n$-qubit Pauli has
eigenvalues $\pm 1$, so
$\rho^\pm$ has spectrum $\{0, 2/d\}$ (each with multiplicity
$d/2$); hence $\rho^\pm \ge 0$.  The only nonzero Bloch
coordinate is $\mathrm{tr}(P_0\,\rho^\pm) = \pm 1$ at
weight~$w_0$, so $W(\rho^\pm) = w_0 \le k$.

\emph{Indistinguishability.}
For any $P \in \mathcal{S}$, $P \neq P_0$ (since
$P_0 \notin \mathcal{S}$), and Pauli orthogonality gives
$\mathrm{tr}(P\,\rho^+) = \mathrm{tr}(P\,\rho^-) = 0$.
Both states produce identical measurement data on
$\mathcal{S}$.

\emph{Orthogonality.}
Since $\rho^+$ and $\rho^-$ commute and have disjoint
eigenvalue support ($\rho^+$ is supported on the $+1$
eigenspace of $P_0$, $\rho^-$ on the $-1$ eigenspace),
$F(\rho^+, \rho^-)
= \bigl(\sum_i \sqrt{p_i q_i}\bigr)^2 = 0$,
since $p_i q_i = 0$ for every eigenvalue pair.
Any estimator must assign the same reconstruction~$\hat\sigma$
to both states.  We show
$F(\hat\sigma,\rho^+) + F(\hat\sigma,\rho^-) \le 1$
for every state~$\hat\sigma$.
Let $\Pi_\pm$ be the projectors onto the $\pm 1$
eigenspaces of~$P_0$, so $\rho^\pm = (2/d)\,\Pi_\pm$
and $\Pi_+ + \Pi_- = I$.
Since $\Pi_\pm$ is a projector,
$\sqrt{\rho^\pm} = \sqrt{2/d}\;\Pi_\pm$.
By Uhlmann's theorem~\cite{Uhlmann1976},
$F(\hat\sigma,\rho^\pm)
= \bigl(\mathrm{tr}\sqrt{\sqrt{\hat\sigma}\,\rho^\pm
\sqrt{\hat\sigma}}\bigr)^2
= (2/d)\,\|\sqrt{\hat\sigma}\,\Pi_\pm\|_1^2$.
Since $\mathrm{rank}(\Pi_\pm) = d/2$, the Schatten-norm
rank inequality
$\|A\|_1 \le \sqrt{\mathrm{rank}(A)}\;\|A\|_F$ gives
$\|\sqrt{\hat\sigma}\,\Pi_\pm\|_1^2
\le (d/2)\,\|\sqrt{\hat\sigma}\,\Pi_\pm\|_F^2
= (d/2)\,\mathrm{tr}(\Pi_\pm\hat\sigma)$.
Summing:
$F(\hat\sigma,\rho^+) + F(\hat\sigma,\rho^-)
\le (2/d)\cdot(d/2)\bigl[\mathrm{tr}(\Pi_+\hat\sigma)
+ \mathrm{tr}(\Pi_-\hat\sigma)\bigr]
= \mathrm{tr}(\hat\sigma) = 1$.
Therefore $\min\{F(\hat\sigma,\rho^+),\,
F(\hat\sigma,\rho^-)\} \le 1/2$: no estimator achieves
worst-case fidelity exceeding~$1/2$.
For $q \ge 3$, the analogous construction uses
Heisenberg--Weyl shift-and-clock operators in place of
$\rho^\pm = (I \pm P_0)/d$: each generalized Pauli operator
$P_0$ has $q$ eigenspaces (one per character $\omega^j$,
$j = 0,\ldots,q-1$) of equal dimension $d/q$, and the
$q$ uniform-on-eigenspace states $\rho^{(j)} = (q/d)\Pi_j$
are pairwise orthogonal, mutually indistinguishable on
$\mathcal{S} \setminus \{P_0\}$, and yield
$\sum_j F(\hat\sigma, \rho^{(j)}) \le 1$ by the same
Schatten/rank argument applied to $q$ projectors of rank
$d/q$, giving worst-case fidelity $\le 1/q$
(Supplemental Material~\cite{supplemental}).

\emph{Delsarte dual certificate.}
For $q \le 3$, where $G(K)$ lies in the Bose--Mesner algebra
(Lemma~\ref{lem:spectral_q_main}),
the spectral obstruction has a structural origin in the
algebra of~$H(n,q)$~\cite{Delsarte1973}.
The Gram matrix decomposes as
$G(K) = \sum_w \lambda_w(K)\, E_w$ with
$\lambda_w = q^n c_w(K)/v_w$, where $c_w(K)$ counts
lattice vectors at parity weight~$w$ and norm~$\le K$.
When $c_{w_0}(K) = 0$, i.e., no lattice mode at
parity weight~$w_0$ has been resolved, the
primitive idempotent
\begin{equation}\label{eq:delsarte-cert}
  E_{w_0} = \frac{v_{w_0}}{q^n}
  \sum_{h=0}^{n}
  \frac{K_{w_0}(h;\, n,\, q)}{v_h}\, D_h
\end{equation}
(where $K_{w_0}(h;\, n,\, q)$ is the Krawtchouk
polynomial~\cite{Delsarte1973}, $v_h = \binom{n}{h}(q{-}1)^h$,
and $D_h$ denotes the distance-$h$ adjacency matrix
of~$H(n,q)$~\cite{BannaiIto1984}) satisfies
$G(K)\, E_{w_0} = \lambda_{w_0}\, E_{w_0} = 0$.
This is a \emph{spectral certificate of incompleteness}:
$E_{w_0}$ lies in $\ker G(K)$, so $G(K)$ cannot be
positive definite, and no estimator can resolve states
differing only in the $w_0$-th weight class.
The certificate is constructive and arises from the
Delsarte linear program adapted to measurement
design~\cite{Delsarte1973,BannaiIto1984}.
The \emph{primal} program asks: given $G(K)$,
maximize $\alpha$ subject to
$G(K) - \alpha\, E_{w_0} \succeq 0$.
Since all matrices lie in the Bose--Mesner algebra
(for $q \le 3$),
this semidefinite constraint reduces to the scalar
inequality $\lambda_{w_0}(K) \ge \alpha$, making the
program a linear program in the eigenvalue
variables~\cite{Delsarte1973}.
(For $q \ge 4$, the weight-block-diagonal structure yields
an analogous block-scalar reduction; see
Lemma~\ref{lem:spectral_q_main}.)  The \emph{dual} LP provides a certificate of
infeasibility: the idempotent $Y = E_{w_0}$ satisfies
$Y \ge 0$ (since $E_{w_0}$ is a projection onto a
subspace of~$\mathbb{R}^{q^n}$), $\langle Y, G(K) \rangle
= \lambda_{w_0}(K) \cdot \mathrm{rank}(E_{w_0}) = 0$
(when $c_{w_0} = 0$), and $\mathrm{tr}(Y) = v_{w_0} > 0$.
By LP strong duality, the primal optimum $\alpha = 0$:
the Gram matrix cannot resolve the $w_0$-th weight class.
This is the measurement-design analogue of
Delsarte's LP bound on code
size~\cite{Delsarte1973,BannaiIto1984}:
$E_{w_0}$ certifies spectral incompleteness of the
measurement set rather than bounding code distance.

\medskip\noindent
\emph{Part~(ii).}
Weight saturation ($\mu_w = 1$ for $w \le k$) gives
$\varepsilon_{\mathrm{flat}} = 0$ by
Theorem~\ref{thm:basin}(i).
Eigenvalue positivity ($\lambda_w > 0$ for all~$w$)
ensures support-completeness
(Supplemental Material, Corollary~1~\cite{supplemental}).
Applying Lemma~\ref{lem:purity_main} to~$\hat\rho$
gives $\varepsilon_{\mathrm{pos}} \le (d{-}1{-}S_k)/d^2$
(under weight saturation, all weight-${\le}\,k$ operators
are measured, so the MLE low-weight coordinates approximate
the true values under the Hoeffding event~\eqref{eq:hoeffding};
the resulting error $|S_k - \hat{S}_k| \le M t$ is
$O(\sqrt{M\ln(2M/\delta)/N_s})$, which is absorbed by the
shot-noise term in~\eqref{eq:spectral-resolution}).
Combining via~\eqref{eq:fidelity-decomp}:
$\|\rho - \hat\rho_{\mathcal{S}}\|_F^2 \le d\,\varepsilon_{\mathrm{flat}} +
d\,\varepsilon_{\mathrm{pos}} + \tfrac{8(1+t_\star^2)}{(1-t_\star)^2}\,M\ln(2M/\delta)/(d\,N_s)
= 0 + (d{-}1{-}S_k)/d + \tfrac{8(1+t_\star^2)}{(1-t_\star)^2}\,M\ln(2M/\delta)/(d\,N_s)$,
which is~\eqref{eq:spectral-resolution}.

\medskip\noindent
\emph{Part~(iii).}
The dichotomy is immediate from (i) and (ii): if
$\mu_{w_0} < 1$ for any $w_0 \le k$, part~(i) gives
worst-case $F \le 1/2$; if $\mu_w = 1$ for all $w \le k$,
part~(ii) gives $\|\rho - \hat\rho_{\mathcal{S}}\|_F^2 = O(d^{-1}) + O(d/N_s)$.
The HS error $O(d^{-1}) \to 0$ as $d\to\infty$.
For random measurement sets, the transition is sharp
(Lemma~\ref{lem:coupon}).

\medskip\noindent
\emph{Part~(iv): spectral progression.}
For $q \le 3$, the Gram matrix lies in the Bose--Mesner algebra of
$H(n,q)$ (Lemma~\ref{lem:spectral_q_main}(i)):
$G(K) = \sum_w \lambda_w(K)\, E_w$ with
$\lambda_w(K) = q^n c_w(K) / v_w$, where $v_w =
\binom{n}{w}(q{-}1)^w$.
For $q \ge 4$, the weight-block-diagonal structure
(Lemma~\ref{lem:spectral_q_main}(ii)) yields the same
parity-weight counts $c_w(K)$ and the same monotonicity argument.
Since $c_w$ is non-decreasing in $K$ (adding lattice shells
can only increase counts), the per-weight budget $M_w(K)$
is non-decreasing.  We establish insufficiency at each
cutoff below~$K^*$.

\emph{Phase~1: spectral activation ($K = 1, \ldots, n{-}1$).}
At cutoff~$K$, parity weight~$w$ requires at least $w$
odd-valued coordinates, each contributing $\ge 1$ to $|m|^2$,
so $|m|^2 \ge w$.  Hence $c_{K+1}(K) = 0$ and the idempotent
$E_{K+1}$~\eqref{eq:delsarte-cert} lies in $\ker G(K)$:
a Delsarte dual certificate of spectral incompleteness.
No estimator can resolve states differing in
weight class~$K{+}1$.

\emph{Phase~2: spectral completeness without saturation
($K = n, \ldots, K^*{-}1$).}
At $K = n$ (specifically $K = 4$ for $n = 4$, $q = 2$),
all weight classes have $c_w > 0$ (the vectors
$(\pm 1, \ldots, \pm 1)$ contribute $c_n = 2^n$).
The Gram matrix $G(n)$ has full rank:
no Delsarte null certificate exists.
However, the $89$ lattice vectors with $|m|^2 \le 4$
decompose by parity weight as
$c_w(4) = \{9, 8, 24, 32, 16\}$ for $(n,q) = (4,2)$
(verified by direct enumeration;
Supplemental Material~\cite{supplemental}).
After greedy redistribution: $M_w = \{1, 12, 28, 32, 16\}$.
Weight~$2$ achieves coverage $28/54 = 52\% < 100\%$,
so $\mu_2 < 1$ and part~(i) applies:
worst-case $F \le 1/2$ for $2$-local states.
Note that $\sum c_w = 89 > 54 = A_2$, so weight-$2$
saturation is achievable by non-spectral operator
selection, but requires abandoning the Krawtchouk
eigenvalue allocation.  The correspondence between
lattice geometry and measurement design has not yet
converged.
For $q = 2$ this phase contains only $K = n = 4$.
For $q \ge 3$, the range $n < K < q^2$ is nonempty and
the same mechanism applies at each cutoff: the
Krawtchouk budget grows monotonically with~$K$
(Lemma~\ref{lem:monotone}) but the mod-$q$
invisible vectors that generate weight-$0$ overflow
first appear at $|m|^2 = q^2$, so
redistribution remains insufficient
until $K = K^*$ (verified computationally for
$q = 3, 5, 7$ in the Supplemental
Material~\cite{supplemental}).

\emph{Phase~3: convergence at $K^*$.}
At $K^* = 5$, the $|m|^2 = 5$ shell contributes~$48$
lattice vectors of parity weight~$1$
(the $(\pm 1, \pm 2, 0, 0)$ family),
pushing $c_1$ from~$8$ to~$56$.
After redistribution: $c_0 = 9$ yields excess~$8$
(only~$1$ identity operator);
$c_1 = 56$ plus the excess gives~$64$, of which~$12$
fill weight~$1$, leaving excess~$52$.
At weight~$2$: $c_2 = 24$ plus excess~$52$ gives~$76$,
of which~$54 = A_2$ fill weight~$2$.
The resulting budget $M_w = \{1, 12, 54, 54, 16\}$
saturates weights~$0$--$2$.
No reallocation is needed: the spectral prescription
and weight saturation coincide.
By the semisimplicity of the Bose--Mesner
algebra~\cite{BannaiIto1984}, the eigenvalue
decomposition is unique, so $K_{\mathrm{sat}}$~\eqref{eq:ksat-def}
is the unique threshold at which spectral geometry
and operational performance converge.
\end{proof}

\begin{remark}[Tightness status]
Concretely, at fixed $n$, $N_s = 10{,}000$ achieves a
high-probability shot-noise term of~$0.070$
at $\delta = 0.01$ (Theorem~\ref{thm:basin}(iii)).
The Uhlmann bound of part~(i) is \emph{tight}: the
Pauli-axis construction achieves $F = 0$ between the
two adversarial states, matching the bound with equality.
The reconstruction bound of part~(ii) is \emph{not known to be
tight}: closing the gap requires a matching lower bound
on Hilbert--Schmidt error, which we leave as an open problem.
Part~(iv) is tight in two senses:
the Delsarte certificates at $K \le n{-}1$ are exact
(zero eigenvalue, not merely small);
and the correspondence failure at $K = n$ is explicit
(deficit of~$26$ operators at weight~$2$).
The eigenvalue condition number satisfies
$\kappa(G(K^*)) = \lambda_{\max}/\lambda_{\min}
= 256/64 = 4$ at $n = 4$, $q = 2$.
More generally:
\begin{equation}\label{eq:kappa}
  \kappa(G(K^*)) = \frac{\max_w c_w(K^*)/v_w}
  {\min_w c_w(K^*)/v_w}
\end{equation}
where $v_w = \binom{n}{w}(q{-}1)^w$.
At the operational $K^*$ (one shell beyond spectral
completeness), $\kappa = 4$ for $n = 2, 3, 4, 5$
and $\kappa \approx 4.9$ for $n = 6$
(Supplemental Material~\cite{supplemental}),
consistent with at most logarithmic growth in~$n$.
The Krawtchouk prescription thus produces a
well-conditioned measurement design.
Whether intra-weight reweighting can
improve~$\kappa$ at fixed~$K^*$ is an open question.
\end{remark}

\pclaim{env:cor:lower_bound}{proved}{q>=2}{cor:lower_bound}%
\begin{corollary}[Operator lower bound]
\label{cor:lower_bound}
Any measurement set $\mathcal{S}$ achieving worst-case fidelity
$F > 1/2$ over the class $\mathcal{W}_k$ of weight-$k$ supported states
must satisfy $\mu_w(\mathcal{S}) = 1$ for every $w = 1,\ldots,k$.
In particular:
\begin{equation}\label{eq:lower_bound}
  |\mathcal{S}| \ge
  \sum_{w=1}^{k}\binom{n}{w}(q^2 - 1)^w.
\end{equation}
At $n = 4$, $q = 2$, $k = 2$: $|\mathcal{S}| \ge 66$.
The $K^*$ prescription achieves this with $M = 137$ operators, of
which $67$ are weight~${\le}\,2$ ($66$~necessary plus identity) and
$70$ are weight~$> 2$ (from the Krawtchouk lattice structure).
\end{corollary}
\begin{proof}
If $\mu_{w_0} < 1$ for any $w_0 \le k$, then by
Theorem~\ref{thm:spectral-char}(i) there exist weight-$k$ supported
states $\rho_1, \rho_2$ that are indistinguishable on
$\mathcal{S}$ yet satisfy $F(\rho_1, \rho_2) = 0$.
No estimator achieves worst-case $F > 1/2$,
contradicting the hypothesis.  Hence $\mu_w = 1$ for all
$w \le k$, which requires $|\mathcal{S}|
\ge \sum_{w=1}^{k} A_w$.
\end{proof}

\pclaim{env:thm:asymptotic}{proved}{*}{thm:asymptotic}%
\begin{theorem}[Asymptotic $K^*$ separation]
\label{thm:asymptotic}
Fix $k \ge 1$ and $q \ge 2$.
Let $M_n = N_n(K^*_n)$ be the $K^*$-prescription operator
count for $n$-qubit systems, and assume
$M_n/(q^{2n}{-}1) \le c < 1$ for some constant~$c$
independent of~$n$ (verified numerically at
$q = 2$ for $n = 2, \ldots, 9$).

\textup{(i)} \emph{Deterministic saturation.}
The $K^*$ prescription achieves weight saturation
$\mu_w = 1$ for all $w \le k$, for every~$n$
(monolithically for $n \le 5$; via the compositional
architecture of Sec.~\ref{sec:8qubit} for $n \ge 6$).

\textup{(ii)} \emph{Random failure.}
Let $\mathcal{S}_{\mathrm{rand}}$ be a uniform random
subset of size $M_n$ drawn from the $q^{2n} - 1$
non-identity Pauli operators.  For each fixed
$w \ge 1$:
\begin{equation}\label{eq:random_failure}
  \Pr\!\bigl(\mu_w(\mathcal{S}_{\mathrm{rand}}) = 1\bigr)
  \le \biggl(\frac{M_n}{q^{2n} - 1}\biggr)^{\!A_w}
  \to 0 \quad\text{as } n \to \infty.
\end{equation}

\textup{(iii)} \emph{Fidelity dichotomy.}
Combining with Theorem~\ref{thm:spectral-char}:
as $n \to \infty$, the $K^*$ prescription (via
$K_{\mathrm{sat}}$) achieves
$\|\rho - \hat\rho_{\mathcal{S}}\|_F^2 = O(d^{-1}) + O(d/N_s)$
for all pure weight-$k$ supported states,
while a same-budget random set yields
worst-case $F \le 1/2$ asymptotically almost surely.
\end{theorem}

\begin{proof}
\emph{Part~(i).}
By construction: the $K^*$ redistribution at cutoff
$K_{\mathrm{sat}}(n)$ achieves $M_w = A_w$ for all $w \le k$
(Theorem~\ref{thm:spectral-char}(iv),
Eq.~\eqref{eq:ksat-def}).
For $q = 2$ and $n \le 5$, $K_{\mathrm{sat}} = K^* = 5$;
for larger~$n$, the compositional architecture reduces
each patch to $n = 4$, where saturation holds.

\emph{Part~(ii).}
This is the coupon-collector argument of
Lemma~\ref{lem:coupon}.  Drawing $M_n$ operators without
replacement from $N = q^{2n} - 1$, the probability that
all $A_w$ weight-$w$ operators are included is
\begin{equation}
  \Pr(\mu_w = 1)
  = \prod_{j=0}^{A_w - 1} \frac{M_n - j}{N - j}
  \le \Bigl(\frac{M_n}{N}\Bigr)^{A_w},
\end{equation}
since $(M_n - j)/(N - j) \le M_n / N$ for all
$j \ge 0$.  By hypothesis, $M_n/N \le c < 1$.
Since $A_w = \binom{n}{w}(q^2{-}1)^w \to \infty$ for
fixed $w \ge 1$, the bound $c^{A_w} \to 0$.

\emph{Part~(iii).}
From part~(ii), $\Pr(\mu_w < 1) \to 1$ for some
$w \le k$.  By Theorem~\ref{thm:spectral-char}(i),
worst-case $F \le 1/2$ for the random set.
Part~(i) and Theorem~\ref{thm:spectral-char}(ii) give
the $K^*$ guarantee.
\end{proof}

\begin{remark}[$M_n/N$ bound]
\tolerance=1500
The condition $M_n/(q^{2n}{-}1) \le c < 1$ holds for all $n \ge 2$
at $q = 2$.  At $K^* = 5$, every lattice point has at most $5$
nonzero coordinates (each contributing $\ge 1$ to $|m|^2$),
so $c_w(5) = 0$ for $w \ge 6$ and
\begin{equation*}
  M_n = 1 + 4n + 12\tbinom{n}{2} + 8\tbinom{n}{3}
        + 16\tbinom{n}{4} + 32\tbinom{n}{5}
\end{equation*}
is a degree-$5$ polynomial in $n$, while $N = 4^n - 1$ is
exponential.  Hence $M_n/N \to 0$ as $n \to \infty$.
For $n \le 3$ the protocol uses $K = n$ (Recipe above),
yielding even fewer operators; the worst case is
$M_2/N_2 = 9/15 = 0.60$ (at $K = 2$), so
$c = 0.60$ satisfies the hypothesis unconditionally.
Numerically: $M_n/N = 0.54$ at $n = 4$,
$0.06$ at $n = 8$, $0.002$ at $n = 12$
(Supplemental Material~\cite{supplemental}).
\end{remark}

\begin{remark}[Linear indistinguishability]
\tolerance=1500
Theorem~\ref{thm:spectral-char}(i) implies that no
state-independent linear criterion can detect weight-class
structure: Pauli orthogonality makes all $M$-element
subsets identical in Fisher eigenvalues, frame potential,
and coherence
(Appendix~\ref{sec:fisher}).
The empirical $\Delta F \approx +0.03$ under linear
inversion versus $\Delta F \approx +0.33$ under MLE
(Supplemental Material~\cite{supplemental}, Sec.~VII\,F) reflects
the ${\sim}11\times$ amplification by positivity-constrained
reconstruction across the dichotomy of
part~(iii).
\end{remark}

\subsection{Universality across Hamming schemes}
\label{sec:universality}

The spectral decomposition extends from qubits ($q = 2$) to
qudits, with a structural boundary at $q = 3$.

\pclaim{env:lem:spectral_q_main}{proved}{*}{lem:spectral_q_main}%
\begin{lemma}[Spectral decomposition and Euclidean--Hamming compatibility]
\label{lem:spectral_q_main}
Let $G(K)$ be the Gram matrix of the Hamming scheme $H(n,q)$
at Euclidean cutoff~$K$.\tolerance=1500
\begin{enumerate}
\item[\emph{(i)}] \emph{(Iff characterization.)}
$G(K)$ lies in the Bose--Mesner algebra of $H(n,q)$ if and only
if $q \le 3$.  When this holds, the Krawtchouk spectral
decomposition gives
\begin{equation}\label{eq:spectral_q_main}
  G(K) = \sum_{w=0}^{n} \lambda_w(K)\, E_w, \quad
  \lambda_w = \frac{q^n \, c_w(K)}{\binom{n}{w}(q{-}1)^w},
\end{equation}
where $E_w$ are the primitive idempotents of $H(n,q)$ and
$c_w(K) = |\{m \in \mathbb{Z}^n : |m|^2 \le K,\;
w(\bar{m}) = w\}|$ counts lattice vectors by $q$-ary parity
weight.
\item[\emph{(ii)}] \emph{(Weight-block structure, all $q$.)}
For every~$q$, coordinate-permutation invariance
($S_n$-symmetry) makes $G(K)$ block-diagonal in Hamming
weight: the parity-weight counts $c_w(K)$ and the
weight-class saturation budget $M_w(K)$ are well-defined
for all~$q$.
\item[\emph{(iii)}] \emph{(Spectral completeness and mass threshold.)}
Spectral completeness ($c_w > 0$ for all~$w$, equivalently
support-completeness; Lemma~\ref{lem:completeness}) is achieved
at $K = n$ for all~$q$.
At $K_{\mathrm{mass}} = q^2$ the spectral mass threshold
(proved for all $q \ge 2$;
Supplemental Material~\cite{supplemental},
Proposition ``Spectral mass threshold'')
is the minimal cutoff at
which the mod-$q$ invisible vectors $(\pm q, 0, \ldots, 0)$
enter, increasing $c_0$ from~$1$ to $1 + 2n$.
For $q = 2$, $K^* = 5 = q^2 + 1$.
For $q \ge 3$, full weight
saturation~\eqref{eq:ksat-def} requires $K > q^2$
(e.g., $K_{\mathrm{sat}} = 14$ at $q = 3$, $n = 4$;
Supplemental Material~\cite{supplemental}).
\end{enumerate}
\end{lemma}
\begin{proof}
\emph{(i), forward.}
For $q \le 3$, every nonzero residue $a \in \mathbb{Z}_q
\setminus \{0\}$ has minimal symmetric lift $\ell(a)$ with
$|\ell(a)|^2 = 1$ (for $q = 2$: $\ell(1) = \pm 1$; for
$q = 3$: $\ell(1) = 1$, $\ell(2) = -1$).
All nonzero residues therefore contribute identically to the
Euclidean cutoff $|m|^2 \le K$, so $G(K)$ is constant on
each Hamming distance class and lies in the Bose--Mesner algebra.
\emph{Converse.}
For $q \ge 4$, residues $1$ and $2$ lift to
$|\ell|^2 = 1$ and $|\ell|^2 = 4$ respectively.
At $n = 1$, $K = 1$, the off-diagonal entries
$G_{0,1} = 1 + 2\cos(2\pi/q) \neq G_{0,2} = 1 + 2\cos(4\pi/q)$,
so $G(K) \notin \mathrm{BM}(H(1,q))$.
Since the lift-norm discrepancy is a property of~$\mathbb{Z}_q$
(independent of~$n$), the $n = 1$ submatrix embeds in every
$H(n,q)$ with $n \ge 1$ and $K \ge 1$, so the failure
persists for all~$n$.
\emph{(ii)} follows from $S_n$-invariance of the Euclidean norm.
\emph{(iii)} is verified by explicit lattice enumeration
(Supplemental Material~\cite{supplemental}, Proposition~1 therein,
which restates this result with the full $q$-ary eigenvalue formula).
\end{proof}

Basin separation (Theorem~\ref{thm:basin}) extends
to qudit systems for \emph{all}~$q$: the proof uses only operator
orthogonality and the purity constraint, both of which hold for
any orthonormal Hermitian operator basis.
Basin separation requires only the
weight-block-diagonal structure of the Gram matrix
($S_n$-invariance), which holds for all~$q$; it does not require
full Bose--Mesner membership.
Identifiability therefore applies to all~$q$, even
though the Krawtchouk diagonalization is restricted to $q \le 3$
(Lemma~\ref{lem:spectral_q_main}).
For $q \ge 3$, the natural character basis
(Heisenberg--Weyl operators) is non-Hermitian; practical qudit
tomography uses the Hermitian generalized Gell-Mann basis instead.
Gell-Mann weight-class sizes are $(q^2{-}1)^w \binom{n}{w}$
rather than the Hamming-scheme multiplicities $(q{-}1)^w \binom{n}{w}$
that appear in Lemma~\ref{lem:spectral_q_main}.
For $q \le 3$, the spectral decomposition and saturation thresholds
(Lemma~\ref{lem:spectral_q_main}) are stated in the character
basis; the operational translation to Gell-Mann coordinates
preserves orthogonality and the purity bound, so the
identifiability argument (Theorem~\ref{thm:basin}(i))
applies with identical logical structure, but the
weight-class counts
and saturation budgets must be recalculated in the Gell-Mann basis
(Supplemental Material~\cite{supplemental}, where this is done
explicitly for $q = 3$).
For $q \ge 4$, the Krawtchouk eigenvalues are not available, but
the weight-block structure still determines the saturation budget
per weight class; the basin separation argument applies unchanged.
For $q \ge 3$, the compositional architecture
(overlapping $d^*$-qudit patches) compensates for the
slower coverage growth at higher~$q$.

%% ========================================================
\section{Experimental setup}
\label{sec:experiment}

The preceding theory makes concrete predictions: structured
selection should outperform random whenever the measurement budget
is underdetermined, with the advantage driven by weight-class
saturation.  We now test these predictions on superconducting
hardware.

All experiments use qubits ($q = 2$); the qudit extension
is verified computationally through $q = 7$
(Sec.~\ref{sec:universality}).

IBM tomography experiments were performed on \texttt{ibm\_fez}
(156-qubit Eagle~r3)~\cite{Kim2023} via the IBM Quantum Platform,
using qubits 140--147 on a linear chain with reported controlled-Z (CZ) error rates
of 0.15--0.19\%.
The Floquet DTC application (Sec.~\ref{sec:dtc}) used
\texttt{ibm\_kingston} (156-qubit Heron~r2).
Multi-vendor replication used a Rigetti Ankaa-3 processor
(84-qubit, square lattice, iSWAP native gate) via the Open Quantum
cloud platform, with controlled-$R_Y$ gates decomposed
into $u_3{+}\text{CX}$ (controlled-NOT) sequences.

State-preparation-and-measurement (SPAM) errors are common to both arms (identical state preparation,
qubits, and computational-basis measurement); only single-qubit
basis rotations differ, so $\Delta F$ isolates operator selection.

For each state, we executed two experiments with identical parameters:
\begin{itemize}
\item \textbf{Structured}: $M$ Pauli operators from the Krawtchouk
prescription ($M = 137$ for $n = 4$; $M = 9$ for $n = 2$).
\item \textbf{Random}: $M$ Pauli operators selected uniformly at
random (fixed seed for reproducibility).
\end{itemize}
Although the 29 structured bases yield $196$ nontrivial Pauli
expectations (Sec.~\ref{sec:method}), only the $M = 137$
Krawtchouk-selected operators are passed to the MLE for both arms,
ensuring that the comparison isolates operator \emph{selection},
not operator \emph{count}.

Each Pauli measurement circuit consists of state preparation,
single-qubit basis rotations ($H$ for $X$, $S^\dagger H$ for $Y$,
identity for $Z$), and computational-basis measurement.  All circuits
were transpiled with optimization level~1 and executed with 1{,}000~shots.
No error mitigation of any kind was applied~\cite{Schwemmer2015} --
no readout correction, no dynamical decoupling, no randomized
compiling, no zero-noise extrapolation, and no probabilistic error
cancellation; both arms use identical raw counts.
Circuits were constructed using Qiskit~2.3 (Python~3.13).
Qubits 140--147 were selected for connectivity, not calibration
quality; four-qubit experiments used qubits 140--143 (the first
four in the chain).  Rigetti qubit identifiers were assigned by
the Open Quantum scheduler and are not user-selectable on
cloud-access platforms; replication on these independent qubits
corroborates the effect on a second architecture.

\paragraph{2-qubit state.}
Bell pair $|\Phi^+\rangle = (|00\rangle + |11\rangle)/\sqrt{2}$
(1~CX), measured with $M = 9$ Pauli operators (all lattice modes
at $n = 2$, $K = 2$).

\paragraph{4-qubit states.}
(i)~Product state $|{+}\rangle^{\otimes 4}$ (0~CX
gates), (ii)~4-qubit W~state
$|W_4\rangle = \tfrac{1}{2}(|1000\rangle + |0100\rangle +
|0010\rangle + |0001\rangle)$ (6~CX), and (iii)~4-qubit GHZ state
$|GHZ_4\rangle = \tfrac{1}{\sqrt{2}}(|0000\rangle + |1111\rangle)$
(3~CX).

\paragraph{Dimensional sweep.}
A same-session run compared W~states at $n = 2$ ($M = 9$) and
$n = 4$ ($M = 137$) under identical calibration conditions
(Sec.~\ref{sec:dim-dep}, Table~\ref{tab:results}).

%% ========================================================
\section{Results}
\label{sec:results}

\subsection{Head-to-head comparison}
\label{sec:head-to-head}

Table~\ref{tab:results} and Figure~\ref{fig:fidelity} summarize the
reconstruction fidelity
$F(\rho_{\rm ideal}, \rho_{\rm reconstructed})$ for each state and
method.

\begin{table}[tb]
\centering
\caption{Reconstruction fidelity: structured vs.\ random Pauli selection
on superconducting hardware.  $M$ denotes the number of Pauli
\emph{operators} measured; under tensor-product basis grouping,
the circuit count is smaller (e.g., $137$ operators $\to 29$
$K^*$ bases or $\sim 50$ random bases at $n = 4$).
Each circuit uses $1{,}000$~shots.
Fidelity is Uhlmann fidelity
$F(\rho,\sigma) = (\mathrm{tr}\sqrt{\sqrt{\rho}\,\sigma\sqrt{\rho}})^2$.}
\label{tab:results}
\resizebox{\columnwidth}{!}{%
\begin{tabular}{lccccc}
\toprule
State & $n$ & $M$ & Struct.\ $F$ & Rand.\ $F$ & $\Delta F$ \\
\midrule
$|{+}\rangle^{\otimes 4}$ & 4 & 137 & 0.996 & 0.980 & $+$0.016 \\
Bell pair & 2 & 9 & 0.518 & 0.465 & $+$0.052 \\
W$_4$ state$^\ddagger$ & 4 & 137 & \textbf{0.872}$\pm$\textbf{0.021}
  & \textbf{0.54}$\pm$\textbf{0.08} & $+$\textbf{0.33}$\pm$\textbf{0.07} \\
GHZ$_4$ state$^\S$ & 4 & 137 & 0.496 & 0.498 & $-$0.002 \\
\midrule
W$_2$ state$^\dagger$ & 2 & 9 & 0.983 & 0.522 & $+$0.461 \\
W$_4$ rerun$^\dagger$ & 4 & 137 & 0.876 & 0.441 & $+$0.435 \\
\midrule
\multicolumn{6}{l}{\emph{Rigetti Ankaa-3 (Open Quantum)}} \\
W$_4$ state$^\|$ & 4 & 137 & 0.815 & 0.568 & $+$0.248 \\
W$_3$ state$^\|$ & 3 & 27 & 0.840 & 0.831 & $+$0.009 \\
\midrule
\multicolumn{6}{l}{\emph{Rigetti Ankaa-3, 8-qubit compositional}} \\
$W_8$ (mean)$^\#$ & 8 & 137 & \textbf{0.668}$\pm$\textbf{0.07}
  & \textbf{0.134}$\pm$\textbf{0.07} & $+$\textbf{0.53}$\pm$\textbf{0.06} \\
$|{+}\rangle^{\otimes 8}$$^\#$ & 8 & 137 & 0.992$\pm$0.002 & -- & -- \\
\bottomrule
\end{tabular}}% end resizebox
\vspace{2pt}
{\footnotesize Noise-robust MLE reconstruction
(Sec.~\ref{sec:reconstruction}), $1{,}000$~shots per circuit
(each operator appears in exactly one basis).
$^\ddagger$Four independent runs (seeds 100--400);
$\pm$ denotes sample standard deviation.
$^\dagger$Same-session dimensional sweep (different random seed).
$^\S$Full-state MLE; automatic reconstructor switching recovers
$F = 0.887$; subspace least-squares gives $F = 0.926$
(Sec.~\ref{sec:ghz}).
$^\|$Same-session basis-grouped comparison
(Sec.~\ref{sec:rigetti}).
$^\#$Compositional architecture, 5 overlapping 4-qubit patches;
$\pm$ denotes standard deviation across patches
(Sec.~\ref{sec:8qubit}; Fig.~\ref{fig:8qubit}).}
\end{table}

\begin{figure}[tb]
\centering
\includegraphics[width=\columnwidth]{fig1_fidelity.pdf}
\caption{Reconstruction fidelity for structured (Krawtchouk $K^*$)
versus random Pauli measurement selection on IBM \texttt{ibm\_fez}.
4-qubit states use $M = 137$ operators ($29$ bases) at 1{,}000 shots per basis.
W-state bars show four-run means; error bars are standard deviations
across runs ($\pm 0.021$ structured, $\pm 0.08$ random).
Other states show single-run bootstrap uncertainties ($1{,}000$~resamples).
The W-state advantage $\Delta F = +0.33 \pm 0.07$ is the largest
observed effect.}
\label{fig:fidelity}
\end{figure}

Structured selection exceeds random for the
product and W~states, consistent with the basin separation
prediction (Theorem~\ref{thm:basin}): the $K^*$ set saturates
all weight classes up to $w = 2$, so every informative Pauli
coordinate is constrained.
The advantage is largest for the
W~state, the most entangled state with nontrivial local structure.
Four independent hardware repetitions with different random seeds
confirm the W-state result: $F(K^*) = 0.872 \pm 0.021$ (stable)
versus $F(\mathrm{rand}) = 0.54 \pm 0.08$ (volatile),
$\Delta F = +0.33 \pm 0.07$ (all four runs positive).
For the Bell pair ($\Delta F = +0.052$, not significant), an
$n = 2$ MLE pathology where the identity operator dominates variance
weights (Supplemental Material~\cite{supplemental}).
RDM decomposition corroborates the mechanism predicted by
Theorem~\ref{thm:basin}(i): $K^*$ gives near-perfect
1-RDM fidelity ($0.985 \pm 0.004$) from weight-$1$ saturation and
a 2-RDM advantage ($+0.097$) from weight-$2$ saturation
(Table~\ref{tab:rdm}; Appendix~\ref{sec:rdm}).

GHZ gives $\Delta F = -0.002$ under full-state MLE,
consistent with zero: full-state MLE reconstructs all $d^2 - 1$
Pauli coordinates simultaneously, so the GHZ signal
(concentrated on $6$ weight-$2$ $ZZ$-type and up to $9$ weight-$4$
correlators) is diluted by the remaining ${\sim}\,122$ near-zero
$K^*$ operators.  However, the $K^*$ data contains the GHZ signal:
a weight-class fingerprint classifier automatically detects the
information-budget mismatch (parity ratio $r = 0.010$) and switches
to subspace estimation, recovering $F = 0.887$ from the same
$137$~operators at zero additional cost (Sec.~\ref{sec:ghz}).
Raw phase correlators confirm ${\sim}\,85$\% coherence
($\langle X^{\otimes 4}\rangle = +0.864$).  Full-state MLE
deficit reflects primarily the reconstructor rather than the
measurement set (analyzed in detail in Sec.~\ref{sec:ghz}).

Uncertainties (``$\pm$'') are sample standard deviations across four
independent runs, not standard error of the mean (SEM).  All four per-run W-state advantages are
positive ($\Delta F = +0.33 \pm 0.07$; $t(3) = 9.3$,
$p < 0.003$; $95$\% CI $[0.22, 0.44]$); run~3's smaller gap
($+0.224$) reflects $IIII$ (the all-identity Pauli string $I^{\otimes 4}$) contamination in its random draw
(Supplemental Material~\cite{supplemental}).
Per-run shot-noise bootstraps ($10^3$~resamples each) give
$\sigma_{K^*}^\mathrm{boot} \approx 0.019$ and
$\sigma_\mathrm{rand}^\mathrm{boot} \approx 0.062$, consistent
with the inter-run dispersions and confirming that the volatility
difference is shot-noise dominated (Supplemental Material~\cite{supplemental}).

\subsection{Weight-allocated baseline}
\label{sec:weight-alloc}

A natural objection to the structured advantage is that a practitioner
who simply saturates the low-weight operators and fills the remaining
budget from higher weights, without any spectral theory, would achieve
comparable performance.
We test this directly
with a three-arm experiment, first in simulation
(Table~\ref{tab:weight-alloc}) and then on hardware
(Table~\ref{tab:weight-alloc-hw}).

A \emph{weight-allocated random} baseline uses the same per-weight
budget $M_w = \{1, 12, 54, 54, 16\}$ as the structured set, but
selects operators \emph{randomly} within each weight class.
A separate \emph{uniform random} baseline draws $137$
operators uniformly from the $255$ nontrivial Paulis (excluding IIII).
Because the structured set includes the identity operator as its
single weight-$0$ element while the uniform random pool does not,
the budget definitions differ by one operator; at $n = 4$ this
amounts to $< 1$\% of the total and does not measurably affect
reconstruction quality (Supplemental Material~\cite{supplemental}).
Both use the same MLE reconstruction
($\zeta = 0.05$, $\alpha = 0.5$).

\begin{table}[tb]
\centering
\caption{Three-arm simulation across states and dimensions.  ``Structured'' is
the Krawtchouk $K^*$ set; ``Alloc.\ random'' uses the same per-weight
budget with random intra-weight selection (20~seeds); ``Uniform random''
draws 137 operators uniformly (20~seeds).  $\Delta(S{-}AR)$ isolates
intra-weight ordering; $\Delta(AR{-}UR)$ isolates weight allocation.
Simulation: depolarizing noise $p = 0.03$, readout error $2$\%, $1{,}000$~shots.}
\label{tab:weight-alloc}
\begin{tabular}{lcccc}
\toprule
State & Struct.\ $F$ & Alloc.\ $F$ & Unif.\ $F$
  & $\Delta(S{-}AR)$ \\
\midrule
W($n{=}5$)  & 0.730 & 0.688 & 0.428 & $+0.042$ \\
W($n{=}4$)  & 0.833 & 0.733 & 0.302 & $+0.100$ \\
GHZ($n{=}4$) & 0.856 & 0.574 & 0.227 & $+0.282$ \\
W($n{=}3$)  & 0.868 & 0.538 & 0.284 & $+0.331$ \\
W($n{=}2$)  & 0.987 & 0.414 & 0.479 & $+0.573$ \\
\bottomrule
\end{tabular}
\end{table}

\begin{table}[tb]
\centering
\caption{Three-arm \emph{hardware} validation of the weight-allocation
mechanism (Theorem~\ref{thm:spectral-char}(iv)) on \texttt{ibm\_fez}
(4-qubit W~state, three seeds, $1{,}000$~shots per circuit).
Circuit execution order is reversed: weight-allocated random first,
uniform random second, structured last, ruling out
calibration-drift confounds.  The uniform random pool
excludes IIII.  The allocation fraction is
$(F_\mathrm{AR} - F_\mathrm{UR})/(F_{K^*} - F_\mathrm{UR})$.
Uniform random's high volatility ($\sigma = 0.325$,
coefficient of variation ${>}\,1$) makes this ratio unstable;
$\Delta(S{-}AR)$ is the more reliable metric.}
\label{tab:weight-alloc-hw}
\resizebox{\columnwidth}{!}{%
\begin{tabular}{lccccc}
\toprule
Seed & $F(K^*)$ & $F(\mathrm{AR})$ & $F(\mathrm{UR})$
     & $\Delta(S{-}AR)$ & Alloc.\ \%$^\dagger$ \\
\midrule
100 & 0.886 & 0.627 & 0.039 & $+0.259$ & 69\% \\
200 & 0.917 & 0.886 & 0.630 & $+0.031$ & 89\% \\
300 & 0.905 & 0.654 & 0.099 & $+0.252$ & 69\% \\
\midrule
Mean & $0.903{\pm}0.016$ & $0.722{\pm}0.142$ & $0.256{\pm}0.325$
     & $+0.181$ & $69$--$89$\% \\
\bottomrule
\end{tabular}}% end resizebox

{\footnotesize $^\dagger$Ill-conditioned; see caption.}
\end{table}

In simulation, for W~states at $n = 4$, $\Delta(S{-}AR) = +0.100$: the
structured set retains an advantage over weight-allocated random,
and weight allocation alone accounts for most
($\Delta(AR{-}UR) = +0.431$, or ${\sim}\,81$\% of the
total structured advantage).
Replacing the $K^*$-prescribed weight-$3$/$4$ split ($54{:}16$)
with a proportional split ($40{:}30$, distributing by pool size)
degrades W-state fidelity by $-0.21$ (Supplemental Material~\cite{supplemental}),
confirming that the Krawtchouk allocation
(Theorem~\ref{thm:spectral-char}(iv)) is not a
na\"ive saturation strategy.

On hardware (Table~\ref{tab:weight-alloc-hw}),
a separate three-arm session
corroborates the \emph{sign} of this decomposition: the per-seed residuals
$\Delta(S{-}AR) \in \{0.031, 0.252, 0.259\}$ are positive for all
seeds, consistent with the simulation estimate of
$+0.100$ (though larger in magnitude, likely reflecting the
hardware noise topology that the uniform depolarizing simulation
underestimates).  Three hardware seeds establish the direction of the
structured advantage beyond weight allocation, but the order-of-magnitude
spread ($+0.031$ to $+0.259$) precludes a precise quantitative
decomposition; additional seeds would be needed to narrow the hardware
estimate of $\Delta(S{-}AR)$.  (The wide range $69$--$89$\% reflects the uniform random arm's
volatility; $\Delta(S{-}AR)$ is the more stable metric.)
Both sessions used the same \texttt{ibm\_fez} processor (Eagle~r3
revision, March~7 and March~11, 2026); the $+0.031$ shift in
$F_{K^*}$ is within typical inter-session drift.
Structured selection is stable across seeds ($\sigma = 0.016$),
while the uniform random arm is wildly volatile ($\sigma = 0.325$,
coefficient of variation $> 1$): a single unlucky draw can miss
critical low-weight operators entirely.
Seed~200 is anomalous in both random arms
($F_\mathrm{UR} = 0.630$, $F_\mathrm{AR} = 0.886$);
the structured set still wins ($F_{K^*} = 0.917$,
$\Delta(S{-}AR) = +0.031$).

\begin{remark}[Hardware-replication prediction across calibration epochs]
Theorem~\ref{thm:spectral-char} characterizes the weight-class
coverage of $K^*$.  Under the physical interpretation that
$F(K^*)$ reconstruction is governed by that coverage rather than by
which specific noise directions a given calibration stresses, it
follows that $F(K^*)$ is expected to be approximately
calibration-invariant across hardware epochs -- a consequence
empirically consistent with the inter-session agreement reported in
Sec.~\ref{sec:results}.  The random-Pauli volatility quantified in
this section ($\sigma = 0.325$, coefficient of variation $> 1$)
supplies the $F(\mathrm{rand})$ side of the comparison.

Physically, calibration asymmetry -- coherent errors or readout
crosstalk that degrade certain operators more than others, the
typical regime for calibrated devices -- is expected to drive the
run-to-run spread of $F(\mathrm{rand})$ and the decrement of its
mean in the same direction: different random draws sample the
asymmetric noise differently, widening
$\mathrm{Var}[F(\mathrm{rand})]$, while the worst-affected operators
lower the average reconstruction fidelity, bringing
$\mathrm{Mean}[F(\mathrm{rand})]$ down.  Combining the approximate
constancy of $F(K^*)$ with this coupled shift in $F(\mathrm{rand})$
yields a testable prediction for future hardware replications:
epochs with elevated random-arm volatility will exhibit
correspondingly elevated mean $\Delta F$.  This suggests the
reported $\Delta F$ value carries a calibration-dependent dynamic
range, and the natural invariants to monitor in a replication are
the stability of $F(K^*)$ and the sign of $\Delta F$ -- the exact
magnitude of $\Delta F$ depends on the calibration's asymmetry
profile.
\end{remark}

\subsection{Dimensional dependence}
\label{sec:dim-dep}

A clear dimensional progression emerges from the three-arm comparison
(Fig.~\ref{fig:dim-dep}).
In detail, $\Delta(S{-}AR)$ measures how much
the specific Krawtchouk operator selection matters beyond weight
allocation: it decreases monotonically from $0.573$ ($n{=}2$)
to $0.042$ ($n{=}5$).

\begin{figure}[tb]
\centering
\includegraphics[width=\columnwidth]{fig5_dimensional_dependence.pdf}
\caption{Intra-weight residual $\Delta(S{-}AR)$ versus number of
qubits under MLE reconstruction.  The residual drops by $13.8\times$
from $n = 2$ to $n = 5$; at $n = 4$ (dashed line),
the Krawtchouk allocation condition ($2n = 2^{n-1}$) holds and weight allocation alone is
nearly sufficient.  Simulation: depol $= 0.03$, readout $= 0.02$,
$1{,}000$~shots, $20$~seeds.}
\label{fig:dim-dep}
\end{figure}

At $n = 4$, the weight-level description becomes \emph{nearly
sufficient}: intra-weight ordering contributes only a small
residual ($\Delta = 0.100$).
This smooth decrease is MLE-specific; under linear inversion the
progression is non-monotonic
(Supplemental Material~\cite{supplemental}).
At $n = 5$ ($M = 333$, $32.6$\% coverage), $K^*$ retains
a $+0.276$ overall advantage over uniform random but the
weight-allocated random gap narrows to $+0.004$ overall
(averaged across $13$~states; $+0.042$ for the W~state alone,
Table~\ref{tab:weight-alloc}), confirming
that the allocation fraction approaches unity as $n$ grows past
the sweet spot, consistent with the asymptotic convergence in
Theorem~\ref{thm:asymptotic}
(Supplemental Material~\cite{supplemental}).
A same-session dimensional sweep on \texttt{ibm\_fez} shows that
$K^*$ selection gives large advantages at both $n = 2$ ($\Delta F =
+0.461$) and $n = 4$ ($\Delta F = +0.435$)
(Table~\ref{tab:results}).  At $n = 4$, the rerun ($F = 0.876$)
reproduces the four-run mean ($F = 0.872 \pm 0.021$) within calibration drift.

Raw gaps $\Delta F$ (structured vs.\ uniform random) are
similar at $n = 2$ and $n = 4$; the Krawtchouk allocation prediction
concerns the \emph{allocation-specific} residual $\Delta(S{-}AR)$,
which decreases by $13.8\times$ from $n{=}2$ to $n{=}5$ under MLE
($28.3\times$ under projected least squares), not the total advantage.

\subsection{GHZ states: reconstructor dependence and subspace estimation}
\label{sec:ghz}

The preceding analysis assumes states with nontrivial local
structure; GHZ states concentrate information entirely in global
correlators, creating a qualitatively different regime.
The full-state MLE deficit ($F \approx 0.50$) has a structural
origin in the Krawtchouk eigenvalue allocation prescribed by
Theorem~\ref{thm:spectral-char}(iv)
(Fig.~\ref{fig:ghz-mismatch}).

A 4-qubit GHZ state
$|GHZ\rangle = (|0000\rangle + |1111\rangle)/\sqrt{2}$ has only
$15$ of $255$ Pauli operators with nonzero ideal expectation value:
$6$~weight-$2$ ZZ correlators ($\langle Z_i Z_j\rangle = +1$),
$1$~weight-$4$ all-$Z$ correlator ($\langle Z^{\otimes 4}\rangle = +1$),
and $8$~weight-$4$ phase correlators (e.g.,
$\langle X^{\otimes 4}\rangle = +1$,
$\langle Y^{\otimes 4}\rangle = +1$).
Weight-$4$ operators (ZZZZ plus $8$ phase correlators) carry
$60$\% of the total GHZ Pauli information
(measured by $\sum_i \langle P_i\rangle^2$ over nontrivial Paulis:
$9/15$), yet it receives only $11.7$\% of the Krawtchouk eigenvalue
mass and is allocated just $M_4 = 16$ of $81$ available operators.
Of the $8$~GHZ phase correlators ($XXXX$, $XXYY$, $XYXY$, $XYYX$,
$YXXY$, $YXYX$, $YYXX$, $YYYY$), the $K^*$~set includes $4$
($XXXX$, $XYXY$, $YXYX$, $YYYY$); all $6$~weight-$2$ ZZ correlators
are present.  The GHZ MLE deficit is therefore driven by the weight-$4$
budget ($16/81$), not by which specific weight-$4$ operators are
selected.

Raw hardware data shows that \emph{phase coherence
is retained}: $\langle X^{\otimes 4}\rangle = +0.864$
($\sim\!85$\% of ideal).
Rather, $F \approx 0.50$ arises because $126$ of the $137$
$K^*$~operators have near-zero GHZ expectation value, and the MLE
settles on a near-maximally-mixed state.
Simulation cross-checks confirm this: varying $\zeta \in [0, 0.2]$
gives $F \in [0.49, 0.50]$.

Precisely the eigenvalue allocation that concentrates
resources in low-weight sectors (boosting the W~state by
$\Delta F = +0.33 \pm 0.07$) necessarily starves the weight-$4$
sector where GHZ information lives.
GHZ is therefore a predicted consequence of the
eigenvalue-mass distribution in
Theorem~\ref{thm:spectral-char}(iii)--(iv): concentrated
low-weight allocation predicts degraded ratio-form reconstruction
for any state whose information concentrates in weight classes
that $K^*$ undersamples.
A full Hradil update recovers $F = 0.887$ on the same data
(Supplemental Material~\cite{supplemental}, Sec.~VII\,E),
confirming that the mismatch is solver-resolvable rather than
information-theoretic.

\begin{figure}[tb]
\centering
\includegraphics[width=\columnwidth]{fig7_ghz_mismatch.pdf}
\caption{GHZ eigenvalue-mass mismatch and subspace recovery.
Top: the GHZ state concentrates $60$\% of its Pauli information
in weight-$4$ operators, but $K^*$ allocates only $11.7$\% of
eigenvalue mass there.  This mismatch is a prediction of the
Krawtchouk framework, not a hardware artifact.
Bottom: four estimators applied to the same $137$ $K^*$~circuits (no additional data).
Full-state MLE gives $F = 0.496$ (mismatch-limited); subspace
estimation in $\{|0000\rangle, |1111\rangle\}$ recovers
$F = 0.926$; direct fidelity estimation (DFE) provides a model-free lower bound $F \geq 0.618$;
the weight-class fingerprint auto-switch (Sec.~\ref{sec:ghz})
recovers $F = 0.887$ at zero additional circuit cost.}
\label{fig:ghz-mismatch}
\end{figure}

\paragraph{Resolution: subspace estimation.}
As established above, $K^*$ covers $11$ of the $15$ operators
with nonzero GHZ coefficients.
Projecting into the GHZ-relevant subspace
$\mathrm{span}\{|0000\rangle, |1111\rangle\}$ reduces the
reconstruction to $3$~real parameters ($a$, $b = 1 - a$,
$|c|\,e^{i\phi}$).  Least-squares fitting the $137$ $K^*$
expectations in this subspace gives
$a = 0.529$, $|c| = 0.426$ (ideal: $a = 0.5$, $|c| = 0.5$), and
\begin{equation}
  F_{\mathrm{sub}} = 0.926\,.
  \label{eq:ghz-sub}
\end{equation}
The subspace estimator succeeds because the $K^*$ set includes all
weight-$0$ and weight-$4$ operators that span the GHZ subspace, a
consequence of the full weight-class coverage guaranteed by
Theorem~\ref{thm:spectral-char}(ii).
Direct fidelity estimation from the $11$~informative operators gives
a model-free lower bound $F_{\mathrm{DFE}} \geq 0.618$ (the missing
$5$~operators, if at the same hardware quality, would raise this to
${\sim}\,0.88$).  The phase correlators independently certify genuine
multipartite entanglement:
$|\langle X^{\otimes 4}\rangle| = 0.864 > 0.50$ (biseparable
bound~\cite{TothGuhne2005}).

$K^*$ thus provides data compatible with multiple
estimators (full-state MLE, subspace estimation, DFE, and
entanglement witnesses).

\paragraph{Automatic reconstructor switching via weight-class fingerprint.}
The GHZ failure mode is self-diagnosing: no extra
circuits required.  Define a weight-class energy fingerprint
$\mathcal{E}_w = \sum_{P:\,\mathrm{wt}(P) = w} \langle P\rangle^2$
from the $137$ $K^*$ expectations.  GHZ-type states have a telling \emph{even-parity signature}:
$\mathcal{E}_1$ and $\mathcal{E}_3$ are nearly zero (odd-weight
Pauli expectations vanish by $Z_2^{\otimes n}$ symmetry), while
$\mathcal{E}_2$ and $\mathcal{E}_4$ dominate.  The parity ratio
$r = \mathcal{E}_{\mathrm{odd}} / \mathcal{E}_{\mathrm{even}}$ separates the
state classes cleanly: $r = 0.010$ for the hardware GHZ,
versus $r = 0.99$ (product) and $r = 0.84$ (W~state).
We trigger subspace estimation whenever $r < 0.05$ and
$\mathcal{E}_4 / \mathcal{E}_{\mathrm{total}} > 0.25$.

On the same $137$ circuits, this rule correctly flags the GHZ
state and recovers $F_{\mathrm{sub}} = 0.887$
(up from $F_{\mathrm{MLE}} = 0.496$).  It leaves the W~state
and product state on the MLE path, where they already do well.
No additional circuits are required.
\emph{Caveat:} the even-parity trigger is specific to GHZ-type
states; other even-parity states (e.g.,
$(|0011\rangle{+}|1100\rangle)/\sqrt{2}$) would also trigger
the classifier but require a different target subspace.
This rule should therefore be understood as a
GHZ-specific heuristic, not a general blind classifier.
Within this scope, the information-budget mismatch that causes
the GHZ deficit also provides the signal for correcting it.

\subsection{Compositional scaling to 8 qubits}
\label{sec:8qubit}

Having characterized $K^*$'s performance on individual 4-qubit
subsystems---including both its strengths and its predicted GHZ
weakness---we now test whether the compositional architecture
(Sec.~\ref{sec:compositional}) extends the advantage to larger
processors.
We ran an 8-qubit chain on Rigetti Ankaa-3 with a three-arm
design: 59~$K^*$ circuits preparing $|W_8\rangle$, 59~$K^*$
circuits preparing $|{+}\rangle^{\otimes 8}$, and 59~random-Pauli
circuits preparing $|W_8\rangle$ (177~circuits total,
1{,}000~shots each, interleaved).
All three arms use global $n$-qubit tensor-product circuits;
per-patch expectation values are extracted identically via
partial-trace marginalisation, so the comparison isolates
the operator-selection strategy (structured vs.\ random)
with the compositional infrastructure held constant.
Five overlapping 4-qubit patches (A--E) were reconstructed
independently using the same $137$-operator $K^*$ set per patch.

\paragraph{Product-state validation.}
All five patches achieve $F > 0.989$ (mean $0.992 \pm 0.002$)
on $|{+}\rangle^{\otimes 8}$, confirming the compositional
infrastructure.  Boundary trace distances
($D_{\mathrm{tr}} = \tfrac{1}{2}\|\rho_A - \rho_B\|_1 = 0.03$--$0.05$) are small, consistent with single-qubit
noise.

\paragraph{Entangled-state scaling.}
$W_8$ results are more interesting.
Per-patch $K^*$ fidelities range from $0.58$ to $0.76$
(mean $0.668 \pm 0.07$), benchmarked against the true 4-qubit
marginal $\rho_\mathrm{patch} = \tfrac{1}{2}|W_4\rangle\langle
W_4| + \tfrac{1}{2}|0000\rangle\langle 0000|$.
Random Pauli selection, by contrast, gives $F = 0.09$--$0.26$
(mean $0.134 \pm 0.07$), barely above the maximally-mixed baseline
$1/d = 0.0625$.
All five bootstrap CIs exclude zero, and the per-patch gap
$\Delta F = +0.53 \pm 0.06$ is larger than what we saw at 4~qubits
($+0.33$); random sampling simply becomes exponentially dilute
at this size (Fig.~\ref{fig:scaling}).
This trend is the behavior predicted by
Theorem~\ref{thm:asymptotic}: at $n = 8$ each weight class
contains $\binom{8}{w} \cdot 3^w$ operators, and random selection
at $M = 137$ covers a vanishing fraction of each class.
$K^*$'s Krawtchouk-prescribed allocation guarantees coverage of
every weight class up to $K^* = 5$, independent of the total
operator pool size; by Theorem~\ref{thm:basin}, this weight
saturation is sufficient for basin separation on each patch.
One cost of the patch decomposition is visible in the boundary
trace distances: $D_\mathrm{tr} = 0.16$--$0.21$ for $W_8$,
roughly $3$--$6\times$ larger than the product-state boundaries
(Fig.~\ref{fig:8qubit}b).
This is expected: each patch recovers its own marginal
faithfully, but cross-patch correlations require global measurements
that the compositional architecture does not include.

\paragraph{Robustness under hardware non-uniformity.}
Per-patch fidelity spread ($0.58$--$0.76$) is not random;
it tracks position-dependent gate errors.
Patches~A and~B contain a qubit with elevated CX error
(identifiable from the product-state arm) and give
$F \approx 0.58$--$0.60$.
Patches~C--E avoid that qubit and give $F \approx 0.69$--$0.76$.
Product-state results are the key diagnostic here:
$F > 0.989$ on every patch (including the noisy ones) means the
measurement infrastructure itself is working fine, and the fidelity
loss is entirely in the $W_8$ entangling gates.
Even on the worst patches, $\Delta F > +0.46$.
Meanwhile, random fidelities cluster near $1/d = 0.0625$ regardless
of patch position. The problem there is exponential dilution at
$n = 8$, not gate quality.

\begin{figure}[tb]
\centering
\includegraphics[width=\columnwidth]{fig8_8qubit_results.pdf}
\caption{Eight-qubit compositional results on Rigetti Ankaa-3
(177~circuits, 1{,}000~shots each).
(a)~Per-patch reconstruction fidelity for the $W_8$ state:
$K^*$ (blue) versus random Pauli (red).  Random fidelities
cluster near the maximally-mixed baseline $1/d$ (dashed);
$K^*$ achieves $\Delta F > +0.46$ on every patch.  Error bars:
95\% bootstrap CIs ($10^3$~resamples).
(b)~Boundary trace distances at the four patch overlaps.
Product-state boundaries ($D_\mathrm{tr} \lesssim 0.05$) confirm
patch consistency; $W_8$ boundaries
($D_\mathrm{tr} \approx 0.18$) reflect the locality cost of
independent per-patch reconstruction for a globally entangled
state.}
\label{fig:8qubit}
\end{figure}

\begin{figure*}[tb]
\centering
\includegraphics[width=\textwidth]{fig9_scaling.pdf}
\caption{$K^*$ advantage $\Delta F$ versus qubit count across two
superconducting platforms.  The advantage grows from $+0.05$ at
$n = 2$ (overcomplete regime) through the $n = 4$ sweet spot
($+0.25$--$+0.33$) to $+0.53$ at $n = 8$ (compositional, five
patches on Rigetti Ankaa-3), as random Pauli selection becomes
exponentially dilute in $4^n$.  Error bar at $n = 4$ (IBM): sample
standard deviation over four runs; at $n = 8$: standard deviation
across five patches.  The $n = 3$ point
($\Delta F = +0.009$, Rigetti) confirms the overcomplete-regime
prediction where selection advantage vanishes.}
\label{fig:scaling}
\end{figure*}

\subsection{Multi-vendor replication (Rigetti Ankaa-3)}
\label{sec:rigetti}

Without modification, the same $137$ $K^*$ operators were run on a
Rigetti Ankaa-3 processor (84~qubits, square-lattice topology,
iSWAP native gate) through the Open Quantum cloud platform.
This is a substantially different architecture from IBM:
the 4-qubit W-state preparation decomposes into $u_3{+}\text{CX}$
sequences at roughly twice the two-qubit gate depth of the IBM
circuit.
The experiment uses tensor-product basis
grouping~\cite{Hamamura2020,Jena2019}: $29$~$K^*$ bases and
$50$~random bases ($79$~circuits total) interleaved in a single
calibration session.
Result: $K^*$ gives $F = 0.815$, random gives $F = 0.568$,
a gap of $\Delta F = +0.248$, consistent with the IBM
result ($+0.33 \pm 0.07$) on a completely different platform
(Table~\ref{tab:results}).
The cross-platform replication is expected: the basin separation
guarantee (Theorem~\ref{thm:basin}) depends only on the weight
saturation of the operator set, not on the hardware topology or
native gate set.
Bootstrap resampling ($10^3$~draws from raw bitstring counts)
puts $\sigma_{K^*} \approx 0.015$ and
$\sigma_\mathrm{rand} \approx 0.010$; every resample gives
$\Delta F > 0$.
Absolute fidelity is lower than IBM ($0.815$ vs.\ $0.872$),
which is expected given the deeper gate decomposition on the
Ankaa-3 square lattice.

$K^*$ compresses $4.7\times$ ($137 \to 29$~bases) versus
$2.7\times$ for the random set ($137 \to 50$), a $1.7\times$
circuit-efficiency advantage from the structured weight distribution
(Fig.~\ref{fig:efficiency}; Appendix~\ref{sec:compression-advantage}).
Under the compositional architecture
(Sec.~\ref{sec:compositional}), this gap persists at every
processor size: $K^*$ requires $\le 116$~basis-grouped circuits for any
linear chain, compared to ${\sim}\,200$ for the random selection at
matched operator count (Fig.~\ref{fig:efficiency}, right).

\begin{figure}[tb]
\centering
\includegraphics[width=\columnwidth]{fig6_circuit_efficiency.pdf}
\caption{Circuit efficiency of $K^*$ selection.
Top: tensor-product bases required at $n = 4$ ($M = 137$ operators
each arm).  $K^*$ compresses into $29$~bases ($4.7\times$) versus
$50$ for random ($2.7\times$), a $1.7\times$ advantage.  Full
tomography ($81$~bases, all $255$ operators) is shown for reference.
Bottom: total circuit count versus processor size under the
compositional architecture.  $K^*$ and random both plateau (constant
patch structure), while full tomography grows as $3^n$.}
\label{fig:efficiency}
\end{figure}

At $n = 3$ the story is different: structured and random give
nearly identical fidelity ($F = 0.840$ vs.\ $0.831$,
$\Delta F = +0.009$).
This is exactly what Theorem~\ref{thm:spectral-char}(iii) predicts:
with $M/4^n = 27/64 \approx 42\%$ of the full operator set, the
regime is close to overcomplete and there is little room for
selection to matter.
Contrast with $n = 4$ ($\Delta F = +0.25$ on Rigetti,
$+0.33$ on IBM) is sharp: at $M/4^n = 137/256 = 54\%$ the
budget is tight enough for weight-class allocation to bite.
Going the other direction, the 8-qubit compositional run on the
same Ankaa-3 processor gives $\Delta F = +0.53 \pm 0.06$ across
five patches (Sec.~\ref{sec:8qubit}; Fig.~\ref{fig:scaling}). The
advantage keeps growing as random selection becomes exponentially
dilute.

\subsection{Operator-selection comparison (simulation)}
\label{sec:shadows}

Having established the structured advantage on hardware across two
platforms and system sizes, we now compare operator-selection
strategies under controlled simulation conditions.
Table~\ref{tab:shadows} compares $K^*$-structured selection against
derandomized classical shadows~\cite{Huang2021} and uniform-random
shadows at matched budget (137~settings, $1{,}000$~shots,
depol $= 0.03$, readout $= 0.02$) under a common MLE reconstructor.

\begin{table}[tb]
\centering
\caption{Simulated comparison of measurement \emph{selection strategies}
under a common MLE reconstructor.
Settings: 137~bases $\times$ 1{,}000~shots, uniform depolarizing noise
(depol $= 0.03$, readout $= 0.02$); no hardware data.}
\label{tab:shadows}
\begin{tabular}{lcccr}
\toprule
State & $F(K^*)$ & $F(\mathrm{DR})$ & $F(\mathrm{RS})$
      & $\Delta(K^*{-}\mathrm{DR})$ \\
\midrule
$|{+}\rangle^{\otimes 4}$ & 0.975 & 0.994 & 0.994 & $-0.019$ \\
Bell                       & 0.670 & 0.550 & 0.536 & $+0.120$ \\
$W$                        & 0.829 & 0.497 & 0.626 & $+0.332$ \\
GHZ                        & 0.756 & 0.506 & 0.507 & $+0.250$ \\
\bottomrule
\end{tabular}
\end{table}

$K^*$ selection beats both shadow-basis selections for all
entangled states.
Under the native shadow channel-inverse estimator, the full protocol
($81$ bases, $255$ Paulis) reaches $F = 0.97$ at $1.86\times$ the
measurement budget; with $40$ random bases, $F = 0.77$ (W) and
$0.69$ (GHZ).  $K^*{+}\text{MLE}$ ($137$~Paulis) exceeds the
$40$-base shadow for W ($0.95$ vs.\ $0.77$); for GHZ,
full-state MLE gives $0.61$ vs.\ shadow's $0.69$, but subspace
estimation from the same $K^*$ data recovers $F = 0.926$
(Sec.~\ref{sec:ghz}).

\paragraph{Broader comparison.}
Table~\ref{tab:sota} and Fig.~\ref{fig:sota} extend the
within-pipeline comparison (common MLE reconstructor, matched
operator budget) to greedy optimal
experimental design (A/D/E-optimal), adaptive multi-round
selection, and weight-allocated random.
All non-adaptive methods select $137$~operators
($1{,}000$~shots, depol~$= 0.03$, readout~$= 0.02$,
common robust MLE reconstructor), except the
derandomized-shadow baseline which selects $137$~bases
and reconstructs from all covered operators.
Mean fidelity is computed across $23$~test states (product,
W, GHZ, $10$~Haar-random pure, $10$~random mixed of rank~$4$)
at $20$~trials each; adaptive methods split the $137$-operator
budget into multiple rounds ($3$ or $5$) with
information-gain-based selection between rounds.

\begin{table}[tb]
\centering
\caption{Mean reconstruction fidelity across $23$~states ($20$~trials).
All non-adaptive methods select $137$~Pauli operators,
$1{,}000$~shots, common MLE reconstructor,
depol~$= 0.03$, readout~$= 0.02$.
Note: the derandomized-shadow (DR-shadow) entry selects $137$ tensor-product
\emph{bases} via greedy set-cover, then reconstructs from the union
of all covered operators ($> 137$ individual Paulis); this gives
DR-shadow a budget advantage in operator count, making the comparison
conservative for $K^*$.
Our DR-shadow implementation omits per-observable weighting
(Supplemental Material~\cite{supplemental}); a fully weighted
implementation may perform better.
Adaptive methods (below second rule) are included as
out-of-regime controls; they split the fixed budget into
multiple rounds, a regime where they are not designed to
operate (Sec.~\ref{sec:shadows}).}
\label{tab:sota}
\resizebox{\columnwidth}{!}{%
\begin{tabular}{lcccccr}
\toprule
Strategy & Product & W & GHZ & Haar & Mixed & Overall \\
\midrule
$K^*$          & 0.991 & \textbf{0.932} & \textbf{0.502} & 0.397 & \textbf{0.513} & \textbf{0.501} \\
D-optimal      & 0.987 & 0.811 & 0.225 & \textbf{0.398} & 0.502 & 0.479 \\
E-optimal      & \textbf{0.992} & 0.833 & 0.225 & 0.317 & 0.466 & 0.429 \\
A-optimal      & \textbf{0.993} & 0.833 & 0.229 & 0.316 & 0.462 & 0.428 \\
WA-random      & 0.888 & 0.414 & 0.352 & 0.242 & 0.391 & 0.347 \\
DR-shadow      & 0.952 & 0.567 & 0.114 & 0.066 & 0.227 & 0.198 \\
\midrule
Adaptive (3R)  & 0.120 & 0.096 & 0.139 & 0.062 & 0.224 & 0.140 \\
Adaptive (5R)  & 0.034 & 0.092 & 0.075 & 0.062 & 0.224 & 0.133 \\
Uniform random & 0.711 & 0.291 & 0.215 & 0.063 & 0.227 & 0.179 \\
\bottomrule
\end{tabular}}% end resizebox

\end{table}

\begin{figure}[tb]
\centering
\includegraphics[width=\columnwidth]{fig4_sota_comparison.pdf}
\caption{Per-state-class fidelity for nine measurement selection strategies
at matched budget ($137$~operators, $1{,}000$~shots, depol $= 0.03$,
readout $= 0.02$; common MLE reconstructor).  Black diamonds mark overall
mean across $23$~states.  $K^*$ ranks first overall;
the advantage is concentrated on entangled states (W, GHZ, mixed).
Numerical values in Table~\ref{tab:sota}.}
\label{fig:sota}
\end{figure}

Under the common MLE reconstructor, $K^*$ ranks first overall,
consistent with the spectral construction's full weight-class
coverage (Theorem~\ref{thm:spectral-char}(ii)),
ahead of greedy D-optimal by $+0.022$
(Wilcoxon signed-rank $p < 10^{-4}$, paired across $20$~trials;
Supplemental Material~\cite{supplemental}, Sec.~VI\,E).
\emph{Note on optimal-experimental-design (OED) divergence:}
Appendix~\ref{sec:fisher} shows that all $M$-element Pauli subsets
are equivalent under state-independent Fisher criteria.
Greedy D/E/A-optimal algorithms in Table~\ref{tab:sota}
nonetheless produce different operator sets because they are
sequential heuristics that break ties differently at each step:
e.g., D-optimal maximizes log-determinant while E-optimal
maximizes the minimum eigenvalue, leading to different greedy
paths through the same degenerate landscape.
Fidelity differences arise primarily from which
weight classes each greedy path happens to cover,
not from any intrinsic advantage of one OED criterion over
another (Supplemental Material~\cite{supplemental}).
$K^*$ also has a computational advantage: the operator set is a
precomputed lookup table ($O(1)$ at measurement time), whereas
greedy D-optimal requires $O(M \cdot 4^n)$ sequential selection
steps (${\sim}\,60$\,s at $n = 4$ on a single CPU core; scaling intractably beyond $n \approx 6$).
Separation is largest on the W~state ($+0.12$ over D-optimal,
$+0.10$ over E-optimal), the GHZ state ($+0.28$ over D-optimal),
and random mixed states ($+0.05$); $K^*$'s weight-class balance
preserves both local correlators (W) and global phase correlations
(GHZ) that greedy OED does not prioritize.
For product states, all non-adaptive strategies perform comparably;
the advantage concentrates on entangled states where weight-class
completeness matters.

Adaptive methods perform below uniform random.
Splitting $137$~operators into $3$--$5$ rounds
leaves each round severely underdetermined
($27$--$45$ measurements for $d^2 - 1 = 255$~parameters),
producing unreliable intermediate reconstructions that
misguide subsequent selection.
Using $K^*$ as a high-quality first round does not help:
a hybrid protocol ($K^*$ round-$1$ + adaptive refinement
for remaining budget) matches pure $K^*$ to within
$\pm 0.001$ at all tested budget levels, including
informationally complete ($M = 255 = 4^n - 1$;
Supplemental Material~\cite{supplemental}).
Adaptive overhead from round-splitting and intermediate
reconstruction provides no marginal benefit over
single-batch measurement when the initial selection is
already well structured.
This failure is specific to the underdetermined regime
($M/4^n < 1$) and our round-splitting implementation;
production adaptive tomography systems~\cite{Huszar2012}
typically operate in the overcomplete regime with
interactive mid-experiment feedback, conditions not
available on cloud quantum-processing-unit (QPU) platforms.
Our comparison therefore isolates the regime where $K^*$
is designed to operate, not adaptive methods in general.

\subsection{Application: Floquet discrete time crystal}
\label{sec:dtc}

The preceding benchmarks compare operator sets under a common
reconstructor.
To illustrate $K^*$'s ability to access off-diagonal
observables that are absent from Z-basis measurement records
(which contain only diagonal expectation values),
we apply the measurement set to a
Floquet discrete time crystal (DTC) on IBM
\texttt{ibm\_kingston} ($n = 4$ qubits, $1{,}000$ shots per
circuit, $137$ $K^*$ operators).

Our Floquet unitary $U_F = U_{\mathrm{kick}}(\phi)\,
U_{\mathrm{XXZ}}(J, \varepsilon)$ acts on an initial N\'eel
state $|0101\rangle$ for $6$ Floquet cycles, with nearest-neighbor coupling $J = 1$ (energy unit), kick angle
$\phi = \pi/2$ (DTC phase), anisotropy $\varepsilon = 0.1$,
and weak bond disorder (amplitude $0.5$).  The transpiled circuit
has depth~$329$ on the Heron-class backend.

\emph{Z-only diagnostics.}
Stroboscopic magnetization shows period-$2$ oscillation
with period-$2$ magnetization contrast $\Delta_{\mathrm{DTC}} = 0.826$, confirming the DTC
signature.  However, the Z-basis order parameter
$\Delta_{\mathrm{DTC}}$ can be replicated by a classical
kicked Ising model with the same symmetry breaking;
it does not certify quantum coherence~\cite{Mi2022,Ippoliti2021}.

\emph{Off-diagonal correlators from $K^*$.}
Because $K^*$ includes weight-$1$ and weight-$2$ operators
at all Pauli types (not just~$Z$), the raw expectation
values directly yield nearest-neighbor XY correlators
($\langle X_i X_j\rangle$, $\langle Y_i Y_j\rangle$,
$\langle X_i Y_j\rangle$ for $j = i + 1$).
Comparing these $9$ correlators against noiseless simulation
of the same Floquet circuit gives Pearson correlation
$r = 0.94$ with $8/9$ sign agreement
(Table~\ref{tab:dtc});
a Fisher $z$-transform $95\%$ confidence interval gives
$r \in [0.72,\, 0.99]$ ($n_{\rm corr} = 9$), and a
permutation test yields $p < 0.002$ for the null hypothesis
of uncorrelated hardware and ideal values.
Only one sign mismatch occurs, at
$\langle X_1 Y_2\rangle$, where the ideal value
($+0.044$) is near zero and the hardware value ($-0.008$)
is within shot noise.

\begin{table}[tb]
\centering
\caption{Nearest-neighbor XY correlators for a 4-qubit
Floquet DTC ($6$ cycles, N\'eel initial state) on IBM
\texttt{ibm\_kingston}.  Ideal values from noiseless
simulation of the same circuit.  Pearson $r = 0.94$,
$8/9$ sign agreement.}
\label{tab:dtc}
\begin{tabular}{lccc}
\hline\hline
\textrm{Pair} & $\langle XX\rangle$ & $\langle YY\rangle$
& $\langle XY\rangle$ \\
\hline
$(0,1)$\,hw   & $-0.111$ & $-0.063$ & $+0.084$ \\
$(0,1)$\,ideal & $-0.123$ & $-0.123$ & $+0.067$ \\[3pt]
$(1,2)$\,hw   & $+0.078$ & $+0.094$ & $-0.008$ \\
$(1,2)$\,ideal & $+0.096$ & $+0.096$ & $+0.044$ \\[3pt]
$(2,3)$\,hw   & $-0.116$ & $-0.129$ & $+0.139$ \\
$(2,3)$\,ideal & $-0.122$ & $-0.122$ & $+0.067$ \\
\hline\hline
\end{tabular}
\end{table}

These off-diagonal correlators are consistent with quantum
coherence in the time-crystalline phase: the $\ell_1$ coherence
$\mathcal{C}_{\ell_1} = \sum_{i \neq j}|\rho_{ij}|$
($= 2.02$ from the $K^*$-reconstructed
density matrix, $\delta\mathcal{C}_{\ell_1} \approx 0.2$ from
shot noise) and partial-transpose negativity
$\mathcal{N} = (\|\rho^{T_A}\|_1 - 1)/2$
($= 0.106$, $\delta\mathcal{N} \approx 0.04$;
computed from the $K^*$-reconstructed density matrix,
uncertainties from $10^3$ bootstrap resamples)
are absent from Z-basis measurement records.
$K^*$ resolves these quantities at no
additional circuit cost beyond the $29$ tensor-product bases
already used for local tomography.
This is a direct consequence of the forward direction of the
spectral correspondence: the Krawtchouk eigenvalue masses
(Lemma~\ref{lem:spectral_q_main}) allocate operators across
all Pauli weight classes, which naturally covers the conserved % vague-verb-justified: domain-standard; eigenvalue mass allocation is a structural property of the Hamming scheme
sectors of the XXZ interaction and provides the off-diagonal
correlators that Z-basis protocols structurally lack.

%% ========================================================
\section{Discussion}
\label{sec:discussion}

The preceding experiments confirm the Krawtchouk framework's
predictions and expose its limits.
We now synthesize these findings, position the method against
existing approaches, and identify open directions.

\subsection{The bidirectional correspondence}
\label{sec:bidirectional}

The forward direction (Krawtchouk eigenvalues $\to$ measurement
allocation; Sec.~\ref{sec:selection}) and the reverse direction
(weight saturation as MLE identifiability criterion;
Theorem~\ref{thm:basin}) were introduced in Sec.~\ref{sec:intro}
and established formally in Secs.~\ref{sec:theory}--\ref{sec:experiment}.
Sections~\ref{sec:head-to-head}--\ref{sec:dtc} confirm that the
bridge operates in practice: the spectral prescription outperforms
alternatives on hardware, and the basin separation guarantee
correctly predicts both where the advantage is largest (W~states)
and where it breaks down (GHZ).
We now position this bidirectional character against analogous
cross-disciplinary bridges in the recent literature.
Recent PRX Quantum publications have established similarly productive
cross-disciplinary bridges: chi-boundedness from graph theory applied
to shadow tomography~\cite{King2025}, algebraic graph theory applied
to noise learnability~\cite{ChenFlammia2026}, and measurement frame
theory applied to shadow estimation~\cite{Innocenti2023}.
Each imports a mathematical structure to solve a quantum problem.
Our correspondence adds a reverse direction: the basin
separation theorem exports a reconstruction-theoretic rigidity
(the dichotomy of Theorem~\ref{thm:spectral-char}(iii))
back to the Hamming-scheme setting, and extends to all $q$-ary
alphabets (Lemma~\ref{lem:spectral_q_main}).

\subsection{Mechanism and scope of the advantage}

Why does structured selection help at all?  Standard design criteria
offer no answer.  Any $M$-element Pauli subset has identical
state-independent Fisher information, frame potential, condition
number, and mutual coherence (Appendix~\ref{sec:fisher}), so
any difference in reconstruction quality must arise through the
state-dependent, nonlinear part of the estimator.

Theorem~\ref{thm:basin} pins it down.  Weight-class completeness
(i.e., $c_w > 0$ for every Hamming weight~$w$;
Lemma~\ref{lem:completeness})
pushes the Hessian's zero eigenvalues into directions where the
target state carries no information.  The positivity constraint
then resolves only those uninformative directions, with a bounded
excursion $\varepsilon_{\mathrm{pos}} = O(p/d)$.  A direct test
confirms the nonlinear coupling: under plain linear inversion,
uniform random \emph{outperforms} $K^*$ ($\bar{F} = 0.574$ vs.\
$0.454$; Supplemental Material~\cite{supplemental}, Sec.~VI\,E;
Appendix~\ref{sec:mle-mechanism} gives the full ablation study).
This is consistent with the finding that randomized measurements are
near-optimal for linear estimation of pure-state
properties~\cite{Huang2020}; any advantage from measurement
selection has to couple through a nonlinear reconstructor, or it
is not there.
In the underdetermined regime ($M < 4^n$), physical reconstruction
(MLE, Bayesian) is mandatory (linear inversion yields unphysical
outputs), so the relevant question is not whether the reconstructor
is nonlinear but which operator selection gives the nonlinear
reconstructor the best raw material.
$K^*$ answers this: it distributes measurement information across
weight classes in the proportions prescribed by the Krawtchouk
eigenvalue masses, which constrain the MLE feasible set most
tightly per circuit.

This non-variational origin is load-bearing: because $K^*$ comes
from algebraic structure rather than from optimizing a
reconstruction objective, the measurement set does not depend on
which reconstructor will be used downstream.  The same $137$
operators feed MLE, linear inversion, subspace estimation, DFE,
and entanglement witnesses, with no redesign step in between.

Where is the advantage largest?  The eigenvalue-mass analysis
(Sec.~\ref{sec:ghz}) predicts this \emph{a~priori}: the advantage
is concentrated on states with nontrivial local structure
($\Delta F = +0.33$ on W), vanishes for GHZ where information
sits in undersampled global correlators, and is small for product
states that both arms already reconstruct well.
The multi-estimator compatibility demonstrated in
Sec.~\ref{sec:ghz} means the measurement campaign is
reusable across reconstruction tasks without redesign.
In return, a modest per-state optimality cost buys
target-agnostic reusability: only $0.29$\% of random draws
match $K^*$'s performance across all state classes simultaneously
(Supplemental Material~\cite{supplemental}, Sec.~V).

\subsection{Comparison to existing methods}
\label{sec:comparison}

Measurement design methods generally fall into two camps.
\emph{Variational} approaches (OED, derandomized shadows)
optimize a cost function over operator subsets;
\emph{sampling-based} approaches (classical shadows, compressed
sensing) draw operators from a distribution.
Both involve design choices (an objective, a target observable set,
a sampling distribution) that $K^*$ does not require.
By contrast, $K^*$ is the algebraically determined output of
Krawtchouk spectral analysis on~$H(n,2)$ at the critical cutoff,
requiring none of these design choices.
This makes the approaches complementary rather than competing
(Supplemental Material~\cite{supplemental}).

\paragraph{Cloud-platform regime.}
$K^*$'s natural habitat is the underdetermined regime
($M < 4^n$), which is the default on cloud platforms where
circuits are non-interactive and billed per shot.
In this regime $K^*$ ranks first among all tested alternatives
under the common MLE pipeline (Table~\ref{tab:sota}), with
negligible precomputation (a fixed lookup rather than the
iterative design cost of greedy OED) and a $1.7\times$
circuit-compression advantage
(Appendix~\ref{sec:compression-advantage}).
That said, $K^*$ is not a universal replacement.
Classical shadows~\cite{Huang2020} predict many linear
observables at $O(\log M / \varepsilon^2)$ cost without full
reconstruction, the right tool when individual expectation values
are the target, especially at large~$n$ where full
local density matrices are not needed.
Adaptive Bayesian methods win when
interactive overcomplete budgets are
available~\cite{Huszar2012,HuangZhu2025,Lange2023};
DFE~\cite{Flammia2011} is the right tool when the target state is
known; compressed sensing~\cite{Gross2010} exploits low-rank priors
that $K^*$ ignores.
At scale, the distinction sharpens: shadows and $K^*$ answer
different questions (linear functionals versus full local
states) and are complementary rather than competing on
next-generation processors.  Comparing $K^*{+}\text{MLE}$ against
shadows${+}$median-of-means conflates the reconstructor choice
with the measurement choice; the fair comparison holds the
reconstructor fixed (Table~\ref{tab:sota}) and varies only the
operator selection.

\paragraph{Scope within the characterization landscape.}
$K^*$ targets one specific task: \emph{measurement design for
underdetermined Pauli tomography}, within a larger characterization
ecosystem.
Its compositional architecture extends to arbitrarily large
processors at $O(1)$ circuit cost via overlapping patches
(Sec.~\ref{sec:8qubit}), and the advantage is predicted to
\emph{grow} with system size: random Pauli selection dilutes as
$4^{-n}$ per patch while $K^*$'s per-patch allocation is
size-independent.  Hardware validation currently reaches $n = 8$
(five patches on Rigetti Ankaa-3,
$\Delta F = +0.53$); larger-scale demonstration awaits access to
${\ge}\,20$-qubit experiments.
Full local density matrices are also relevant to quantum advantage
benchmarking: cross-entropy (XEB) scores require comparison against
an ideal output distribution, and verifying that a prepared
subsystem matches its target requires tomographic reconstruction
at a cost that $K^*$ keeps constant per patch.
Generative state-learning models~\cite{Akhound-Sadegh2024} solve
a different problem entirely (sampling from rather than
reconstructing the state), though they could in principle use $K^*$
as a measurement front end.
Gate-level benchmarking~\cite{Erhard2019} is complementary:
it diagnoses individual gate errors, not full-state properties.
Patch structure also aligns naturally with quantum error % vague-verb-justified: domain-standard; QEC subsystem alignment is a structural consequence of locality
correction: each code block is a local subsystem of the size
$K^*$ targets, so the same constant-cost protocol characterizes
logical-qubit patches regardless of code distance.
State-independence is deliberate (blind certification
requires it), but it means adaptive methods will outperform $K^*$
whenever strong priors and interactive protocols are available.

\paragraph{Beyond weight counting.}
Weight allocation captures ${\sim}\,81$\% of the structured
advantage, a fact the eigenvalue-mass analysis predicts
(Sec.~\ref{sec:dim-dep}); the Krawtchouk framework provides
four additional predictions that weight counting alone cannot make
(Sec.~\ref{sec:intro}): the exact threshold $K^* = 5$, the GHZ
deficit, the dimensional peak at $n = 4$, and the
saturated-vs-subsampled weight dichotomy.
An additional ${\sim}\,19$\% from intra-weight operator selection
is confirmed by the hardware scaling: $\Delta F$ grows from $+0.33$
at $n = 4$ to $+0.53$ at $n = 8$, consistent with finer
intra-weight discrimination mattering more as the operator pool
grows.
Krawtchouk spectral analysis determines the allocation from
eigenvalue masses $\lambda_w$; any approach that reproduces the
allocation empirically (e.g., by exhaustive search over weight
splits) would need to explain \emph{why} those particular weights
work.  The spectral framework provides that explanation and extends
to other dimensions and qudit alphabets (Supplemental Material~\cite{supplemental}), determining
the full measurement prescription from a single integer~$K^*$ with
no free parameters.
Absence of optimization is a structural feature, not a
practical shortcut: because the allocation is uniquely determined
by Krawtchouk eigenvalue masses, there is no design freedom to
exploit and no objective function to choose.
This contrasts with greedy OED ($O(M \cdot 4^n)$ sequential steps)
and derandomized shadows (combinatorial optimization over basis
assignments), where the result depends on which optimality
criterion is selected.

\subsection{Predictions}
\label{sec:predictions}

Eigenvalue-mass analysis generates three testable predictions:
(1)~the W-state advantage will persist on
trapped-ion~\cite{Haffner2005,Wright2019} hardware;
(2)~the full-state MLE GHZ deficit will persist on any hardware
(the mismatch is structural), but the weight-class fingerprint
classifier will automatically detect and recover
$F > 0.85$ wherever phase coherence is preserved;
(3)~enriching weight-$4$ allocation does \emph{not} recover GHZ
fidelity (simulation gives
$F(\text{enriched}) < F(\text{default})$ under all tested noise
channels).
Rigetti replication (Sec.~\ref{sec:rigetti}) confirms the
advantage on a second superconducting platform; the trapped-ion
prediction remains open.

\subsection{Limitations and open questions}
\label{sec:limitations}

Basin separation (Theorem~\ref{thm:basin}) applies
directly to weight-$k$ supported states and, via the approximate extension, to
states with small high-weight tails.  For states with large weight
support (e.g., GHZ, where $W(\rho) = n$), the theorem does not
guarantee high fidelity, consistent with the observed GHZ deficit
($F \approx 0.50$).  A formal extension to non-local states would
require incorporating the weight-class fingerprint classifier
(Sec.~\ref{sec:ghz}) into the theoretical framework.

Intra-weight operator ordering is non-unique: the Krawtchouk
eigenvalues prescribe per-weight totals, not which specific
operators within each weight class to measure.  A three-arm test
(Sec.~\ref{sec:weight-alloc}) shows this choice has minor effect
at $n = 4$.  Predictive power of the eigenvalue masses operates
through the per-weight allocation, not individual operator
choice.
Hardware results span two superconducting vendors; trapped-ion
verification is open.
Dimensional-dependence predictions (Sec.~\ref{sec:dim-dep})
rest on IBM data alone; the Rigetti $n = 3$ overcomplete result
($\Delta F = +0.009$) is consistent but not independently
decisive.
At 8 qubits, the compositional test confirms the $K^*$ advantage for
entangled states ($\Delta F = +0.53$, Sec.~\ref{sec:8qubit}) but
reveals that boundary consistency degrades for globally entangled
states ($D_\mathrm{tr} \approx 0.18$ for $W_8$ vs.\
$0.04$ for product); entangled-state composability beyond
patch-local marginals remains an open problem for all local
tomography methods.
Because the compositional architecture uses a fixed 4-qubit patch
size, per-patch reconstruction quality is independent of processor
scale; only boundary consistency and state-preparation depth change
with~$n$.
Compositional design is not merely a convenience but a
practical necessity.
At $K = 5$, weight-$1$ and weight-$2$ saturation holds for all
$n \ge 4$ (proved; Supplemental Material~\cite{supplemental}),
but full saturation to $\lfloor n/2 \rfloor$ holds only for
$n \le 5$.
For $n \ge 6$, the monolithic cutoff $K_{\mathrm{full}}(n)$
required for full saturation grows with~$n$, and the operator
budget scales rapidly:
$137$ ($n{=}4$) $\to$ $1{,}341$ ($n{=}6$, $K_{\mathrm{full}}{=}6$)
$\to$ $12{,}369$ ($n{=}8$, $K_{\mathrm{full}}{=}7$).
By contrast, the compositional architecture uses $n - 3$
overlapping $4$-qubit patches ($2$ at $n{=}5$, $5$ at $n{=}8$;
Sec.~\ref{sec:8qubit}), with a per-patch budget of $137$
operators.  At $n = 8$ this gives $5 \times 137 = 685$
operators, an $18\times$ saving over the monolithic budget.
This quantitative gap motivates the patch decomposition as the
scaling mechanism for $K^*$.
Our matched-budget comparison (Table~\ref{tab:sota}) uses our own
implementations of each strategy rather than established
reference codebases; while the implementations follow the
published specifications faithfully, this is a limitation
relative to comparisons using independently validated software.
The comparison is simulation-only
at $n = 4$; at $n = 5$, the $K^*$-vs-random gap narrows as the
allocation fraction approaches unity
(Supplemental Material~\cite{supplemental}).  At $n = 5$, the comparison uses
$M = 333$ operators at $1{,}000$~shots ($333{,}000$ total shots
vs.\ $137{,}000$ at $n = 4$); equalizing the total shot budget
would reduce per-operator precision and could widen the gap.
The $K^*$ advantage requires nonlinear reconstruction
(Supplemental Material~\cite{supplemental}); more realistic noise models
(crosstalk, leakage) remain untested in simulation, though
hardware results incorporate these effects implicitly.
The raw-fidelity advantage is also noise-regime dependent:
in the near-noiseless limit (e.g.\ two-qubit gate error rate $\varepsilon_{2Q}$ satisfying
$\varepsilon_{2Q}\!\cdot\!\text{depth}
\lesssim 10^{-3}$, characteristic of post-fault-tolerance regimes),
the reconstruction becomes information-limited and
random or shadow-based estimators may match or exceed
$K^*$'s raw fidelity on generic states.
$K^*$'s structural advantages, certifiable bounds via the
basin-separation theorem and direct access to off-diagonal
correlators (Sec.~\ref{sec:dtc}), persist independently of
noise level; on every superconducting hardware platform shipping
in 2026, the noise regime sits well above this crossover and
$K^*$ dominates raw fidelity by $\Delta F = +0.016$ to $+0.33$
(Tables~\ref{tab:results},~\ref{tab:weight-alloc-hw}).
Calibration drift is ruled out by reversed execution order and
stability across seven runs ($\sigma = 0.024$).
Extended robustness checks (ordering sensitivity,
shot allocation, statistical power, Bayesian alternatives,
Cram\'{e}r--Rao bounds) are detailed in
Supplemental Material~\cite{supplemental}.

\paragraph{Formal verification and audited proof statements.}
The supplementary \texttt{lean4-formal/} repository contains
machine-checked Lean~4 proofs with no admitted goals
(no \texttt{sorry}), including, to our knowledge, the first
formal proof that the Krawtchouk spectral correspondence holds
if and only if $q \le 3$
(\texttt{bose\_mesner\_combinatorial\_iff}).
Two further lemmas verify weight-saturation arithmetic
(\texttt{weight12\_saturation\_K5\_all\_n}) and MLE-update
contractivity (\texttt{R\_closed\_moves\_when\_b\_ne\_y}).
Beyond the proofs themselves, three mechanized audits
(hypothesis ablation, definition mutation testing, \texttt{tex}-to-Lean
coverage; details in the Lean\,4 verification appendix below)
defend against the standard failure modes of formalization work:
every hypothesis is load-bearing and no labels are unclassified
or stale.

%% ========================================================
\section{Summary and Outlook}
\label{sec:conclusion}

The bidirectional correspondence between Hamming-scheme
algebra and quantum measurement design
(Sec.~\ref{sec:bidirectional}) rests on three theoretical pillars:
\pclaim{prose:summary:basin}{proved}{*}{thm:basin}%
basin separation (Theorem~\ref{thm:basin}), which gives weight
saturation a quantum-operational reading as an MLE identifiability
criterion;
a purity bound (Lemma~\ref{lem:purity_main}), which closes the
fidelity guarantee at $\varepsilon_{\mathrm{pos}} = O(p/d)$;
\pclaim{prose:summary:iff}{proved}{*}{lem:spectral_q_main}%
and an iff characterization
(Lemma~\ref{lem:spectral_q_main}), which shows the Krawtchouk
spectral decomposition holds exactly when $q \le 3$, while the
weight-block-diagonal structure that drives basin separation
extends to all~$q$.
Together these yield a fully determined prescription ($M = 137$
operators at $n = 4$, compiling into 29 tensor-product circuits)
whose only input is the qubit count~$n$.

Experimentally, the correspondence delivers consistent
advantages across platforms, system sizes, and state classes
(Table~\ref{tab:results}; Fig.~\ref{fig:scaling}).
Structured selection outperforms random for every entangled
state tested, with the advantage growing from
$\Delta F = +0.052$ at $n = 2$ to $+0.53 \pm 0.06$ at $n = 8$
as random Pauli sampling becomes exponentially dilute
(${\sim}\,4^{-n}$ per weight class).
The eigenvalue allocation correctly predicts where the
framework struggles: GHZ states concentrate information in
undersampled weight classes, and automatic detection plus
estimator switching recovers high fidelity from the same
circuits (Sec.~\ref{sec:ghz}).
Weight-$2$ completeness accounts for the dominant share of
the advantage (three-arm hardware test: $76 \pm 12$\%;
Appendix~\ref{sec:mle-mechanism}).
$K^*$ matches or exceeds all tested non-adaptive alternatives
at matched budget (Table~\ref{tab:sota}), and a compositional
architecture extends the protocol to arbitrary linear-chain
lengths at constant per-patch circuit cost.
Two open questions remain: whether the basin separation theorem
extends to non-local states via a weight-class fingerprint
classifier that adapts the estimator per subsystem, and whether the
reconstruction dichotomy (Theorem~\ref{thm:spectral-char}(iii))
yields tighter linear-programming constraints on covering-code
parameters in the Delsarte theory.
We publish the complete operator set, all raw
data, and the full simulation and reconstruction
code~\cite{supplemental} for independent replication.

\paragraph*{Operating regime and longevity.}
Protocol value is robust to the noise-floor trajectory
(quantitative crossover analysis in Supplemental
Material~\cite{supplemental}, Sec.~VIII).
As gate fidelities improve toward the post-fault-tolerance regime
($\varepsilon_{2Q}\!\cdot\!\text{depth} \lesssim 10^{-3}$), the
raw-fidelity advantage over random Pauli sampling fades on generic
states (Sec.~\ref{sec:limitations}).  Three structural properties,
however, depend on the algebraic design of the operator set rather
than on the noise level, and therefore remain relevant
independent of gate fidelity:
basin separation (Theorem~\ref{thm:basin}) provides worst-case
fidelity bounds conditioned only on weight saturation, complementing
the average-case guarantees of shadow estimation~\cite{Huang2020}
and making $K^*$ a natural choice for safety-critical deployments
where per-state certification is required;
direct single-shot access to off-diagonal correlators
(Sec.~\ref{sec:dtc}) enables coherence diagnostics that $Z$-basis
or shadow protocols cannot produce without additional
post-processing variance, a permanent advantage for Hamiltonian
learning, time-crystal certification, and mid-circuit measurement
experiments;
and the compositional architecture's circuit budget is constant
per patch, growing only with system topology, not with $n$,
yielding an $18\times$ operator-count saving at $n = 8$ that grows
monotonically with system size independent of per-gate fidelity.
Even after fault tolerance, code-block characterization continues
to operate on $4$--$9$ physical qubits at the surface-code patch
scale, the regime where $K^*$ was constructed.
The protocol therefore transitions from a near-term
performance tool to a long-term verification primitive embedded
in the fault-tolerant stack, rather than being superseded.

%% ========================================================
\section*{Verification and AI disclosure}

The author formulated all conjectures, chose all high-level proof
strategies, and made all substantive mathematical decisions.
Every algebraic identity, bound, and numerical prediction was
verified through an eight-tier automated pipeline spanning four
independent computation backends:
SageMath/FLINT provides exact arithmetic over GMP integers and
Pari number fields;
SymPy independently re-derives each bound in rational arithmetic;
NumPy checks numerical consistency of all tabulated values against
their authoritative data files;
and Lean~4 (Mathlib, v4.29.0-rc8) machine-verifies the formal
statements with zero \texttt{sorry} obligations
(Appendix, ``Lean\,4 formal verification status'').
Cross-tier implication chains confirm that conclusions in each tier
follow from the certified outputs of earlier tiers, and a
manuscript-claim layer independently parses every numerical value
from the \LaTeX{} source and verifies it against its data file.
No AI tools were used in any verification tier.
The verification scripts and Lean~4 proofs were developed
independently of the proof-formalization sessions.

Claude (Anthropic, Opus~4.6) was used as a proof-formalization
tool: drafting LaTeX proofs from the author's specified
conjectures and strategies, and for copyediting manuscript prose.
Claude occasionally selected specific algebraic techniques
(e.g., Schur-lemma averaging, Le~Cam's two-point method) to
execute the strategies the author specified.
A provenance record documenting the author's input and Claude's
output for each theorem is available upon request.
The Python verification scripts (tiers~2--6) and the
certified outputs of all eight tiers are archived
at~\cite{zenodo}.

%% ========================================================
\begin{acknowledgments}
The author thanks the IBM Quantum and Rigetti developer
communities, and the Open Quantum team, for technical support
during the hardware experiments.
Quantum hardware experiments used IBM Quantum
(IBM Quantum Platform) and Rigetti Computing (Open Quantum).
This work received no specific grant from any funding agency,
public or private.
\end{acknowledgments}

\section*{Data availability}
All data supporting the findings of this study (raw hardware
measurement counts in JSON, simulation outputs, reconstruction
code, the complete 137-operator measurement set, the eight-tier
verification results, and SageMath/FLINT exact-arithmetic certified
values) are archived at~\cite{zenodo} and described in the Supplemental
Material~\cite{supplemental}.
The Lean~4 formal-proof source (Mathlib~v4.29.0-rc8,
0~\texttt{sorry}), SageMath/FLINT exact-arithmetic scripts,
bundled QPU data for zero-configuration external
verification, and a master orchestrator for all verification tiers
are maintained at~\cite{github_verification}.
A Jupyter notebook (\texttt{notebook/k\_star\_demo.ipynb},
runnable in-browser via Binder) reproduces every paper claim
in ${\sim}\,3$~minutes with loud PASS/FAIL assertions; the
tiered verification onramp (\texttt{VERIFICATION.md}) guides a
30~s manifest peek, 2~min axiom-table read, 10~min notebook
run, or 30~min Lean rebuild-from-scratch.

\section*{Competing interests}
The author declares no competing financial or non-financial interests.

%% ========================================================
\appendix
\section{Per-weight operator budget}
\label{app:pauli-set}

Table~\ref{tab:budget} summarizes the per-weight operator allocation
for $n = 4$, $K^* = 5$.  The full list of $137$ operators (as
4-character strings over $\{I,X,Y,Z\}$) and the $29$ tensor-product
bases are provided in the Supplemental Material~\cite{supplemental}.

\begin{table}[h]
\centering
\caption{Per-weight operator budget at $n = 4$, $K^* = 5$.
$A_w = \binom{n}{w} 3^w$ is the number of available weight-$w$
operators; $M_w$ is the number selected.  Weights $0$--$2$ are
fully saturated ($M_w = A_w$); weights $3$--$4$ are budget-limited.}
\label{tab:budget}
\begin{tabular}{cccccc}
\toprule
$w$ & $\binom{4}{w}$ & $3^w$ & $A_w$ & $M_w$ & Status \\
\midrule
0 & 1 & 1  &   1 &  1 & saturated \\
1 & 4 & 3  &  12 & 12 & saturated \\
2 & 6 & 9  &  54 & 54 & saturated \\
3 & 4 & 27 & 108 & 54 & budget-limited \\
4 & 1 & 81 &  81 & 16 & budget-limited \\
\midrule
  &   &    & 256 & 137 & \\
\bottomrule
\end{tabular}
\end{table}

%% ========================================================

\section{Support-completeness criterion}

\pclaim{env:lem:completeness}{proved}{q=2}{lem:completeness}%
\begin{lemma}[Support-completeness criterion ($q = 2$)]
\label{lem:completeness}
If\/ $\mathrm{rank}\,G(K^*) = 2^n$, then for every Hamming weight
$w \in \{0,\ldots,n\}$, at least one lattice vector $m$ with
$|m|^2 \le K^*$ has parity weight $w(\bar{m}) = w$.
Consequently, all $2^n$ support classes of the Pauli algebra are
populated by the lattice at~$K^*$.
\end{lemma}

\begin{proof}
The Krawtchouk eigenvalues satisfy
$\lambda_w(K) = 2^n \, c_w(K) / \binom{n}{w}$, where
$c_w(K) = |\{m \in \mathbb{Z}^n : |m|^2 \le K,\; w(\bar{m}) = w\}|$
counts lattice vectors by parity weight
(see Appendix~\ref{sec:layers} for the full derivation).
Hence $\lambda_w > 0$ implies $c_w > 0$.
Since $\mathbb{Z}^n$ is invariant under coordinate permutations,
$c_w > 0$ means that all $\binom{n}{w}$ parity patterns of weight~$w$
are realized (any one pattern can be permuted to any other).
Hence every support class $s \in \mathbb{Z}_2^n$ appears as the parity
vector of some lattice mode at~$K^*$.
\end{proof}

%%%%%%%%%%%%%%%%%%%%%%%%%%%%%%%%%%%%%%%%%%%%%%%%%%%%%%%%%%%%%%%%
%% Lean4 Formal Verification Status
%% Added: 2026-03-27
%%%%%%%%%%%%%%%%%%%%%%%%%%%%%%%%%%%%%%%%%%%%%%%%%%%%%%%%%%%%%%%%

\bigskip
\noindent\textbf{Lean\,4 formal verification status.}
All 11 formal statements above have been machine-verified in Lean~4
(Mathlib, v4.29.0-rc8) with 0~\texttt{sorry} obligations.
The user-facing trust base consists of four textbook axioms
(A1--A4 below) declared as \texttt{axiom} rather than proved,
owing to gaps in Mathlib's density-operator library; these are
the surface of a broader \texttt{axiom} layer whose
adversarial vetting is described after the list.
The full $\#$\texttt{print axioms} graph is machine-checked:
every \texttt{axiom} is load-bearing.
The consolidated statement file is
\texttt{KstarFormal/Statements.lean};
the full source and a CI pipeline enforcing sorry-freeness
on every commit are maintained
at~\cite{github_verification}.

\medskip
\noindent\emph{Architecture.}
Results that are purely algebraic or combinatorial
(Lemmas~1--6, Corollaries~1--2, Theorems~1 and~3)
are proved without any custom axioms.
The two impossibility conclusions ($F(\rho_+,\rho_-)=0$ and
$\min\{F(\hat\sigma,\rho_+),F(\hat\sigma,\rho_-)\}\le\tfrac12$)
are derived from axioms A2 and A3 via Theorem~2.
The positivity bound
$\varepsilon_{\mathrm{pos}}\le(d{-}1{-}S_k)/d$ (Corollary~1)
takes the Parseval identity and Weyl perturbation bound as
hypotheses and is proved by pure algebra, with no custom axioms.
The four axioms are declared in \texttt{KstarFormal/Axioms.lean}:

\begin{enumerate}\setlength\itemsep{0pt}\tolerance=9999\emergencystretch=2em
\item[\textbf{A1.}] \texttt{witness\_states\_are\_valid}:
  $(I\pm P_0)/d$ are density matrices
  \emph{(Nielsen \& Chuang, \S2.4)}
\item[\textbf{A2.}] \texttt{fidelity\_orthogonal\_zero}:
  $F(\rho_+,\rho_-)=0$ for orthogonal-support states
  \emph{(Uhlmann 1976)}
\item[\textbf{A3.}] \texttt{fidelity\_witness\_traceproj\_decomp}:
  $\exists\, p_\pm \ge 0$ with $p_+ {+}\, p_- = 1$ and
  $F(\hat\sigma,\rho_\pm)\le p_\pm$
  \emph{(Fuchs \& van de Graaf 1999; Schatten-rank decomposition)}.
  The familiar bound
  $F(\hat\sigma,\rho_+){+}F(\hat\sigma,\rho_-)\le 1$
  is derived as a one-line theorem
\item[\textbf{A4.}] \texttt{fidelity\_nonneg}:
  $F(\rho,\sigma)\ge 0$
  \emph{(by definition)}
\end{enumerate}

\begin{sloppypar}\noindent
Running \texttt{\#print axioms} in Lean reveals exactly which of A1--A4 each result depends on.
\end{sloppypar}

\medskip
\noindent\emph{Adversarial vetting of the trust base.}
Axioms A1--A4 are the user-facing surface of the formal
development's \texttt{axiom} declarations. Every declaration
has been audited on six aspects (load-bearing, hypothesis
tightening, generalization stress, definitional anchoring,
import minimization, universe-polymorphism stress) plus a
mutation-based test suite.
Full details and all audit tooling are in
Supplemental Material~\cite{supplemental}.

\medskip
\noindent\emph{Statement cross-reference.}

\smallskip
\noindent
{\footnotesize
\begin{tabular}{@{}llll@{}}
\hline
\textbf{Statement} & \textbf{Lean name} & \textbf{Status} & \textbf{Axioms} \\
\hline
Lemma~1     & \texttt{lem1\_hessian\_diagonal}              & proved       & ---        \\
Theorem~1   & \texttt{thm1\_basin\_separation}              & proved       & ---        \\
Corollary~1 & \texttt{cor1\_approx\_locality\_stmt}         & proved       & ---        \\
Lemma~2     & \texttt{lem2\_purity\_bound}                  & proved       & ---        \\
Lemma~3     & \texttt{lem3\_eigenvalue\_monotonicity\_n4}   & proved       & ---        \\
Theorem~2   & \texttt{thm2\_spectral\_characterization}     & proved$^*$   & A2, A3     \\
Lemma~4     & \texttt{lem4\_hypergeometric\_bound}          & proved       & ---        \\
Corollary~2 & \texttt{cor2\_operator\_lower\_bound\_n4}     & proved       & ---        \\
Theorem~3   & \texttt{thm3\_asymptotic\_separation}         & proved       & ---        \\
Lemma~5     & \texttt{lem5\_spectral\_decomposition\_n4}    & proved       & ---        \\
Lemma~6     & \texttt{lem6\_support\_completeness\_n4}      & proved       & ---        \\
\hline
\end{tabular}}

\smallskip
\noindent
${}^*$\,Conclusion theorem proved from standard axioms A2, A3;
algebraic and combinatorial content fully machine-verified.

\medskip
\begin{sloppypar}\raggedright
\noindent\emph{Scope of formalization.}
The Lean development machine-verifies all eleven formal statements
of the paper (0~\texttt{sorry}).  Lemma~1 (Hessian diagonality),
Lemma~2 (purity bound), Lemma~3 (eigenvalue monotonicity),
Lemma~4 (coupon-collector), Corollary~1 (approximate locality),
and Theorem~1 parts (i)--(iii) are proved at
\emph{universal}~$n$ (or universal dimension~$d$, whichever
the statement ranges over) using only the standard Lean core
axioms \texttt{propext}, \texttt{Classical.choice},
\texttt{Quot.sound}; no custom axiom is introduced.
Theorem~2 parts (i)--(iii) are proved at universal~$n$ modulo two
named textbook axioms (Uhlmann~1976,
\texttt{fidelity\_orthogonal\_zero}; and Fuchs--van~de~Graaf~1999,
\texttt{fidelity\_witness\_traceproj\_decomp}), both declared
with citations in \texttt{KstarFormal.Axioms} and traceable via
\texttt{\#print axioms}.  Theorem~2 part~(iv), and the concrete
numerical witnesses at $n=4$ (lattice count $N=137$, condition
number $\kappa=4$, eigenvalue list $\lambda_w(5)=[144,224,64,128,256]$),
are certified by \texttt{native\_decide}; the $K^*=5$ saturation
for weight classes~$1$ and~$2$ is additionally verified at
$n\in\{4,5,6,7,8\}$, and additionally at \emph{universal}~$n\ge4$
via the compositional patch theorem
\texttt{compositional\_saturation\_w12} (every weight-$\le2$ Pauli
support embeds in a 4-qubit patch, Lean-core axioms only).
The universal-$n$ closed forms
$c_0(5,n)=1+2n$, $c_1(5,n)=2n(2n{-}1)$, $c_2(5,n)=2n(n{-}1)$
are proved symbolically by structural induction on~$n$.
Theorem~3 is stated abstractly over the ratio-decreasing
hypothesis and asymptotic-limit conclusion.
The universal-$n$ wrappers live in \texttt{KstarFormal.Universal},
\texttt{KstarFormal.UniversalSpectralChar},
\texttt{KstarFormal.UniversalKstar},
\texttt{KstarFormal.Combinatorics.CwClosedForm}, and
\texttt{KstarFormal.UniversalKstarCompositional}.
Per-claim axiom footprints are machine-generated (see
supplemental material, Sec.~\ref{sm:lean-universal}).
\end{sloppypar}


%% ============================================================
\section{Three mathematical layers of $K^*$}
\label{sec:layers}

The $K^*$ construction involves three distinct mathematical layers
that should be carefully distinguished:

\begin{enumerate}
\item \textbf{Spectral geometry.}  The Krawtchouk polynomials arise
as eigenfunctions of the Hamming association scheme $H(n,2)$.  The
eigenvalue structure
$K_j(k;n) = \sum_l (-1)^l \binom{k}{l}\binom{n-k}{j-l}$
diagonalizes the distance kernel on the Hamming cube
$\mathbb{Z}_2^n$.  These are the spherical harmonics of the
discrete sphere $\{0,1\}^n$, genuine discrete spectral geometry.

\item \textbf{Lattice point counting.}  The selection rule
$|m|^2 \leq K^*$ counts points in $\mathbb{Z}^n$ inside a ball,
giving $N_n(K^*)$ operators.  This is Gauss's circle problem
generalized to $n$ dimensions.  The value $K^* = 5$ is the minimal cutoff at which
weight saturation holds ($\mu_w = 1$ for $w \le 2$);
the Gram matrix $G(K)$ first achieves full rank at
$K = n = 4$, but without weight saturation
(Sec.~\ref{sec:spectral-necessity}).

\item \textbf{Weight-class budget allocation.}  Given $N_n(K^*)$
operators, the distribution across weight classes
$\{w = 0, \ldots, n\}$ follows from lattice geometry: weight-$w$
operators correspond to lattice points with exactly $w$ nonzero
mod-2 coordinates.  This is purely combinatorial once $N_n(K^*)$ is
fixed (Table~\ref{tab:budget} gives the per-weight breakdown
for $n = 4$, $K^* = 5$).
\end{enumerate}

A natural objection: \emph{``Isn't weight-class completeness
sufficient?  Can't I just saturate low weights without the
spectral framework?''}  This heuristic contains four undetermined
choices that the Krawtchouk spectral analysis resolves:
\begin{itemize}
\item \emph{Which weights to saturate.}  Weight-${\le}\,1$
(13~operators) is insufficient; weight-${\le}\,3$ (175) exceeds
any practical budget.  The Krawtchouk eigenvalue masses show
weights~0--2 carry 65\% of spectral mass in only 67~operators; the
saturation boundary at $w = 2$ is where the lattice budget first
falls below the Pauli multiplicity ($c_3 = 32 < A_3 = 108$).
\item \emph{The total budget.}  $M = 137 = N_4(K^*)$ is a theorem:
$K^* = 5$ is the minimal cutoff satisfying both full Gram rank and
weight-$1$ dominance (Supplemental Material~\cite{supplemental}).  At $K = 4$,
weight-$1$ carries only 9\% of eigenvalue mass; at $K^* = 5$, it
carries 40.9\%.
\item \emph{The high-weight split.}  The remaining 70 operators are
allocated as $54{:}16$ (weight-$3$:weight-$4$), determined by
parity-weight counts $c_3 = 32$, $c_4 = 16$.  The three-arm test
(Sec.~\ref{sec:weight-alloc}) confirms this outperforms proportional
filling ($40{:}30$) in simulation
(Supplemental Material~\cite{supplemental}).
\item \emph{The optimal dimension.}  The Krawtchouk allocation
condition ($2n = 2^{n-1}$, unique at $n = 4$) predicts the weight
allocation becomes most sufficient at $n = 4$
(main text, Sec.~\ref{sec:dim-dep}).
\end{itemize}
Weight-class saturation is a \emph{consequence} of the Krawtchouk
spectral structure, not an alternative to it.
The single derived parameter $K^*$ jointly determines the total
budget $M = N_n(K^*)$, the saturation threshold
$w_{\mathrm{sat}}$, the per-weight allocation $\{M_w\}$, and
the dimensional sweet spot.  The remaining steps (intra-weight
operator ordering and tensor-product basis compression via greedy
set cover) are deterministic and parameter-free but do introduce
additional structure beyond $K^*$ alone
(Supplemental Material~\cite{supplemental}, Sec.~IV).


%% ============================================================
\section{Circuit compression via tensor-product bases}
\label{sec:circuits}

\subsection{Greedy set cover}

The main text presents the $K^*$ set as requiring 137 measurement
circuits (one per selected Pauli operator).  In practice, Pauli
operators are measured via \emph{tensor-product bases}: each qubit
is measured in one of $\{X, Y, Z\}$, yielding a basis labeled by
the choice string $b \in \{X,Y,Z\}^n$.  For $n = 4$ qubits, the
$3^4 = 81$ possible tensor-product bases each yield $2^4 - 1 = 15$
nontrivial Pauli expectation values (all non-identity tensor
products of $\{I, b_i\}$ across qubits).

A greedy set-cover algorithm selects a small number of
tensor-product bases that cover all 136 nontrivial $K^*$
operators:

\begin{table}[ht!]
\caption{Circuit compression of the $K^*$ measurement set at
$n = 4$.  The 29 tensor-product bases that cover all $K^*$
operators also provide 60 additional operators for free.  The
value the covering number $\bar\chi = 29$ is the global optimum of the underlying
set-cover problem, certified by branch-and-bound on a 81-variable
binary integer program (universe $|U|=137$ K* operators, sets
$|S|=3^4=81$ tensor-product bases over $\{X,Y,Z\}^4$); the HiGHS
MIP gap closes to zero at termination.  The solver-independent
\texttt{.mps} model and witness assignment are archived in
\texttt{tier3-numpy/artifacts/clique\_cover/}.}
\label{tab:circuits}
\begin{ruledtabular}
\begin{tabular}{lr}
Metric & Value \\
\hline
$K^*$ operators (incl.\ $IIII$) & 137 \\
Tensor-product bases needed & \textbf{29} \\
Compression ratio (nontrivial) & $4.7\times$ \\
Total estimable operators & \textbf{196 of 255 (76.9\%)} \\
Free additional operators & 60 \\
Missing operators & 7 weight-3, 52 weight-4 \\
\end{tabular}
\end{ruledtabular}
\end{table}

Each tensor-product basis measures 15 operators simultaneously, so
the 29 bases yield $29 \times 15 = 435$ operator--basis pairs.
After removing duplicates (operators appearing in multiple bases),
196 unique nontrivial operators are estimable, or 76.9\% of the full
255, up from the nominal 53.3\% $K^*$ coverage.

The 60 free additional operators (those estimable from the 29 bases
but not in the $K^*$ set) are all at weight~3 (47 operators) and
weight~4 (13 operators), precisely the underdetermined sectors
where additional data is most valuable.

\subsection{Practical impact}

The circuit cost drops from $137 \times N_s$ to $29 \times N_s$
shots, a \textbf{79\% reduction}.  At $N_s = 1000$ shots per
circuit, this is $29{,}000$ vs.\ $137{,}000$ total shots.

For the compositional architecture (main text, Sec.~\ref{sec:compositional}), the
improvement is even larger: $\le 4 \text{ colors} \times 29
\text{ circuits} = 116$ circuits (before basis deduplication)
for any linear chain, compared
to $4 \times 137 = 548$ in the per-operator scheme.

\subsection{Compressibility advantage over random selection}
\label{sec:compression-advantage}

The structured weight distribution of the $K^*$ set yields a
secondary advantage: superior circuit compressibility.
An equal-size random selection ($137$~operators, seed $= 42$)
requires $50$~tensor-product bases under the same greedy set
cover ($2.7\times$ compression), compared with $29$~bases for
$K^*$ ($4.7\times$).  This $1.7\times$ compressibility gap
arises because $K^*$ operators concentrate on low Krawtchouk
weights, whose non-identity components align into fewer shared
measurement bases.

The gap widens with $n$:

\begin{center}
\begin{tabular}{crrrrr}
\toprule
$n$ & $M$ & $K^*$ bases & Rand bases & $K^*$ ratio & Gap \\
\midrule
2 &   9 &  4 &  5 & $2.2\times$ & $1.25\times$ \\
3 &  27 & 13 & 11 & $2.1\times$ & $0.85\times$ \\
4 & 137 & 29 & 50 & $4.7\times$ & $1.72\times$ \\
5 & 333 & 65 & 112 & $5.1\times$ & $1.72\times$ \\
\bottomrule
\end{tabular}
\end{center}

At $n = 3$ the random set happens to compress better (the small
operator count admits a fortuitous packing); from $n = 4$ onward
$K^*$ is consistently $1.7\times$ more circuit-efficient.

For the Rigetti Ankaa-3 same-session comparison (main text,
Sec.~\ref{sec:rigetti}), basis grouping reduced the total circuit count from
$274$ (two individual arms) to $29 + 50 = 79$, a $71\%$ cost
reduction while measuring identical operator sets.
The resulting fidelities ($F_{K^*} = 0.815$,
$F_{\mathrm{rand}} = 0.568$, $\Delta F = +0.248$) confirm that
basis grouping does not degrade reconstruction quality.

\subsection{Connection to classical shadows}
\label{sec:sm-shadows}

The 29-base compressed $K^*$ protocol is operationally equivalent
to a \emph{derandomized classical shadow} scheme~\cite{Huang2020,
Huang2021}.  In randomized shadows, one samples random
tensor-product bases, measures, and reconstructs observables via
median-of-means estimation.  Derandomized variants select a small,
fixed set of bases that covers a target observable list, which is exactly
what the $K^*$ greedy set cover produces.

The key distinction is the selection criterion: shadow
derandomization uses greedy optimization against a
\emph{user-supplied} observable list, while $K^*$ derives the
observable list from spectral geometry (all operators with
$|m|^2 \leq K^*$) and then applies greedy compression.  The two
frameworks converge on the same operational object, namely a small set of
tensor-product bases providing efficient multi-observable
estimation, but $K^*$ provides a \emph{principled,
state-independent rule} for choosing which observables to target.

This positions $K^*$ as a \textbf{spectral-geometric front end}
for shadow tomography: use Krawtchouk structure to decide
\emph{what} to measure, then use shadow estimation machinery to
decide \emph{how} to allocate shots and reconstruct.


%% ============================================================
\section{Reduced density matrix fidelities}
\label{sec:rdm}

\subsection{Motivation}

The full-state fidelity $F \approx 0.36$ (linear inversion) or
$F \approx 0.87$ (MLE) reported in the main text measures
reconstruction quality over the entire $2^4$-dimensional Hilbert
space.  For near-term applications (variational quantum eigensolver, molecular energy
estimation, entanglement detection), the operationally relevant
metrics are the fidelities of \emph{reduced density matrices}
(RDMs): how well are the 1- and 2-body marginals recovered?

Because the $K^*$ set includes \emph{all} 12 weight-1 and all 54
weight-2 operators, it provides informationally complete coverage
of every 1-body and 2-body RDM by construction.
This is a direct consequence of weight-class saturation
(Theorem~\ref{thm:basin}): the same algebraic property that
guarantees basin separation also guarantees exact RDM recovery,
providing independent operational evidence for the theorem's
practical relevance.

\subsection{Hardware results}

From the four hardware W-state runs on \texttt{ibm\_fez} (linear
inversion with eigenvalue clipping):

\begin{table}[ht!]
\caption{RDM fidelities from four hardware W-state runs.
Uncertainties are standard deviations across runs (1-RDM: 16
values = 4 qubits $\times$ 4 runs; 2-RDM: 24 values = 6 pairs
$\times$ 4 runs).}
\label{tab:rdm}
\begin{ruledtabular}
\begin{tabular}{lccc}
Metric & Structured & Random & Advantage \\
\hline
Full-state $F$ & $0.362 \pm 0.005$ & $0.332 \pm 0.020$ & $+0.030$ \\
1-RDM $F$ & $\mathbf{0.985 \pm 0.004}$ & $0.969 \pm 0.022$ & $+0.016$ \\
2-RDM $F$ & $\mathbf{0.771 \pm 0.031}$ & $0.675 \pm 0.045$ & $+0.097$ \\
\end{tabular}
\end{ruledtabular}
\end{table}

\subsection{Analysis}

The near-perfect 1-RDM fidelity ($F = 0.985$) follows directly
from weight-1 saturation: every single-qubit Bloch vector is fully
determined by the $K^*$ measurements.  The 2-RDM advantage of
$+0.097$ reflects weight-2 saturation: all two-body correlations
are measured, while the random arm misses approximately half
(expected 28.8/54 weight-2 operators).

Three additional observations:
\begin{enumerate}
\item The random arm achieves $F = 0.969$ for 1-RDMs because
weight-1 operators are overrepresented at small $n$ (12 out of 255,
but only 4.7\% of the space; random selection of 137 operators
captures $\sim$6.4 on average).
\item The random arm's standard deviation is consistently higher
than structured ($4$--$5.5\times$ for full-state and 1-RDM,
$1.5\times$ for 2-RDM), reinforcing the stability advantage.
\item Full-state fidelity substantially understates reconstruction
quality for local observables: the jump from $F = 0.36$ (full
state) to $F = 0.985$ (1-RDM) reflects the fact that local
information is fully determined while global correlations are
underdetermined.
\end{enumerate}

\subsection{Reframing the value proposition}

This reframes $K^*$ from a full-state tomography method (where
137/255 operators leave the state underdetermined) to an
\emph{informationally complete method for all 1- and 2-body
observables}, with partial information about higher-order
correlations.  For the majority of near-term quantum computing
applications, local observable quality is what matters.


%% ============================================================
\section{Fisher information and frame theory}
\label{sec:fisher}

A natural question is whether standard measurement-design criteria
(A-optimality, D-optimality, frame potential, coherence) favor the
$K^*$ set over random Pauli subsets.  The answer is \textbf{no}:
these criteria are \emph{completely uninformative} for Pauli
tomography.

\subsection{Fisher information equivalence}
\label{sec:fisher-equiv}

For Pauli tomography at the state-independent design point
($\rho = I/d$), each selected operator $P_j$ contributes
independently to the Fisher information matrix:
\begin{equation}
F_{jj} = \frac{N_s \, d^2}{\sigma^2}
\quad \text{if } P_j \in \mathcal{S}, \qquad
F_{jj} = 0 \quad \text{otherwise},
\end{equation}
with all off-diagonal elements zero (Pauli orthogonality).
Consequently:
\begin{itemize}
\item All nonzero eigenvalues of $F$ are identical
($= N_s d^2 / \sigma^2$).
\item $\det(F) = 0$ for any proper subset ($M < d^2 - 1$).
\item The pseudoinverse trace $\mathrm{tr}(F^+)$ is identical for
all $M$-element subsets: $\mathrm{tr}(F^+) = M \sigma^2 / (N_s d^2)$.
\item A-optimality, D-optimality (pseudoinverse), and E-optimality
are \emph{identical} for all Pauli subsets of the same size.
\end{itemize}

$K^*$ has \textbf{no state-independent Fisher advantage}.

\subsection{State-dependent Fisher information}

At $\rho \neq I/d$, the Fisher information per shot for operator
$P_j$ depends on $\langle P_j \rangle_\rho$.  For states with
local structure (W state, product states), low-weight operators
carry larger expectations.  The $K^*$ set captures all such
operators:

\begin{table}[ht!]
\caption{Weight-class coverage: $K^*$ vs.\ expected random
(137 of 255).  Each row shows the number of operators included
per weight class and, in the last row, the probability that
a random draw of 137 operators includes \emph{all} weight-1
Paulis.}
\label{tab:weight-coverage}
\begin{ruledtabular}
\begin{tabular}{lcc}
Metric & $K^*$ & Random (expected) \\
\hline
Weight-1 & 12/12 (100\%) & 6.4/12 (53\%) \\
Weight-2 & 54/54 (100\%) & 28.8/54 (53\%) \\
$P(\text{all wt-1 included})$ & 1.000 & 0.0004 \\
$P(\text{all wt-0/1/2})$ & 1.000 & $5 \times 10^{-23}$ \\
\end{tabular}
\end{ruledtabular}
\end{table}

The probability that uniform random selection of 136 nontrivial
operators (from 255) includes all 66 weight-1/2 operators is
$\sim 5 \times 10^{-23}$: essentially zero. % vague-verb-justified: numerical; 5e-23 is negligible to machine precision

\subsection{Frame theory equivalence}
\label{sec:frame}

Frame-theoretic quality metrics are likewise identical for all
Pauli subsets:
\begin{itemize}
\item \textbf{Frame potential}: $\mathrm{FP} = \sum_{i,j}
|\mathrm{tr}(P_i P_j)|^2 = M \cdot d^2 = 137 \times 256 =
35{,}072$ for any 137-element subset (Pauli orthogonality forces
all cross-terms to zero).
\item \textbf{Sensing matrix condition number}: $\kappa = 1$ for
any Pauli subset (all nonzero singular values equal 1).
\item \textbf{Coherence}: $\mu = 0$ (all pairwise correlations
vanish).
\end{itemize}

Standard compressed-sensing quality metrics (RIP constants,
coherence bounds) cannot distinguish $K^*$ from any other
137-element Pauli subset.

\subsection{The advantage is nonlinear: MLE amplification mechanism}
\label{sec:mle-mechanism}

The $K^*$ advantage arises entirely from the interaction between
weight-class completeness and \emph{nonlinear} MLE reconstruction.
Under linear inversion (pseudoinverse), the advantage is small
($\Delta F \approx +0.03$).  Under MLE, it is large
($\Delta F \approx +0.33$).

Section~\ref{sec:head-to-head} demonstrates this empirically: the structured arm
converges (fidelity plateaus by iteration $\sim$300, per-step
changes $\|\Delta\rho\|_F \sim 0.03$) while the random arm
oscillates without stabilizing ($\|\Delta\rho\|_F \sim 0.3$).
The mechanism can be understood through three observations:

\begin{enumerate}
\item \textbf{Local consistency constrains the Bloch polytope.}
Complete weight-1/2 data provides an overconstrained portrait of
all single-qubit and two-qubit reduced density matrices.  With all
67 weight-${\le}\,2$ operators present (identity plus all weight-1/2), the $k$-RDM consistency
conditions interlock to define a small feasible region in the
$255$-dimensional Bloch body ($d^2 - 1$ real parameters).

\item \textbf{Positivity propagates local constraints globally.}
The MLE's positivity constraint ($\rho \geq 0$) forces the global
state to be compatible with \emph{all} 1- and 2-body marginals
simultaneously, a strong constraint that narrows the likelihood
surface to a well-defined basin.  When local data is incomplete,
multiple incompatible local portraits are consistent with
positivity, creating a rugged or flat likelihood landscape.

\item \textbf{Missing local operators create reconstruction
degeneracies.}  If even one weight-1 operator is absent, the
corresponding single-qubit Bloch vector has an unconstrained
component.  The MLE must infer it from higher-weight operators and
positivity alone, a poorly conditioned inverse problem.  This
explains the oscillatory convergence: the MLE cycles between
locally consistent solutions that differ in the unconstrained
subspace.
\end{enumerate}

This mechanism also explains why the weight-allocated random
baseline (same weight-class budget, random intra-class selection)
captures $\sim$81\% of the structured advantage (main text,
Sec.~\ref{sec:weight-alloc}): it achieves the same weight-class completeness, so the
MLE amplification operates similarly.  The residual $\sim$19\%
advantage of $K^*$ arises from the specific intra-class operator
selection that the spectral-geometric criterion optimizes.

The basin separation theorem (Theorem~\ref{thm:basin} in the main
text) formalizes this mechanism.  The following empirical evidence
provides independent confirmation.

\subsubsection{Weight-class ablation}

The three-arm hardware experiment (Sec.~\ref{sec:dim-dep}) shows
that $76 \pm 12$\% of the $K^*$ advantage originates in the
weight-class allocation.  To isolate which weight classes are
load-bearing, we remove each individually in simulation
(50~random states, $n = 4$, 3\% depolarizing, $10^3$~shots;
Table~\ref{tab:ablation}):

\begin{table}[ht!]
\caption{Weight-class ablation: fidelity impact of removing each
weight class from the $K^*$ operator set ($n = 4$, $M = 137$).}
\label{tab:ablation}
\begin{ruledtabular}
\begin{tabular}{lcrr}
Configuration & $M$ & $F$ (mean$\pm$std) & $\Delta F$ \\
\hline
Full $K^*$ (baseline) & 137 & $0.412 \pm 0.099$ & $-$ \\
Remove weight-0 (1 op) & 136 & $0.105 \pm 0.126$ & $-0.307$ \\
Remove weight-1 (12 ops) & 125 & $0.420 \pm 0.095$ & $+0.008$ \\
Remove weight-2 (54 ops) & 83 & $0.370 \pm 0.073$ & $-0.043$ \\
Remove weight-3 (54 ops) & 83 & $0.386 \pm 0.106$ & $-0.027$ \\
Remove weight-4 (16 ops) & 121 & $0.432 \pm 0.069$ & $+0.019$ \\
\end{tabular}
\end{ruledtabular}
\end{table}

The identity operator (weight-0) is overwhelmingly the most
important single measurement ($\Delta F = -0.31$), consistent with
its role as the trace constraint in MLE\@.  Removing weight-4
operators \emph{improves} full-state MLE fidelity, confirming that
high-weight operators contribute noise in the underdetermined MLE
regime.  (Note: this does not mean weight-4 operators are
uninformative; they carry the GHZ phase signal and enable
$F = 0.926$ under subspace estimation; see main text Sec.~\ref{sec:ghz}.)  (Note: $M$ varies across rows because each ablation removes
a different number of operators.  The weight-2 and weight-3 rows
both have $M = 83$, so their comparison is $M$-controlled:
weight-2 removal hurts more, consistent with the weight-2 dominance
identified below.)

\subsubsection{Progressive weight-class addition}

The three-arm hardware experiment (Sec.~\ref{sec:dim-dep})
tells us that weight allocation accounts for $76 \pm 12$\% of
the $K^*$ advantage, but it cannot tell us \emph{which}
weight-class transition matters most.
Cumulative simulation (adding one weight class at a time,
Table~\ref{tab:progressive}) pins this down:

\begin{table}[ht!]
\caption{Progressive addition of weight classes to the measurement
set ($n = 4$, same noise conditions as Table~\ref{tab:ablation}).}
\label{tab:progressive}
\begin{ruledtabular}
\begin{tabular}{lccc}
Weights included & $M$ & $F$ (mean) & $\Delta F$ (incremental) \\
\hline
$w = 0$ only & 1 & 0.084 & $-$ \\
$w \leq 1$ & 13 & 0.172 & $+0.089$ \\
$w \leq 2$ & 67 & $\mathbf{0.376}$ & $\mathbf{+0.204}$ \\
$w \leq 3$ & 121 & 0.441 & $+0.065$ \\
$w \leq 4$ (full $K^*$) & 137 & 0.422 & $-0.019$ \\
\end{tabular}
\end{ruledtabular}
\end{table}

The big jump ($\Delta F = +0.204$) happens at weight-2 inclusion.
This is where the MLE positivity constraint has enough local
information ($1$- and $2$-body marginals) to propagate upward
and reconstruct higher-weight correlations.
What happens when weight-4 operators are added is perhaps
surprising: fidelity drops slightly ($-0.019$).
In the underdetermined regime, high-weight operators whose
true expectations are near zero contribute mostly shot noise, and
the MLE iteration cannot distinguish this noise from
signal, so it gets pulled away from the estimate that
positivity propagation would have produced.

This is worth pausing on.
It means measurement selection for Pauli tomography
cannot be analyzed with standard frame-theoretic or RIP-based
tools, which assume that more operators always help.
The relevant criterion is weight-class coverage \emph{interacting
with the reconstruction algorithm}, a qualitatively different
regime from generic compressed sensing.

\bibliography{refs}

\end{document}
